\documentclass[11pt]{article}

\usepackage[T1]{fontenc}
\usepackage[utf8]{inputenc}
\usepackage[english]{babel}
\usepackage[margin=1in]{geometry}
\usepackage{lmodern}
\usepackage{microtype}
\usepackage{mathtools}
\usepackage{amsmath,amssymb,amsthm}
\usepackage{array}
\usepackage{booktabs}
\usepackage{tabularx}
\usepackage{float}
\usepackage{graphicx}
\usepackage{tikz}
\usetikzlibrary{positioning}
\usepackage{pdflscape}
\usepackage{caption}
\usepackage{enumitem}
\usepackage{fancyhdr}
\usepackage[numbers]{natbib}
\usepackage{doi}
\usepackage{xurl}
\usepackage{hyperref}
\usepackage{cleveref}
\usepackage{orcidlink}
\usepackage{titling}



\newcommand{\T}{\mathcal T}
\newcommand{\N}{\mathcal N}
\newcommand{\R}{\mathbb R}
\providecommand{\C}{}
\renewcommand{\C}{\mathbb C}
\newcommand{\Dual}{\R[\varepsilon]/(\varepsilon^2)}

\newcommand{\SAU}{\mathsf{SAU}}
\newcommand{\Use}{\operatorname{Useful}}
\newcommand{\Sat}{\operatorname{Sat}}
\newcommand{\Strict}{\operatorname{Strict}}
\newcommand{\StrictUsefulTrace}{\operatorname{StrictUsefulTrace}}
\newcommand{\ValSAU}{\operatorname{ValSAU}}
\newcommand{\Prof}{\operatorname{Profile}}
\newcommand{\NF}{\operatorname{NF}}
\newcommand{\Inst}{\operatorname{Inst}}
\newcommand{\Erase}{\operatorname{Erase}}
\newcommand{\Transfer}{\operatorname{Transfer}}
\newcommand{\Cover}{\operatorname{Cover}}
\newcommand{\Status}{\operatorname{Status}}
\newcommand{\Diag}{\operatorname{Diag}}

\newcommand{\down}{\downarrow}
\newcommand{\ndown}{\not\!\downarrow}
\newcommand{\FactDefect}{\delta_{\mathrm{fact}}}
\newcommand{\SplitDefect}{\delta^{\mathrm{split}}}
\newcommand{\OverreadDefect}{\delta_{\mathrm{overread}}}
\newcommand{\NoSmuggleDefect}{\delta_{\mathrm{nosmuggle}}}
\newcommand{\AuditDefect}{\delta_{\mathrm{audit}}}

\newcommand{\Imag}{\mathsf{Imag}}
\newcommand{\Jet}{\mathsf{Jet}}
\newcommand{\Galois}{\mathsf{Gal}}
\newcommand{\Coh}{\mathsf{Coh}}

\newcommand{\Gal}{\operatorname{Gal}}

\newcommand{\RootComp}{\mathsf{RootComp}}
\newcommand{\Coeff}{\mathsf{Coeff}}
\newcommand{\InvDesc}{\mathsf{InvDesc}}
\newcommand{\AggInv}{\mathsf{AggInv}}
\newcommand{\ObsLed}{\mathsf{ObsLed}}
\newcommand{\ProjShad}{\mathsf{ProjShad}}
\newcommand{\ProdProf}{\mathsf{ProdProf}}

\newcommand{\Repart}{\operatorname{Re}}
\newcommand{\End}{\operatorname{End}}
\newcommand{\Fix}{\operatorname{Fix}}
\newcommand{\id}{\operatorname{id}}
\newcommand{\im}{\operatorname{im}}
\newcommand{\ev}{\operatorname{ev}}
\newcommand{\st}{\operatorname{st}}
\newcommand{\Der}{\operatorname{Der}}
\newcommand{\Tr}{\operatorname{Tr}}
\newcommand{\Nm}{\operatorname{Nm}}
\newcommand{\Aut}{\operatorname{Aut}}
\newcommand{\Orb}{\operatorname{Orb}}
\newcommand{\Stab}{\operatorname{Stab}}
\newcommand{\rank}{\operatorname{rank}}
\newcommand{\Per}{\operatorname{Per}}
\newcommand{\Ob}{\operatorname{Ob}}
\newcommand{\Hom}{\operatorname{Hom}}
\newcommand{\Real}{\operatorname{Real}}
\newcommand{\Core}{\operatorname{Core}}

\newcommand{\Trace}{\operatorname{Trace}}
\newcommand{\AccTrace}{\operatorname{AccTrace}}
\newcommand{\StrictTrace}{\operatorname{StrictTrace}}
\newcommand{\EssProm}{\operatorname{EssProm}}
\newcommand{\Anc}{\operatorname{Anc}}
\newcommand{\Hid}{\operatorname{Hid}}
\newcommand{\Vis}{\operatorname{Vis}}
\newcommand{\Bd}{\operatorname{Bd}}
\newcommand{\Head}{\operatorname{Head}}
\newcommand{\Sort}{\operatorname{Sort}}
\newcommand{\Supp}{\operatorname{Supp}}
\newcommand{\Boundary}{\operatorname{Boundary}}
\newcommand{\Red}{\operatorname{Red}}
\newcommand{\Basic}{\operatorname{Basic}}
\newcommand{\Gen}{\operatorname{Gen}}
\newcommand{\Inv}{\operatorname{Inv}}
\newcommand{\RedSAU}{\operatorname{RedSAU}}
\newcommand{\BasicSAU}{\operatorname{BasicSAU}}
\newcommand{\MinProm}{\operatorname{MinProm}}
\newcommand{\InitProm}{\operatorname{InitProm}}

\newcommand{\Pass}{\mathsf{pass}}
\newcommand{\Fail}{\mathsf{fail}}
\newcommand{\Lit}{\operatorname{Lit}}

\newcommand{\FATCD}{\mathsf{FATCD}}
\newcommand{\Raw}[1]{\operatorname{Raw}\!\left(#1\right)}
\newcommand{\Audit}[1]{\operatorname{Audit}\!\left(#1\right)}
\newcommand{\Closed}[1]{\operatorname{Closed}\!\left(#1\right)}
\newcommand{\Dec}[1]{\operatorname{Dec}\!\left(#1\right)}
\newcommand{\Residual}[1]{\operatorname{Residual}\!\left(#1\right)}
\newcommand{\Chan}[2]{\operatorname{Chan}_{#1}\!\left(#2\right)}
\newcommand{\ScopedExactSix}[1]{\operatorname{ScopedExactSix}\!\left(#1\right)}

\newcommand{\citeFoundationsI}{\citep{TsiokosMathTheory}}
\newcommand{\citeFoundationsII}{\citep{TsiokosFoundationsII}}
\newcommand{\citeFoundationsIII}{\citep{TsiokosFoundationsIII}}
\newcommand{\citeWignerHamming}{\citep{Wigner1960,Hamming1980}}
 
\newtheorem{theorem}{Theorem}[section]
\newtheorem{lemma}[theorem]{Lemma}
\newtheorem{proposition}[theorem]{Proposition}
\newtheorem{corollary}[theorem]{Corollary}
\theoremstyle{definition}
\newtheorem{definition}[theorem]{Definition}
\newtheorem{remark}[theorem]{Remark}
\floatstyle{ruled}
\newfloat{algorithm}{tbp}{loa}
\floatname{algorithm}{Algorithm}

\urlstyle{same}
\captionsetup{
  font=small,
  labelfont=bf,
  labelsep=period
}
\captionsetup[table]{skip=6pt,position=top}
\captionsetup[figure]{skip=6pt}
\captionsetup[algorithm]{skip=6pt,position=top}
\crefname{algorithm}{algorithm}{algorithms}
\Crefname{algorithm}{Algorithm}{Algorithms}
\crefname{theorem}{theorem}{theorems}
\Crefname{theorem}{Theorem}{Theorems}
\crefname{lemma}{lemma}{lemmas}
\Crefname{lemma}{Lemma}{Lemmas}
\crefname{proposition}{proposition}{propositions}
\Crefname{proposition}{Proposition}{Propositions}
\crefname{corollary}{corollary}{corollaries}
\Crefname{corollary}{Corollary}{Corollaries}
\crefname{definition}{definition}{definitions}
\Crefname{definition}{Definition}{Definitions}
\crefname{remark}{remark}{remarks}
\Crefname{remark}{Remark}{Remarks}
\setlist[enumerate]{leftmargin=*, itemsep=2pt, topsep=4pt}
\setlength{\droptitle}{-3.2em}
\pretitle{\begin{center}\LARGE\bfseries}
\posttitle{\par\end{center}\vspace{0.5em}}
\preauthor{\begin{center}\normalsize}
\postauthor{\par\end{center}\vspace{-0.4em}}
\predate{\begin{center}\small}
\postdate{\par\end{center}\vspace{0.6em}}
\setlength{\emergencystretch}{2em}
\clubpenalty=10000
\widowpenalty=10000
\displaywidowpenalty=10000
\interfootnotelinepenalty=10000
\allowdisplaybreaks[2]
\fancypagestyle{firstpage}{\fancyhf{}
  \fancyfoot[C]{\scriptsize Automorph Inc., Wilmington, DE, USA \quad \texttt{ioannis@automorph.io} \quad \textcopyright\ 2026 \quad CC-BY 4.0}
  \renewcommand{\headrulewidth}{0pt}
  \renewcommand{\footrulewidth}{0pt}
}

\title{The Usefulness of Non-Descending Objects: A Six Birds Theory of Mathematical Applicability}
\author{Ioannis Tsiokos\,\orcidlink{0009-0009-7659-5964}}
\date{\today}
\hypersetup{
  pdftitle={The Usefulness of Non-Descending Objects: A Six Birds Theory of Mathematical Applicability},
  pdfauthor={Ioannis Tsiokos},
  hidelinks
}

\begin{document}

\maketitle
\thispagestyle{firstpage}

\begin{abstract}
What makes a mathematical object useful when it does not live in the layer
where the problem is posed?  Imaginary numbers are not real, infinitesimals are
not real scalars, roots of irreducible polynomials need not lie in the base
field, and cohomology classes are not vertex potentials.  This paper studies
that pattern by separating a lower layer, where a problem and its accepted
answers are stated, from a promoted layer, where additional objects may be
available.  We define theory layers, promotion bridges, audited descending
expressions, and the Strict Audited Utility (SAU) certificate, and then show
that SAU is extracted from strict useful traces rather than stipulated.

The main results are a theorem stack: Essential Boundary Necessity (Theorem
A), which forces hidden promoted support and an audited lower boundary; SAU
Boundary Normal Form (Theorem B); a Partial Converse and Tautology Barrier
(Theorem C); SAU Soundness (Theorem D); Basic Profile Classification (Theorem
E); four canonical witnesses from imaginary numbers, dual numbers and first
jets, Galois roots, and graph cohomology; Minimality and universal promotion
(Theorem F); and a Profile Transfer theorem (Theorem G).  Together they show
that, in finite typed audited promoted-use traces, strict mathematical
usefulness forces hidden non-descent plus audited boundary descent.  The claim
is deliberately scoped: the paper does not claim that all mathematics is
useful, does not literalize hidden objects downstairs, and does not claim that
the Lean development independently proves every paper theorem.  The Lean
material is a traceability harness, with substantive proof language reserved
for direct, computational, and definitional proof rows.
\end{abstract}

\medskip
\noindent\textbf{Keywords.} mathematical applicability; non-descending objects; audited descent; Strict Audited Utility; theory layers; promotion bridges; profile transfer; graph cohomology.

\section{Introduction: The Use of Impossible Objects}\label{sec:intro}

\subsection{Useful but illegal objects}
\label{subsec:intro-illegal-objects}

The number $i$ is not a real number, but $i^2=-1$ is.  A calculation may
therefore introduce an element outside \(\R\) while still returning a statement
inside \(\R\); \(\C\) supplies the extra room, but the equation \(i^2=-1\) is
read back in the original real setting.  The same pattern appears for
infinitesimals: in the dual-number algebra \(\Dual\), the symbol
\(\varepsilon\) satisfies \(\varepsilon^2=0\), and the first-order identity
\(f(x+\varepsilon)=f(x)+f'(x)\varepsilon\) exposes the ordinary derivative as a
coefficient.

It appears again in algebra.  If one starts over \(\mathbb Q\), a splitting
field for \(t^2-2\) contains an element \(\alpha=\sqrt 2\) that is not in
\(\mathbb Q\), while its minimal polynomial \(t^2-2\) lies in
\(\mathbb Q[t]\).  It also appears in graph cohomology: a cochain \(a\) may
fail to be exact, but its class \([a]\in H^1\) records a genuine obstruction
visible to the original ledger of vertex potentials.

These four examples are classical mathematics; the paper does not claim
novelty for complex numbers, dual numbers, splitting fields, or graph
cohomology themselves.  Their common shape is what matters here: the
calculation uses an object that is useful for a problem about the original
setting, even though the object does not live in that setting, and it pairs
that object with an expression that the original setting can recognize.  The
case studies in \Cref{sec:imaginary,sec:calculus,sec:galois,sec:cohomology}
return to these examples with their full mathematical structure.

For the rest of the paper, call the original setting the lower layer: the
mathematical world in which the problem, proposed answers, and accepted
outputs are stated.  Call the richer ambient setting the promoted layer: the
larger structure introduced to compute, explain, or certify something about
that lower layer.  The next question is therefore not whether the new object
corresponds to reality, but what, precisely, can be carried back to the
original setting.

\subsection{The question swap}
\label{subsec:intro-question-swap}

The preceding examples suggest a more useful question than ``does this object
correspond to something real?''  The working question is mathematical: what
fails to be expressible in the lower layer, what is expressible there, and
what derivation connects the two?  In this paper, descent means being
expressible, definable, representable, or recoverable in the lower layer.
Non-descent means the failure of that property.

As in the four examples above, the point is not merely that a larger setting
has been introduced.  The imaginary case separates a promoted object from its
lower-layer square; the dual-number case separates the infinitesimal direction
from the derivative coefficient; the splitting-field case separates a root
from its lower-layer polynomial data; and the cohomology case separates a
local cochain from the class that records its obstruction.  In each case, the
useful calculation has two sides: something does not descend, while some
related expression does.

This phrasing turns a philosophical worry into structured mathematical data.
The object's non-descent is part of the data, and so is the derivation that
recovers a lower-layer expression from the promoted calculation.  The next
subsection states the paper's central preview: strict useful mathematics
requires non-descent paired with a controlled descent expression, and the
later theorems analyze both pieces.

\subsection{Contribution preview}
\label{subsec:intro-contribution-preview}

The paper's scoped thesis is that, in finite typed audited promoted-use
traces, strict mathematical usefulness---what this paper calls
\emph{Strict Audited Utility} (SAU), formally defined in
\Cref{sec:layers}---forces hidden non-descent plus audited boundary descent.
Here an audit is an explicit checkable derivation connecting a higher-layer
computation back to a legitimate lower-layer answer, and a boundary is the
single derivation step at which descent happens.  The anti-tautology
requirement is built in: the pairing is extracted from strict useful traces
rather than stipulated as a definition.

The first three previewed results, developed in \Cref{sec:necessity}, explain
why this structure is forced.  Theorem A, Essential boundary necessity, says
that a strict useful trace forces a non-descending support and an audited
descent expression at its boundary.  Theorem B, SAU boundary normal form, says
that this boundary expression has a canonical reduced shape with one primitive
head and a defined support set.  Theorem C, Partial converse and tautology
barrier, says that a controlled re-derivation also runs in the reverse
direction and forces the certificate to be extracted rather than stipulated.

The next two results separate licensed conclusions from tempting
overstatements.  Theorem D, SAU soundness, proved in
\Cref{sec:soundness-failure}, says that an accepted certificate licenses only
the audited lower-layer answer, never literalization of the hidden object.
Theorem E, Basic profile classification, proved in \Cref{sec:classification},
sorts the four primitive head types: root-composite, coefficient or
standard-part, invariant descent, and obstruction ledger.  These are not four
names for one phenomenon; they distinguish four boundary shapes.

The last two results explain why the four case studies are not merely
illustrative.  Theorem F, Minimality and universal promotion, says that the
imaginary, calculus, Galois, and cohomology cases in
\Cref{sec:imaginary,sec:calculus,sec:galois,sec:cohomology} each admit a
smallest or universal promoted package satisfying its profile, meaning a
reusable proof shape of non-descent paired with an audited descent expression.
Theorem G, Profile transfer, proved in \Cref{sec:transfer}, says that a
validated profile is portable to another case under matching audit hypotheses,
while not identifying objects across domains.

\subsection{Scope and roadmap}
\label{subsec:intro-scope-roadmap}

The scope is deliberately narrow.  This paper does not claim that all
mathematics is useful, does not classify every useful mathematical structure,
and does not claim that hidden promoted objects literally belong to the lower
layer.  It also does not claim novelty for complex numbers, dual numbers,
splitting fields, trace and norm, or graph cohomology themselves; those
classical constructions are inputs, not outputs.  The novelty asserted here is
the audited-descent synthesis and the theorem stack built around it.  The
paper relates to Foundations I--III for inherited machinery and provenance
\citeFoundationsI, \citeFoundationsII, \citeFoundationsIII, but does not
depend on them as prerequisites: the native vocabulary used here is defined
locally.

\Cref{sec:layers} gives the formal language of theory packages, lower
problems, promotion bridges, and certificates.  \Cref{sec:necessity} proves
Theorems A--C by extracting necessity, normal form, and the anti-tautology
converse from strict useful traces.  \Cref{sec:soundness-failure} proves
Theorem D and states the failure-mode discipline that controls what a
certificate licenses.  \Cref{sec:classification} proves Theorem E and
separates the four primitive head types.
\Cref{sec:imaginary,sec:calculus,sec:galois,sec:cohomology} develop the four
case studies as canonical witnesses.  \Cref{sec:transfer} proves Theorems F
and G by explaining minimality and profile transfer across cases.
\Cref{sec:outlook} states the contribution and its limits.  The next section
begins that route by moving from the informal motivation here to the formal
vocabulary.
 \section{Theory Layers and Promoted-Use Traces}\label{sec:layers}

\subsection{Theory packages and instruments}
\label{subsec:layers-packages-instruments}

A theory layer is an ordinary mathematical setting together with the way it is
being observed.  It may be a base field \(K\) before adjoining a root, the
real numbers \(\R\) before passing to \(\C\), value-only data before first
jets, or the cochain ledger before \(H^1\).  In each case there is not only a
carrier of objects, but also a lens: an observation map or interface recording
which distinctions the layer can see.

\begin{definition}[Theory package]
\label{def:sau-theory-package}
A theory package is a tuple
\[
\mathcal T=(Z,f,\Sigma_f,E,\mathcal A),
\]
where \(Z\) is a carrier, \(f:Z\to O\) is a lens or interface, \(\Sigma_f\) is the family of distinctions representable through \(f\), \(E\) is a completion or packaging operation when one is present, and \(\mathcal A\) is the inherited package audit functional. In the full finite-lens case,
\[
\Sigma_f=\{f^{-1}(B):B\subseteq \operatorname{im}(f)\}.
\]
The paper may use weaker or profile-specific \(\Sigma_f\), but every use must be declared under the surrounding instrument.
\end{definition}

This definition is deliberately richer than a bare set or a bare structure:
the lens \(f\) says which distinctions are visible, and \(\mathcal A\) records
the declared audit functional for that package.  The accompanying instrument
is the policy record that says when a higher-layer computation counts as an
accepted lower-layer answer.

\begin{definition}[Instrument]
\label{def:sau-instrument}
An instrument is a finite typed record
\[
I=(\mathsf{vis}_I,\Theta_I,\mathsf{BridgePolicy}_I,\mathsf{AuditPolicy}_I,
\mathsf{SourcePolicy}_I,\mathcal N_I),
\]
where \(\mathsf{vis}_I\) partitions content into
\[
\{\texttt{visible},\texttt{suppressed},\texttt{outside},\texttt{unknown}\},
\]
\(\Theta_I\) gives threshold rules, \(\mathsf{BridgePolicy}_I\) gives admissible bridge types, \(\mathsf{AuditPolicy}_I\) gives required audit records, \(\mathsf{SourcePolicy}_I\) gives source-of-truth requirements, and \(\mathcal N_I\) is the nonclaim register. A claim accepted under \(I\) is an instrument-indexed claim, not an instrument-free truth claim.
\end{definition}

\subsection{Lower problem profiles and saturation}
\label{subsec:layers-profiles-saturation}

A lower problem profile is the organized version of a question asked inside
the lower layer.  It records the problem, the candidates the layer is allowed
to try, the kind of answer that would count, the defect by which a candidate is
judged, the threshold for passing, and the surrounding nonclaim register.

\begin{definition}[Lower problem profile]
\label{def:sau-lower-profile}
A lower problem profile is a tuple
\[
\mathfrak P_0=(\mathcal T_0,I_0,P,\mathsf{Cand}_0(P),\mathsf{Ans}_0(P),
\delta_0,\Theta_0,\mathcal N_0),
\]
where \(\mathcal T_0\) is the lower theory package, \(I_0\) is the lower instrument, \(P\) is the posed problem, \(\mathsf{Cand}_0(P)\) is the class of lower-layer candidate answers or certificates, \(\mathsf{Ans}_0(P)\) is the lower-layer accepted answer type, \(\delta_0\) is the problem defect, \(\Theta_0\) is the pass threshold, and \(\mathcal N_0\) records what the lower profile does not claim.
\end{definition}

Saturation is the corresponding negative condition: after fixing the lower
interface and the allowed candidates, the lower layer cannot meet the problem's
own success test.

\begin{definition}[Lower-layer saturation]
\label{def:sau-saturation}
The lower problem profile \(\mathfrak P_0\) is saturated, written
\[
\operatorname{Sat}(\mathfrak P_0),
\]
when no candidate available at the fixed lower interface passes the problem threshold:
\[
\forall s\in\mathsf{Cand}_0(P),\qquad
\mathsf{Eval}_{\Theta_0}(\delta_0(s))\neq\texttt{pass}.
\]
In exact zero-defect cases this is
\[
\forall s\in\mathsf{Cand}_0(P),\qquad \delta_0(s)\neq 0.
\]
Saturation is always relative to the declared lower interface, instrument, and candidate class; it is not an absolute impossibility claim about all possible mathematics.
\end{definition}

An accepted lower answer is the positive counterpart.  It is not just a correct
object in isolation, but an answer that is visible at the lower interface,
passes the declared defect threshold, and respects the policies attached to
the lower problem.

\begin{definition}[Accepted lower answer]
\label{def:sau-accepted-answer}
An object \(a\) is an accepted lower answer to \(P\), written
\[
a\in\mathsf{Ans}_0(P),
\]
when \(a\) is visible or admissibly bridged under \(I_0\), passes the threshold \(\Theta_0\), satisfies the source and audit policy of \(I_0\), and respects the lower nonclaim register \(\mathcal N_0\). Accepted answers may be solution witnesses, decision witnesses, invariant values, obstruction certificates, or audited negative answers, depending on the problem profile.
\end{definition}

The imaginary case supplies the first finite preview.  There the real-shadow
lens \(q:C_4\to Y_3\) sends a four-element promoted carrier to a three-element
shadow carrier; the lower problem asks for a half-step of real reversal, and
the lower candidates do not pass that test, so the lower problem profile is
saturated.  This paragraph is only a preview of the example; the full
construction appears in \Cref{sec:imaginary}.

\subsection{Promotion bridges}
\label{subsec:layers-promotion-bridges}

A generic bridge is any declared way of comparing an enriched setting with the
lower one, often by an embedding together with a way to pull information back.
A promotion bridge is more specific: it is the typed package that moves from a
lower theory package to a promoted theory package and returns only the
lower-layer descendants allowed by its declared channel and checks.  Thus it
records not only the promoted ambient structure, but also the interfaces,
visibility data, gates, audits, and nonclaims governing the return trip.

\begin{definition}[Promotion bridge]
\label{def:sau-promotion-bridge}
A promotion bridge from a lower theory package \(\mathcal T_0\) to a promoted theory package \(\mathcal T_1\) is a finite typed record
\[
B_{0\to1}=(\mathcal T_0,\mathcal T_1,I_B,S,q,\pi_0,\pi_1,U,C,
\mathcal G_B,A_B,\mathcal L_B,\mathcal N_B),
\]
where:
\begin{itemize}
\item \(I_B\) is the bridge instrument;
\item \(S\) is the carrier of compared states, witnesses, or presentations;
\item \(q\) is the descent, shadow, quotient, standard-part, invariant, coefficient, trace, or obstruction channel appropriate to the profile;
\item \(\pi_0:S\to O_0\) is the old interface and \(\pi_1:S\to O_1\) is the promoted interface;
\item \(U\) is the promoted object, generator, operation, infinitesimal, root, local datum, or cocycle;
\item \(C\) is the expression former selecting the composite, invariant, coefficient, trace, norm, standard part, or obstruction class that is intended to descend;
\item \(\mathcal G_B\) records gate results;
\item \(A_B\) records audits;
\item \(\mathcal L_B\) records admissible visibility bridges, while \(q\) records the named descent or shadow channel;
\item \(\mathcal N_B\) is the bridge nonclaim register.
\end{itemize}
The bridge is a typed record. It is not accepted merely because \(\mathcal T_1\) exists.
\end{definition}

In this record, \(q\) is the named channel by which a lower-layer descendant is
read off, while \(\pi_0\) and \(\pi_1\) say what the lower and promoted
interfaces see on the same compared state.  The component \(U\) is the hidden
promoted datum used in the calculation, \(C\) is the expression-forming rule
whose output is meant to return downstairs, \(A_B\) stores the checks for that
return, \(\mathcal G_B\) stores the gate outcomes, \(\mathcal L_B\) stores the
admissible visibility comparisons, and \(\mathcal N_B\) states what the bridge
does not license.

The running examples fit this template informally.  The inclusion
\(\R\hookrightarrow\C\) with real-shadow audit, the inclusion
\(\R\hookrightarrow\Dual\) with coefficient audit, and an inclusion
\(K\subset L\) with trace or norm audit are promotion bridges in the sense
developed in \Cref{sec:imaginary,sec:calculus,sec:galois}.

\subsection{Descent, non-descent, and strictness}
\label{subsec:layers-descent-nondescent-strictness}

Exact descent is the typed version of the phrase ``the promoted update has a
well-defined lower shadow.''  For a map \(q:X\to Y\) and an update \(F:X\to X\),
the notation \(F\downarrow_q\) means that there is a lower update
\(F^\sharp:Y\to Y\) with \(qF=F^\sharp q\).  Equivalently, applying \(F\)
upstairs and then observing through \(q\) gives the same lower result as first
observing through \(q\) and then applying \(F^\sharp\); the full formal
definition is recorded in \Cref{app:definitions}.

\[
\begin{array}{ccc}
X & \xrightarrow{\;F\;} & X\\[-1mm]
{\scriptstyle q}\!\downarrow & & \downarrow\!{\scriptstyle q}\\[-1mm]
Y & \xrightarrow{\;F^\sharp\;} & Y
\end{array}
\qquad
qF=F^\sharp q.
\]

The split-pair descent defect, written
\(\delta^{\mathrm{split}}_1(q,F)\), is the failure mode for that square.  It
records pairs of promoted points with the same lower image whose \(F\)-updates
have different lower images.  Thus descent fails exactly when two points that
look identical downstairs split apart after the promoted update; the formal
set-theoretic definition is deferred to \Cref{app:definitions}.

Strictness is a different comparison.  The factorization defect
\(\delta_{\mathrm{fact}}(\pi_0,\pi_1)\) counts pairs of compared states with
the same lower view under \(\pi_0\) but different promoted views under
\(\pi_1\).  A promotion is strict when this set is nonempty: the promoted layer
can distinguish something that the lower layer identifies.  The formal
definition of this defect also appears in \Cref{app:definitions}.

Finally, \(U\not\downarrow_q\) means that the promoted datum \(U\) itself is
not acceptable as a lower-layer object, operation, or claim through the
declared lower interface.  The body needs this definition because the later
arguments distinguish the promoted datum from the lower expression that is
eventually accepted.

\begin{definition}[Non-descending promoted structure]
\label{def:sau-nondescending-structure}
A promoted structure \(U\) is non-descending relative to \((q,I_B)\), written
\[
U\not\downarrow_q,
\]
when the declared lower interface cannot accept \(U\) as a lower-layer object, operation, or claim. In exact update profiles this means
\[
\delta^{\mathrm{split}}_1(q,U)\neq\varnothing.
\]
In non-update profiles this means that the relevant lower representability, visibility, or membership defect is nonzero; for example \(\varepsilon\notin\mathbb R\), \(\alpha\notin K\), or a local datum fails to glue globally.
\end{definition}

\subsection{The SAU certificate}
\label{subsec:layers-sau-certificate}

An audited descending expression is the return leg of a promoted computation.
It says that, under a specific promotion bridge, the expression \(C(U)\) formed
from the promoted datum \(U\) descends to a lower answer \(a^\sharp\), and that
the bridge checks licensing this return have passed.

\begin{definition}[Audited descending expression]
\label{def:sau-audited-desc-expression}
Let \(B_{0\to1}\) be a promotion bridge with promoted structure \(U\) and expression former \(C\). The expression \(C(U)\) descends with audit to a lower answer \(a^\sharp\), written
\[
C(U)\downarrow_{B} a^\sharp,
\]
when:
\begin{enumerate}
\item \(C(U)\) has the profile-appropriate descent property, such as \(qC(U)=F^\sharp q\), coefficient extraction, standard-part projection, invariant descent, trace/norm descent, or obstruction-class descent;
\item \(a^\sharp\in\mathsf{Ans}_0(P)\) under the lower problem profile;
\item the bridge audit \(A_B\) passes;
\item \(\delta_{\mathrm{overread}}=\varnothing\) for the claim using \(a^\sharp\);
\item \(\delta_{\mathrm{nosmuggle}}(B)=\varnothing\);
\item the nonclaim register states that only the descended expression \(C(U)\), not the whole hidden structure \(U\), is claimed at the lower layer.
\end{enumerate}
\end{definition}

The no-overreading and no-smuggling clauses are forward references to the
soundness discipline of \Cref{sec:soundness-failure}: no-overreading means
claiming only the audited descended expression, not the hidden promoted object,
and no-smuggling means that the bridge does not silently insert the desired
answer into the lower layer.  With these checks in place, a Strict Audited
Utility (SAU) certificate is the bundled six-field record of saturation,
strictness, non-descent, audited descent, accepted lower answer, and nonclaim.

\begin{definition}[Strict Audited Utility certificate]
\label{def:sau-certificate}
Let \(\mathfrak P_0\) be a lower problem profile and \(B_{0\to1}\) a promotion bridge. A Strict Audited Utility certificate, abbreviated
\[
\mathsf{SAU}(\mathfrak P_0,B_{0\to1},U,C,a^\sharp),
\]
is a finite typed record containing the following six fields:
\begin{enumerate}
\item \textbf{Saturation field:}
\[
\operatorname{Sat}(\mathfrak P_0).
\]
\item \textbf{Strictness field:}
\[
\delta_{\mathrm{fact}}(\pi_0,\pi_1)\neq\varnothing.
\]
\item \textbf{Non-descent field:}
\[
U\not\downarrow_q.
\]
\item \textbf{Audited descent field:}
\[
C(U)\downarrow_B a^\sharp.
\]
\item \textbf{Problem-resolution field:}
\[
a^\sharp\in\mathsf{Ans}_0(P)
\]
with \(\mathsf{Eval}_{\Theta_0}(\delta_0(a^\sharp))=\texttt{pass}\), when the lower problem uses a defect threshold.
\item \textbf{Audit/no-overread/no-smuggling field:}
\[
\delta_{\mathrm{overread}}=\varnothing,
\qquad
\delta_{\mathrm{nosmuggle}}(B)=\varnothing,
\qquad
\delta_{\mathrm{audit}}(B)=\varnothing,
\qquad
\mathsf{AuditGate}_B:\texttt{pass},
\]
with an explicit nonclaim record \(\mathcal N_B\).
\end{enumerate}
\end{definition}

The certificate pattern can be read in its shortest form as
\[
\boxed{U\ndown_q,\qquad C(U)\down_B a^\sharp.}
\]
The rest of the paper studies what this record forces in
\Cref{sec:necessity}, what it licenses in \Cref{sec:soundness-failure}, how
its shapes are classified in \Cref{sec:classification}, and how validated
shapes transfer across cases in \Cref{sec:transfer}.  The next section asks
whether this pattern is forced by strict useful traces rather than stipulated.
 \section{Necessity and Normal Form}\label{sec:necessity}

\subsection{Promoted-use traces and the anti-tautology setup}
\label{subsec:necessity-traces-setup}

The SAU certificate of \Cref{sec:layers} could otherwise be dismissed as a
definition designed to make the central theorem true.  The next two sections
avoid that circularity by starting from strict useful traces, which do not
mention SAU certificates.  The aim is to show that a boundary expression and a
non-descending support are forced by the trace itself; the certificate is then
extracted from the trace, not stipulated in advance.

\begin{definition}[Finite promoted-use trace]
\label{def:sau-promoted-use-trace}
A finite promoted-use trace for a lower problem profile \(\mathfrak P_0\) is a tuple
\[
\mathfrak T=(D,V,E,r,o,\lambda,\upsilon,\mathcal B,A_{\mathfrak T},\mathcal N,v_\ast),
\]
where:
\begin{enumerate}
\item \(D=(V,E)\) is a finite directed acyclic graph.  An edge \(u\to v\) means that the node \(v\) depends on the node \(u\).
\item \(r\) records the rule, operation, bridge, or audit step used at each non-input node.
\item \(o(v)\) is the mathematical object, assertion, certificate, or intermediate expression carried by \(v\).
\item \(\lambda(v)\in\{0,1\}\) is the declared layer label: \(0\) for lower-layer nodes and \(1\) for promoted-layer nodes.
\item \(\upsilon(v)\in\{\texttt{visible},\texttt{suppressed},\texttt{outside},\texttt{unknown}\}\) is the visibility status under the declared instrument.
\item \(\mathcal B\) is the finite family of declared bridge or descent steps used by the trace.
\item \(A_{\mathfrak T}\) is the finite audit/provenance/replay record for the trace.
\item \(\mathcal N\) is the nonclaim register.
\item \(v_\ast\in V\) is the final answer node, whose object is denoted
\[
a^\sharp_{\mathfrak T}:=o(v_\ast).
\]
\end{enumerate}
The trace is typed: every edge and every rule application must respect the source and target layers recorded by \(\lambda\), the visibility status recorded by \(\upsilon\), and the bridge permissions recorded by \(\mathcal B\).
\end{definition}

A trace is a typed directed acyclic graph for a derivation that uses promoted
content.  Each node carries an object or expression, a layer label, and a
visibility status; bridge steps and audit records are part of the data; the
final answer is the distinguished node \(v_\ast\).

\begin{definition}[Accepted promoted-use trace]
\label{def:sau-accepted-promoted-use-trace}
A finite promoted-use trace \(\mathfrak T\) is accepted, written
\[
\AccTrace(\mathfrak T),
\]
when all of the following hold:
\begin{enumerate}
\item the final answer is a lower accepted answer:
\[
a^\sharp_{\mathfrak T}\in\mathsf{Ans}_0(P),
\qquad
\mathsf{Eval}_{\Theta_0}(\delta_0(a^\sharp_{\mathfrak T}))=\texttt{pass};
\]
\item every bridge in \(\mathcal B\) belongs to the declared admissible bridge family;
\item every use of suppressed content is either internal to the promoted layer or is discharged by an admissible bridge before a lower-layer claim is made;
\item the audit record \(A_{\mathfrak T}\) passes the declared source, provenance, replay, and threshold checks;
\item the nonclaim register \(\mathcal N\) blocks the claim that hidden promoted nodes are lower-layer objects unless such a claim is separately bridged and audited.
\end{enumerate}
Acceptance is weaker than an SAU certificate: it says that a promoted derivation produced an accepted lower answer without audit failure.  It does not yet identify the normal-form pair \((U,C(U))\).
\end{definition}

Acceptance is a precondition, not the certificate itself.  It says the trace
has produced a lower answer through admissible bridge and audit steps, but it
does not yet identify the hidden object, the descended expression, or the
normal form that later sections will extract.

\begin{definition}[Strict useful trace]
\label{def:sau-strict-useful-trace}
A finite promoted-use trace \(\mathfrak T\) is a strict useful trace for the lower problem profile \(\mathfrak P_0\), written
\[
\operatorname{StrictUsefulTrace}(\mathfrak T,\mathfrak P_0),
\]
when
\[
\AccTrace(\mathfrak T),
\qquad
\operatorname{Sat}(\mathfrak P_0),
\qquad
\operatorname{StrictTrace}(\mathfrak T),
\qquad
\EssProm(\mathfrak T)
\]
all hold.
This is a pre-SAU notion.  It does not assume an SAU certificate; it only assumes an accepted finite promoted derivation whose promoted content is strict and essential.
\end{definition}

A strict useful trace combines acceptance with lower-layer saturation, trace
strictness, and essential promotion.  In plain terms, the lower layer cannot
solve the problem by its own candidates, the promoted content does not collapse
to a lower derivation, and the promoted part is essential rather than
decorative.  This definition still does not mention SAU.

\begin{lemma}[Boundary existence]
\label{lem:sau-boundary-existence}
If \(\operatorname{StrictUsefulTrace}(\mathfrak T,\mathfrak P_0)\) holds, then
\[
\Bd(\mathfrak T)\neq\varnothing.
\]
\end{lemma}

\begin{proof}
By acceptance, the final answer \(a^\sharp_{\mathfrak T}=o(v_\ast)\) is a
lower-layer accepted answer.  By essential promotion, some promoted node is in
the dependency cone of that final answer and is necessary for producing it.
Following a directed path from such a promoted node toward \(v_\ast\), the layer
label must eventually change from promoted to lower, because the endpoint is a
lower answer.  Since the trace is typed, the first step where promoted content
is discharged into a lower-layer descendant is one of the declared bridge or
audit steps.  That step is a boundary node, so \(\Bd(\mathfrak T)\) is nonempty.
This lemma is tracked by the formalization harness as a projection over the
strict-useful-trace and boundary predicates.
\end{proof}

The next subsection turns this existence statement into the structural
necessity theorem: the boundary contains a non-descending support and an audited
descent expression.

\subsection{Essential boundary necessity}
\label{subsec:necessity-essential-boundary}

\Cref{lem:sau-boundary-existence} says that the boundary set is nonempty.  The
next step is to identify what a boundary node carries: hidden promoted material
upstream, an audit or bridge step at the boundary, and a lower-visible
descendant that can continue toward the final answer.

\begin{definition}[Audit boundary]
\label{def:sau-audit-boundary}
Let \(\mathfrak T\) be an accepted trace.  A node \(b\in V\) is an audit-boundary node when:
\begin{enumerate}
\item \(b\in \Vis_0(\mathfrak T)\);
\item \(b\in\Anc(v_\ast)\), so it contributes to the final answer;
\item at least one hidden promoted node lies upstream of \(b\):
\[
\Hid_1(\mathfrak T)\cap \Anc(b)\neq\varnothing;
\]
\item every incoming dependence from hidden promoted content to the lower-visible status of \(b\) is discharged by a declared audit or bridge in \(\mathcal B\).
\end{enumerate}
Write \(\Bd(\mathfrak T)\) for the finite set of such boundary nodes.
\end{definition}

Thus a boundary node \(b\) is lower-visible, lies on a path to the final answer,
has hidden promoted ancestors, and is precisely where the trace records that
the dependence on hidden content has been discharged by an admissible audit or
bridge.

\begin{definition}[First boundary and hidden support]
\label{def:sau-first-boundary-support}
A first boundary for a hidden node \(h\in\Hid_\ast(\mathfrak T)\) is a boundary node \(b\in\Bd(\mathfrak T)\cap\Anc(v_\ast)\) such that \(h\leadsto b\) and no lower-visible node strictly between \(h\) and \(b\) on the selected dependency path is already fed by \(h\).  Because \(D\) is finite and \(v_\ast\) is lower-visible in an accepted trace, at least one such first boundary exists for every essential hidden node.

The hidden support of \(b\), written \(H_b\), is the finite promoted subtrace consisting of all hidden promoted ancestors of \(b\) that feed \(b\) without passing through an earlier lower-visible audit-boundary node.  The boundary expression of \(H_b\), written \(C_b(H_b)\), is the object carried by \(b\) together with the audited rule data that turns \(H_b\) into \(o(b)\).
\end{definition}

The first boundary for a hidden node is the earliest lower-visible boundary
descendant on a selected dependency path.  Its hidden support \(H_b\) is the
connected promoted subtrace feeding that boundary before any earlier audit
boundary intervenes, and \(C_b(H_b)\) is the expression extracted there by the
audit data.

The necessity theorem is the dependency-cone version of this reading.  Start
with an essential hidden promoted node, move forward until its first boundary,
and read the boundary audit as the descended expression.

\begin{theorem}[Essential Boundary Necessity]
\label{thm:sau-essential-boundary-necessity}
Let \(\mathfrak T\) be a finite strict useful trace for a saturated lower problem profile \(\mathfrak P_0\).  Then there exist a hidden promoted subtrace \(H\), an audited boundary expression \(C_H(H)\), and a lower-visible boundary answer \(b^\sharp\) such that
\[
H\not\downarrow_{\mathcal T_0},
\qquad
C_H(H)\downarrow_B b^\sharp.
\]
Moreover, \(b^\sharp\) lies in the lower dependency cone of the final answer; hence there is a lower-visible audited continuation from \(b^\sharp\) to the final answer \(a^\sharp_{\mathfrak T}\).
\end{theorem}

\begin{proof}
By \Cref{lem:sau-boundary-existence}, the boundary set \(\Bd(\mathfrak T)\) is
nonempty.  Since \(\mathfrak T\) is strict useful, \(\EssProm(\mathfrak T)\)
also holds, so some hidden promoted node \(h\in\Hid_1(\mathfrak T)\) contributes
nontrivially to the final accepted answer \(a^\sharp_{\mathfrak T}\).  The graph
\(D\) is finite, and \(v_\ast\) is lower-visible in an accepted trace; therefore
there is a first audit-boundary descendant \(b\) of \(h\) on a dependency path
to \(v_\ast\).

By \Cref{def:sau-first-boundary-support}, this \(b\) determines a hidden support
\(H_b\) and a boundary expression \(C_b(H_b)\).  By
\Cref{def:sau-audit-boundary}, every dependence of \(b\) on hidden promoted
content is discharged by a declared audit or bridge in \(\mathcal B\).  Since
the audit record \(A_{\mathfrak T}\) passes in an accepted trace, the expression
carried at \(b\) has a typed audit chain to a lower-visible answer; call that
answer \(b^\sharp\).  This is the asserted descent
\[
C_b(H_b)\downarrow_B b^\sharp.
\]

It remains to identify the non-descending support.  If the hidden support
\(H_b\) already descended through the lower theory \(\mathcal T_0\), then the
essential promoted dependence feeding \(b\) could be replaced by lower-layer
data before the boundary.  That would contradict trace strictness and essential
promotion: the promoted content would no longer be essential to the accepted
lower answer.  Hence \(H_b\not\downarrow_{\mathcal T_0}\).  Taking
\[
H=H_b,
\qquad
C_H(H)=C_b(H_b),
\]
gives the desired pair.  Finally, because \(b\in\Anc(v_\ast)\), the descended
answer \(b^\sharp\) lies in the lower dependency cone of the final answer, and
the remaining trace from \(b^\sharp\) to \(a^\sharp_{\mathfrak T}\) is
lower-visible after all intervening boundary dependencies have been discharged.
This theorem is tracked by the formalization harness as a projection over the
strict-useful-trace, audit-boundary, and hidden-support predicates.
\end{proof}

The next subsection extracts a normal form from the triple
\((H,C_H,b^\sharp)\) and identifies the primitive head of the boundary
expression.

\subsection{SAU boundary normal form}
\label{subsec:necessity-boundary-normal-form}

The necessity theorem gives an existence statement.  The normal-form step
compresses that data into a canonical shape: contract the hidden support to a
single node \(U_H\), retain the first audited boundary expression
\(C_H(U_H)\downarrow_B b^\sharp\), and append the lower-visible verifier from
\(b^\sharp\) to \(a^\sharp_{\mathfrak T}\).

\begin{definition}[SAU boundary normal form]
\label{def:sau-boundary-normal-form}
An accepted trace is in SAU boundary normal form when, after deleting irrelevant nodes outside the dependency cone of the final answer, it has the schematic shape
\[
\mathcal T_0
\xleftarrow{\ q\ }
\mathcal T_1,
\qquad
\left[U\not\downarrow_q,
\quad
C(U)\downarrow_B a^\sharp\right],
\]
possibly followed by a lower-visible verifier from \(a^\sharp\) to the final accepted answer.  In a multi-generator case this becomes
\[
\left[\vec U\not\downarrow_q,
\quad
\vec C(\vec U)\downarrow_B \vec a^\sharp\right].
\]
Write
\[
\NF_{\SAU}(\mathfrak T)
\]
for the boundary normal form obtained by contracting hidden promoted support and retaining the first audited lower-visible boundary.
\end{definition}

The schematic shape factors the trace into a single non-descending hidden
support, a single audited boundary expression, and a lower-visible continuation.
It is a compressed trace, not a new derivation.

\begin{theorem}[SAU Boundary Normal Form]
\label{thm:sau-boundary-normal-form}
Every finite strict useful trace \(\mathfrak T\) is audit-equivalent to a trace in SAU boundary normal form.  More precisely, there exist a hidden support node \(U_H\), an audited expression former \(C_H\), and a lower-visible boundary object \(b^\sharp\) such that
\[
\mathfrak T\simeq_{\mathrm{audit}}
\NF_{\SAU}(\mathfrak T)
=
\left[U_H\not\downarrow_{\mathcal T_0},
\quad
C_H(U_H)\downarrow_B b^\sharp
\quad\leadsto\quad
 a^\sharp_{\mathfrak T}\right].
\]
If the chosen boundary is already the final answer node, then \(b^\sharp=a^\sharp_{\mathfrak T}\).
\end{theorem}

\begin{proof}
Apply \Cref{thm:sau-essential-boundary-necessity} to obtain a hidden promoted
subtrace \(H\), an audited boundary expression \(C_H(H)\), and a lower-visible
boundary answer \(b^\sharp\) with
\[
H\not\downarrow_{\mathcal T_0},
\qquad
C_H(H)\downarrow_B b^\sharp.
\]
Now contract the hidden support cone of the chosen first boundary to a single
node \(U_H\).  This contraction does not alter the lower-visible answer or the
audit replay: it simply packages the upstream hidden promoted subtrace as one
support node.  The audit chain from that support to \(b^\sharp\) is retained as
the single boundary expression former \(C_H\).

Delete nodes outside the dependency cone of \(v_\ast\), and contract the
downstream lower-visible audited path from \(b^\sharp\) to
\(a^\sharp_{\mathfrak T}\) into the final verifier.  These operations preserve
the accepted lower answer, the bridge and audit records, and the nonclaim
register, so the compressed trace is audit-equivalent to the original trace.
This is exactly \(\NF_{\SAU}(\mathfrak T)\).  If the chosen boundary is already
the final node \(v_\ast\), then the boundary answer is the final answer:
\(b^\sharp=a^\sharp_{\mathfrak T}\).  The theorem is tracked by the
formalization harness as a projection over the audit-equivalence relation and
the boundary structure.
\end{proof}

The normal form is the data from which an SAU certificate is assembled.  The
next definition packages the inherited witnesses from the trace without adding
new axioms.

\begin{definition}[Extracted SAU package]
\label{def:sau-extracted-package}
Let \(\mathfrak T\) be a strict useful trace and let \(\NF_{\SAU}(\mathfrak T)\) be a chosen boundary normal form
\[
\left[U_H\not\downarrow_{\mathcal T_0},
\quad
C_H(U_H)\downarrow_B b^\sharp
\quad\leadsto\quad
 a^\sharp_{\mathfrak T}\right].
\]
The extracted raw SAU package \(\mathfrak C(\mathfrak T)\) is the raw proof package whose fields are:
\begin{align*}
\mathsf{SatW}&:=\operatorname{Sat}(\mathfrak P_0),\\
\mathsf{StrictW}&:=\delta_{\mathrm{fact}}(\pi_0,\pi_1)\neq\varnothing,\\
\mathsf{NonDescW}&:=U_H\not\downarrow_{\mathcal T_0},\\
\mathsf{DescW}&:=C_H(U_H)\downarrow_B b^\sharp,\\
\mathsf{SolveW}&:=\text{the accepted lower continuation from }b^\sharp\text{ to }a^\sharp_{\mathfrak T},\\
\mathsf{AuditW}&:=\text{the audit, provenance, bridge, and nonclaim records inherited from }\mathfrak T.
\end{align*}
When the boundary node is itself the final answer node, \(b^\sharp=a^\sharp_{\mathfrak T}\); otherwise \(b^\sharp\) is the descended boundary datum and the final lower verifier is included in \(\mathsf{SolveW}\).
\end{definition}

The fields \(\mathsf{SatW}\), \(\mathsf{StrictW}\), \(\mathsf{NonDescW}\),
\(\mathsf{DescW}\), \(\mathsf{SolveW}\), and \(\mathsf{AuditW}\) are the six
certificate witnesses realized as inherited trace data.  The extraction is by
compression and witness collection, not by an axiom asserting usefulness.

The extracted raw package still owes a converse: when does such a package
validate to an SAU certificate, and does that close the tautology objection?
The next subsection gives the partial converse and the tautology barrier.

\subsection{Partial converse and the tautology barrier}
\label{subsec:necessity-partial-converse-tautology-barrier}

\Cref{thm:sau-essential-boundary-necessity} and
\Cref{thm:sau-boundary-normal-form} produce an extracted package
\(\mathfrak C(\mathfrak T)\) from a strict useful trace.  The remaining
anti-tautology question is whether that extracted package validates when its
inherited audits pass.  The partial converse answers yes: the package is not
postulated as useful, but validated from the witnesses already present in the
trace.

\begin{theorem}[Partial converse to SAU]
\label{thm:sau-partial-converse}
Let \(\mathfrak T\) be a finite strict useful trace with explicit provenance and declared audit boundaries.  Then there exists an extracted raw SAU proof package \(\mathfrak C(\mathfrak T)\).  Moreover, if the inherited audits and lower verifier pass under the declared instruments, then
\[
\ValSAU(\mathfrak C(\mathfrak T))
\]
and hence, by the soundness theorem of \Cref{sec:soundness-failure},
\[
\Use(B_{0\to1},P):\texttt{accepted}.
\]
\end{theorem}

\begin{proof}
The normal-form construction gives the raw package
\(\mathfrak C(\mathfrak T)\) with the six witness fields listed in
\Cref{def:sau-extracted-package}.  Saturation, strictness, non-descent,
descent, lower solving, and audit data are inherited from the original strict
useful trace and its boundary normal form.  If the inherited audits and the
lower verifier pass under the declared instruments, those fields validate as an
SAU proof package.  The soundness theorem of
\Cref{sec:soundness-failure}, referenced in the statement, then licenses the
accepted lower-layer usefulness claim.  This theorem is tracked by the
formalization harness as a projection over the extracted-package construction
and the soundness predicate.
\end{proof}

Combining the existence theorem, the normal form, and the partial converse gives
the advertised round trip: strict useful trace, extracted normal form,
validated certificate, accepted lower answer.

\begin{theorem}[Necessity--soundness chain]
\label{thm:sau-necessity-soundness-chain}
For finite typed audited traces, the following implication chain holds:
\[
\operatorname{StrictUsefulTrace}(\mathfrak T,\mathfrak P_0)
\Longrightarrow
\exists\mathfrak C(\mathfrak T)\text{ with }\ValSAU(\mathfrak C(\mathfrak T))
\Longrightarrow
\Use(B_{0\to1},P):\texttt{accepted},
\]
provided the extracted package passes the inherited audit checks.
Equivalently, within this scope, strict useful mathematics has the normal form
\[
\boxed{
\text{non-descending promoted support}
\quad+\quad
\text{audited descending boundary}
}
\]
up to audit-equivalent compression.
\end{theorem}

\begin{proof}
The first implication is the conjunction of
\Cref{lem:sau-boundary-existence},
\Cref{thm:sau-essential-boundary-necessity}, and
\Cref{thm:sau-boundary-normal-form}: a strict useful trace has a boundary,
that boundary carries a non-descending support and an audited descent
expression, and the trace is audit-equivalent to the corresponding normal
form.  The extracted-package definition then supplies
\(\mathfrak C(\mathfrak T)\).  The second implication is
\Cref{thm:sau-partial-converse}, which validates the extracted package under
the inherited audit checks and then applies soundness.  The boxed expression is
a summary of the normal-form factorization.  This theorem is tracked by the
formalization harness as a projection over the boundary-existence,
normal-form, extracted-package, and soundness predicates.
\end{proof}

The chain is the section's anti-tautology point.  It says that SAU-shaped data
are forced from accepted strict promoted use under the stated finite typed
audit hypotheses, not merely selected as a convenient definition.

\begin{corollary}[Tautology barrier]
\label{cor:sau-tautology-barrier}
The SAU framework is not only the sufficient condition
\[
\SAU\Rightarrow\Use.
\]
For finite typed audited strict-use traces, the reverse extraction also holds:
\[
\operatorname{StrictUsefulTrace}
\Rightarrow
\exists\SAU\text{-shaped certificate}.
\]
Thus SAU is a normal form for strict audited promoted use, not merely a name for useful extensions.
\end{corollary}

The corollary reads SAU as a normal-form description of the trace, not as a
definition that grants its own conclusion.  It is tracked by the formalization
harness as a projection over the necessity theorem, the normal form, and the
partial converse.

With necessity, normal form, and the tautology barrier in place, the next
section turns to soundness and failure-mode discipline: what a validated SAU
certificate licenses, and what it blocks.
 \section{SAU Soundness and Failure Modes}\label{sec:soundness-failure}

\subsection{Raw and validated proof packages}
\label{subsec:soundness-raw-validated}

\Cref{sec:layers} defined an SAU certificate as a six-field record, and
\Cref{sec:necessity} explained how such a package can be extracted from a
strict useful trace.  We now separate the candidate record from the checked
record: a raw proof package records the six obligations and proposed witnesses,
while a validated proof package is the same record after the declared
instruments have checked those witnesses.

\begin{definition}[Raw SAU proof package]
\label{def:sau-raw-proof-package}
A raw Strict Audited Utility proof package is a finite record
\[
\mathfrak C=(\mathfrak P_0,B_{0\to1},U,C,a^\sharp;
\mathsf{SatW},\mathsf{StrictW},\mathsf{NonDescW},\mathsf{DescW},\mathsf{SolveW},\mathsf{AuditW})
\]
whose six witness fields are intended to prove the following six obligations:
\begin{align*}
\mathsf{SatW}&:\operatorname{Sat}(\mathfrak P_0),\\
\mathsf{StrictW}&:\delta_{\mathrm{fact}}(\pi_0,\pi_1)\neq\varnothing,\\
\mathsf{NonDescW}&:U\not\downarrow_q,\\
\mathsf{DescW}&:C(U)\downarrow_B a^\sharp,\\
\mathsf{SolveW}&:a^\sharp\in\mathsf{Ans}_0(P)
\text{ and }\mathsf{Eval}_{\Theta_0}(\delta_0(a^\sharp))=\texttt{pass},\\
\mathsf{AuditW}&:\delta_{\mathrm{overread}}=\varnothing,
\quad\delta_{\mathrm{nosmuggle}}(B)=\varnothing,
\quad\delta_{\mathrm{audit}}(B)=\varnothing,
\quad\mathsf{AuditGate}_B:\texttt{pass},
\end{align*}
with an explicit nonclaim register.
\end{definition}

The six fields \(\mathsf{SatW}\), \(\mathsf{StrictW}\),
\(\mathsf{NonDescW}\), \(\mathsf{DescW}\), \(\mathsf{SolveW}\), and
\(\mathsf{AuditW}\) correspond to the six SAU obligations already visible in
the certificate pattern.  At this stage they are only recorded witnesses: a raw
package is a candidate record, not yet a certified usefulness claim.

Validation is the step from candidate to checked package.  The relevant tests
are not informal judgments; they are run through the lower instrument and the
bridge instrument declared for the problem.

\begin{definition}[Validated SAU proof package]
\label{def:sau-validated-proof-package}
A raw proof package \(\mathfrak C\) is validated when every witness field is checked under the declared instruments \(I_0\) and \(I_B\), every bridge used by the proof belongs to the declared admissible bridge family, and every nonclaim required by \(\mathcal N_0\) and \(\mathcal N_B\) is recorded.  We write
\[
\operatorname{ValSAU}(\mathfrak C)
\]
for this finite verification result.
\end{definition}

Thus \(\operatorname{ValSAU}(\mathfrak C)\) means that each witness field has
passed the declared checks, every bridge is admissible, and the nonclaim
register has been recorded.  The next subsection states what such a validated
SAU proof package licenses.

\subsection{SAU soundness}
\label{subsec:soundness-sau-soundness}

With raw and validated proof packages in place, the soundness theorem states
what validation licenses.  The output is an instrument-indexed accepted
usefulness claim: it is a claim made under the declared lower and bridge
instruments, not an instrument-free assertion that the hidden promoted
structure has become lower-layer content.

\begin{theorem}[Strict Audited Utility Soundness]
\label{thm:sau-soundness}
Let \(\mathfrak C\) be a validated raw SAU proof package for
\((\mathfrak P_0,B_{0\to1},U,C,a^\sharp)\).  Then the following instrument-indexed claim is accepted:
\[
\Gamma;\mathcal T_0\leadsto\mathcal T_1;I_B
\vdash
\operatorname{Useful}(B_{0\to1},P):\texttt{accepted}.
\]
More explicitly, the accepted claim has the following content:
\begin{enumerate}
\item \textbf{Lower-layer saturation.} No candidate at the fixed lower interface solves the problem:
\[
\forall s\in\mathsf{Cand}_0(P),
\qquad
\mathsf{Eval}_{\Theta_0}(\delta_0(s))\neq\texttt{pass}.
\]
\item \textbf{Promoted answer.} The promoted calculation supplies a lower-layer accepted answer:
\[
a^\sharp\in\mathsf{Ans}_0(P),
\qquad
\mathsf{Eval}_{\Theta_0}(\delta_0(a^\sharp))=\texttt{pass}.
\]
\item \textbf{Strictness.} The promoted interface is not a relabeling of the lower interface:
\[
\delta_{\mathrm{fact}}(\pi_0,\pi_1)\neq\varnothing,
\]
with explicit witness pair
\[
(s,s')\in\delta_{\mathrm{fact}}(\pi_0,\pi_1).
\]
\item \textbf{Non-descent of the promoted structure.} The promoted structure itself is not accepted as lower-layer content:
\[
U\not\downarrow_q.
\]
\item \textbf{Audited descent of the selected expression.} The expression selected from the promoted structure descends under the bridge:
\[
C(U)\downarrow_B a^\sharp.
\]
\item \textbf{No overreading.} The accepted claim asserts only the audited lower descendant \(a^\sharp\), not the whole hidden promoted structure \(U\), at the lower layer.
\end{enumerate}
Consequently, \(B_{0\to1}\) is a strictly useful promoted extension of \(\mathcal T_0\) for \(P\) in the sense recorded by the canonical useful-extension definition in \Cref{app:definitions}.
\end{theorem}

\begin{proof}
Each component of the accepted claim is read from the corresponding witness
field of the validated package.  The saturation clause is \(\mathsf{SatW}\);
the accepted lower answer is \(\mathsf{SolveW}\); strictness is
\(\mathsf{StrictW}\); non-descent is \(\mathsf{NonDescW}\); audited descent is
\(\mathsf{DescW}\); and the no-overreading and bridge checks are part of
\(\mathsf{AuditW}\).  Validation says that these witnesses have been checked
under the declared instruments, so the displayed sequent is just the packaged
instrument-indexed form of those six accepted items.  This theorem is tracked
by the formalization harness as a projection over the validated-package
witnesses.
\end{proof}

The theorem licenses an accepted lower-layer claim about \(a^\sharp\).  It does
not license the stronger claim that the hidden promoted structure \(U\) belongs
to the lower layer.

\begin{corollary}[No collapse of hidden structure]
\label{cor:sau-no-collapse}
Under the hypotheses of Theorem~\ref{thm:sau-soundness}, the accepted usefulness claim does not entail
\[
U\in\mathcal T_0
\]
or any equivalent lower-layer literalization of the promoted structure, unless that stronger claim is separately bridged and audited.
\end{corollary}

\begin{proof}
The soundness theorem licenses only the audited descended expression and the
accepted lower answer.  A claim that \(U\) itself is lower-layer content is a
separate literalization claim and is not supplied by the validated package
unless an additional bridge and audit establish it.  This corollary is tracked
by the formalization harness as a projection over the validated-package
witnesses and the no-overreading clause.
\end{proof}

Strictness blocks a second overstatement.  The promoted layer may be useful, but
the strictness witness says it is not merely a notational renaming of the lower
interface.

\begin{corollary}[No reduction to notation]
\label{cor:sau-no-notation}
Under the hypotheses of Theorem~\ref{thm:sau-soundness}, the promoted layer is not merely a notational variant of the lower layer at the declared interface.
\end{corollary}

\begin{proof}
The validated package includes the strictness witness
\(\delta_{\mathrm{fact}}(\pi_0,\pi_1)\neq\varnothing\).  Thus two states can
agree at the lower interface while differing at the promoted interface, so the
promotion is not just a change of notation.  This corollary is tracked by the
formalization harness as a projection over the strictness witness in the
validated package.
\end{proof}

The usefulness is also nontrivial in the fixed problem profile.  Saturation says
that the lower candidates did not already pass, while audited descent supplies
an accepted lower answer.

\begin{corollary}[Nontrivial usefulness]
\label{cor:sau-nontrivial-usefulness}
Under the hypotheses of Theorem~\ref{thm:sau-soundness}, the promoted layer supplies an accepted lower answer that the fixed lower interface does not supply.
\end{corollary}

\begin{proof}
The saturation witness rules out a passing lower candidate at the fixed lower
interface, and the solve witness supplies the accepted lower answer
\(a^\sharp\).  The audited descent witness explains how that answer is obtained
from the promoted computation.  This corollary is tracked by the formalization
harness as a projection over the saturation, descent, and solve witnesses.
\end{proof}

When validation fails, the next question is where it fails.  The next subsection
introduces the failure-field machinery and the diagnostic classifier.

\subsection{Failure fields and the diagnostic classifier}
\label{subsec:soundness-failure-classifier}

Soundness is useful only if failed validation can be localized.  The six failure
fields record exactly which SAU obligation is missing: saturation, strictness,
non-descent, descent, solving, or audit safety.

\begin{definition}[SAU failure fields]
\label{def:sau-failure-fields}
For a raw candidate package
\[
\mathfrak C=(\mathfrak P_0,B_{0\to1},U,C,a^\sharp;
\mathsf{SatW},\mathsf{StrictW},\mathsf{NonDescW},\mathsf{DescW},\mathsf{SolveW},\mathsf{AuditW}),
\]
define the six SAU failure fields as follows:
\[
F_{\mathrm{sat}}(\mathfrak C)=
\begin{cases}
\varnothing,&\operatorname{Sat}(\mathfrak P_0),\\
\{\star\},&\text{otherwise},
\end{cases}
\]
\[
F_{\mathrm{strict}}(\mathfrak C)=
\begin{cases}
\varnothing,&\delta_{\mathrm{fact}}(\pi_0,\pi_1)\neq\varnothing,\\
\{\star\},&\delta_{\mathrm{fact}}(\pi_0,\pi_1)=\varnothing,
\end{cases}
\]
\[
F_{\mathrm{nondesc}}(\mathfrak C)=
\begin{cases}
\varnothing,&U\not\downarrow_q,\\
\{\star\},&\text{otherwise},
\end{cases}
\]
\[
F_{\mathrm{desc}}(\mathfrak C)=
\begin{cases}
\varnothing,&C(U)\downarrow_B a^\sharp,\\
\{\star\},&\text{otherwise},
\end{cases}
\]
\[
F_{\mathrm{solve}}(\mathfrak C)=
\begin{cases}
\varnothing,&a^\sharp\in\mathsf{Ans}_0(P)
\text{ and }\mathsf{Eval}_{\Theta_0}(\delta_0(a^\sharp))=\texttt{pass},\\
\{\star\},&\text{otherwise},
\end{cases}
\]
and
\[
F_{\mathrm{audit}}(\mathfrak C)=
\delta_{\mathrm{overread}}\cup \delta_{\mathrm{nosmuggle}}(B)
\cup F_{6}^{\mathrm{audit}}(\mathfrak C),
\]
where
\[
F_{6}^{\mathrm{audit}}(\mathfrak C)=
\begin{cases}
\varnothing,&\delta_{\mathrm{audit}}(B)=\varnothing
\text{ and }\mathsf{AuditGate}_B:\texttt{pass},\\
\{\star\},&\text{otherwise}.
\end{cases}
\]
The SAU defect vector is
\[
\mathbf F_{\SAU}(\mathfrak C)
=
(F_{\mathrm{sat}},F_{\mathrm{strict}},F_{\mathrm{nondesc}},
F_{\mathrm{desc}},F_{\mathrm{solve}},F_{\mathrm{audit}}).
\]
We write
\[
\mathbf F_{\SAU}(\mathfrak C)=\varnothing
\]
when all six components are empty.
\end{definition}

The six failure modes are summarized in \Cref{tab:failure-modes}.  The table is
prose-facing; the displayed definition above is the formal diagnostic vector.

\begingroup\begin{table}[tbp]
\centering
\footnotesize
\caption{SAU failure modes, blocked claims, and repair routes.}
\label{tab:failure-modes}
\begin{tabularx}{\textwidth}{@{}>{\raggedright\arraybackslash}p{0.16\textwidth}>{\raggedright\arraybackslash}X>{\raggedright\arraybackslash}X>{\raggedright\arraybackslash}X@{}}
\toprule
Failure mode & Missing gate & Blocked claim & Repair route \\
\midrule
Saturation failure &
The lower problem has not been shown saturated at the declared defect threshold. &
The promoted extension is necessary for the fixed lower problem. &
Tighten the lower problem, defect, or threshold, or supply the missing saturation argument. \\
Strictness failure &
The promoted interface has not been shown to distinguish states collapsed by the lower view. &
The promotion is a strict extension rather than a relabeling. &
Choose a comparison carrier and lens that separate promoted states the lower interface identifies. \\
Hidden-structure failure &
The hidden promoted structure is accepted by the lower interface. &
Non-descent of the promoted support. &
Replace \(U\) by a genuinely non-descending generator, or change the lens and bridge. \\
Descent failure &
The selected expression \(C(U)\) has not been shown to descend through the bridge. &
The audited descent \(C(U)\downarrow_B a^\sharp\). &
Choose a different expression former or a declared audit channel that does descend. \\
Solve failure &
The descended answer \(a^\sharp\) does not pass the lower verifier. &
Acceptance of \(a^\sharp\) as the lower answer. &
Re-examine \(a^\sharp\), adjust the lower verifier, or supply a stronger solving argument. \\
Audit failure &
An overreading, smuggling, provenance, or replay check remains open. &
Soundness, no-collapse, no-notation, and no-overreading. &
Repair the bridge audit by recording the nonclaim register, source policy, and replay data. \\
\bottomrule
\end{tabularx}
\end{table}
 \endgroup

The diagnostic tag set turns the formal vector into readable labels.  The tags
do not replace the failure fields; they help a reader see why a candidate
package did not validate.

\begin{definition}[SAU diagnostic tag set]
\label{def:sau-diagnostic-tag-set}
The diagnostic tag set of a raw package is the finite set
\[
\operatorname{Diag}_{\SAU}(\mathfrak C)\subseteq
\left\{
\begin{gathered}
\texttt{accepted},\texttt{convenient},\texttt{non\_strict},\texttt{not\_hidden},\\
\texttt{unreturned},\texttt{irrelevant},\texttt{failed\_audit},\\
\texttt{smuggled},\texttt{overread}
\end{gathered}
\right\}
\]
assigned by the following rules:
\begin{enumerate}
\item \(\texttt{accepted}\in\operatorname{Diag}_{\SAU}(\mathfrak C)\) only when the claim status is \texttt{accepted}, \(\mathbf F_{\SAU}(\mathfrak C)=\varnothing\), and the package validates.
\item If \(F_{\mathrm{sat}}\neq\varnothing\), add \(\texttt{convenient}\).  The promotion may still compress or clarify, but it is not certified as necessary for a saturated lower defect.
\item If \(F_{\mathrm{strict}}\neq\varnothing\), add \(\texttt{non\_strict}\).  This mirrors the canonical promotion status \texttt{non\_strict}; it is recorded here only as a diagnostic tag.
\item If \(F_{\mathrm{nondesc}}\neq\varnothing\), add \(\texttt{not\_hidden}\).  The alleged promoted structure is not shown to be genuinely non-descending.
\item If \(F_{\mathrm{desc}}\neq\varnothing\), add \(\texttt{unreturned}\).  The promoted calculation has not produced an accepted lower descendant.
\item If \(F_{\mathrm{solve}}\neq\varnothing\), add \(\texttt{irrelevant}\).  The descendant does not solve, decide, or certify the declared lower problem.
\item If \(F_6^{\mathrm{audit}}\neq\varnothing\), add \(\texttt{failed\_audit}\); the canonical claim status is then \texttt{blocked}.
\item If \(\delta_{\mathrm{nosmuggle}}(B)\neq\varnothing\), add \(\texttt{smuggled}\); this mirrors the canonical promotion status \texttt{smuggled}.
\item If \(\delta_{\mathrm{overread}}\neq\varnothing\), add \(\texttt{overread}\); this mirrors the canonical claim status \texttt{overread}.
\end{enumerate}
These tags are not a new status family.  They are readable labels derived from the canonical claim, promotion, and gate statuses together with the SAU defect vector.
\end{definition}

Some invalid packages are mathematically recognizable even before one reads the
full defect vector.  The counterfeit patterns name those common failures.

\begin{definition}[Counterfeit SAU patterns]
\label{def:sau-counterfeit-patterns}
A candidate bridge has one of the following counterfeit profiles when it resembles SAU informally but lacks a required field:
\begin{enumerate}
\item \textbf{Ornamental extension:}
\[
\operatorname{Strict}(B),\quad U\not\downarrow_q,
\quad\text{but}\quad C(U)\not\downarrow_B a^\sharp.
\]
The promoted layer is richer but unreturned.
\item \textbf{Mere notation:}
\[
C(U)\downarrow_B a^\sharp,
\quad\text{but}\quad
\delta_{\mathrm{fact}}(\pi_0,\pi_1)=\varnothing.
\]
The promoted expression descends, but no strictness has been shown.
\item \textbf{Irrelevant descent:}
\[
C(U)\downarrow_B a^\sharp,
\quad\text{but}\quad
\mathsf{Eval}_{\Theta_0}(\delta_0(a^\sharp))\neq\texttt{pass}.
\]
The descendant exists but does not answer the problem.
\item \textbf{Smuggled proof:}
\[
\delta_{\mathrm{nosmuggle}}(B)\neq\varnothing.
\]
The proof uses content not admitted by the bridge.
\item \textbf{Literalized hidden object:}
\[
U\not\downarrow_q
\quad\text{and}\quad
\Lit_0(U)
\]
without an independent bridge for \(U\).
\end{enumerate}
\end{definition}

The classifier records the immediate consequence of each failed witness.  It is
not a new proof obligation; it is the readable form of the six failure fields.

\begin{proposition}[SAU failure classifier]
\label{prop:sau-failure-classifier}
Let \(\mathfrak C\) be a raw SAU proof package.  If one of the six witness fields fails, then the strict audited utility claim is not accepted at SAU strength.  More specifically:
\begin{enumerate}
\item If \(\operatorname{Sat}(\mathfrak P_0)\) fails, the SAU-strength claim status is not \texttt{accepted}; any convenience or compression reading must be recorded outside the claim-status family.
\item If \(\delta_{\mathrm{fact}}(\pi_0,\pi_1)=\varnothing\), the promotion status is \texttt{non\_strict} at the declared interface.
\item If \(U\not\downarrow_q\) fails, the alleged hidden promoted structure was already lower-layer accepted or was not shown to be hidden.
\item If \(C(U)\downarrow_B a^\sharp\) fails, the promoted calculation does not return an accepted lower descendant.
\item If the problem-resolution field fails, the descendant does not solve, decide, or certify the lower problem.
\item If the audit/no-overread/no-smuggling field fails, the claim status is \texttt{blocked} or \texttt{outside\_scope}; a no-smuggling failure may also give the promotion status \texttt{smuggled}.
\end{enumerate}
\end{proposition}

\begin{proof}
Each clause is obtained by reading the corresponding component of
\(\mathbf F_{\SAU}(\mathfrak C)\) and translating the failed witness into the
claim-status or promotion-status consequence recorded by the diagnostic tag
set.  This proposition is tracked by the formalization harness as a projection
over the six failure-field components and the diagnostic classifier.
\end{proof}

The acceptance boundary says that validation has exactly one entry condition:
all six components of the failure vector must be empty under the declared
instruments.

\begin{theorem}[Acceptance boundary for SAU]
\label{thm:sau-acceptance-boundary}
A raw package \(\mathfrak C\) can be validated as an SAU proof package only if
\[
\mathbf F_{\SAU}(\mathfrak C)=\varnothing.
\]
Conversely, if the six witness fields are verified under the declared instruments and
\[
\mathbf F_{\SAU}(\mathfrak C)=\varnothing,
\]
then \(\mathfrak C\) validates to an SAU proof package:
\[
\operatorname{ValSAU}(\mathfrak C).
\]
\end{theorem}

\begin{proof}
The forward direction is immediate from validation: a validated package has
passed each of the six witness checks, so each failure component is empty.  For
the converse, the six verified witness fields supply saturation, strictness,
non-descent, descent, solving, and audit safety, while
\(\mathbf F_{\SAU}(\mathfrak C)=\varnothing\) records that none of the
corresponding failure conditions remains open.  Hence the package validates.
This theorem is tracked by the formalization harness as a projection over the
validated-package predicate and the six failure fields.
\end{proof}

The boundary immediately explains why one nonempty component is enough to block
SAU-strength acceptance.

\begin{corollary}[Failure blocks SAU-strength acceptance]
\label{cor:sau-failure-blocks-acceptance}
If at least one component of \(\mathbf F_{\SAU}(\mathfrak C)\) is nonempty, then the strict audited utility claim
\[
\Gamma;\mathcal T_0\leadsto\mathcal T_1;I_B\vdash
\Use(B_{0\to1},P):\texttt{accepted}
\]
is not licensed by Theorem~\ref{thm:sau-soundness}.
\end{corollary}

\begin{proof}
If a component of the failure vector is nonempty, then
\(\mathbf F_{\SAU}(\mathfrak C)\neq\varnothing\).  By
\Cref{thm:sau-acceptance-boundary}, the package cannot validate as an SAU proof
package, so \Cref{thm:sau-soundness} does not license the accepted usefulness
claim.  This corollary is tracked by the formalization harness as a projection
over the acceptance-boundary theorem and the soundness predicate.
\end{proof}

Finally, the exact-six role vector records the same idea at the inherited
six-role level: a role is usable only when it is present, aligned, thresholded,
visible or bridged, audited, and nonclaim-safe.

\begin{definition}[Exact-six role failure vector]
\label{def:sau-exact-six-failure-vector}
For a candidate exact-six upgrade with role witnesses
\[
(W_1,W_2,W_3,W_4,W_5,W_6),
\]
define
\[
\mathbf R_6=(R_1,R_2,R_3,R_4,R_5,R_6),
\]
where \(R_i=\varnothing\) exactly when \(W_i\) is present, role-aligned, thresholded, level-indexed, visible or bridged, audited, and nonclaim-safe; otherwise \(R_i=\{\star\}\).
\end{definition}

The exact-six vector is the role-level analogue of the SAU defect vector.  The
next subsection uses these diagnostics to state the no-overreading boundary for
literalization claims.

\subsection{No-overreading and lower-layer literalization}
\label{subsec:soundness-no-overreading}

The soundness theorem licenses an accepted lower claim about the audited
descendant \(a^\sharp\).  The no-overreading discipline is the corresponding
negative rule: the same certificate does not license the stronger claim that the
hidden promoted structure \(U\) itself is a lower-layer object or operation.

\begin{definition}[Lower-layer literalization claim]
\label{def:sau-literalization-claim}
A lower-layer literalization claim for a promoted structure \(U\), written
\[
\Lit_0(U),
\]
is a claim that asserts one of the following at the lower layer without an additional admissible bridge for \(U\) itself:
\begin{enumerate}
\item \(U\in\mathcal T_0\);
\item \(U\downarrow_q\);
\item \(U\) is directly visible under \(I_0\);
\item the lower layer may use all of \(U\), rather than only the audited descendant \(C(U)\).
\end{enumerate}
A literalization claim is bridge-backed only when the bridge family \(\mathcal L_B\) contains an admissible bridge for the literalized content of \(U\), not merely for \(C(U)\).
\end{definition}

The four listed literalization claims are the main overreadings that a useful
promotion invites.  They confuse an audited descendant with unrestricted
membership or direct visibility of the hidden support.

\begin{theorem}[No-overreading of hidden promoted structure]
\label{thm:sau-no-overreading-hidden-structure}
Assume a bridge satisfies the SAU non-descent field
\[
U\not\downarrow_q
\]
and the audited descent field
\[
C(U)\downarrow_B a^\sharp.
\]
Then the lower-layer literalization claim \(\Lit_0(U)\) is blocked unless an additional admissible bridge for \(U\) itself is supplied.  In particular, the audited descent of \(C(U)\) does not by itself license \(U\in\mathcal T_0\) or \(U\downarrow_q\).
\end{theorem}

\begin{proof}
By audited descent, the certificate supplies a bridge-backed derivation from the
selected expression \(C(U)\) to the lower answer \(a^\sharp\).  That audit
discharges the selected expression only.  A literalization claim
\(\Lit_0(U)\) asserts a stronger statement about \(U\) itself: lower-layer
membership, direct visibility, descent of \(U\), or unrestricted lower use of
all of \(U\).  Such a statement requires a separate admissible bridge in
\(\mathcal L_B\) for the literalized content of \(U\).  Without that bridge, the
claim is not licensed by the audited descent of \(C(U)\).  This theorem is
tracked by the formalization harness as a projection over the no-overread defect
being empty.
\end{proof}

The four case profiles instantiate the same rule in different mathematical
forms.  In the composite case of \Cref{sec:imaginary}, the audited descent of
\(G^2\) does not license \(G\downarrow_q\) or \(i\) as a lower-layer object; see
\Cref{subsec:imag-certificate-completion-no-overreading}.  In the coefficient case of
\Cref{sec:calculus}, the descended derivative does not license a lower-layer
claim that the infinitesimal \(\varepsilon\) is a real scalar.  In the invariant
case of \Cref{sec:galois}, trace and norm descent do not license membership of
a non-base-field root in the base field.  In the obstruction case of
\Cref{sec:cohomology}, the audited cohomology class does not license exactness
or membership of the hidden cochain in the original lower ledger.

This completes the soundness and failure-mode discipline: a validated
certificate licenses the audited lower answer, while the failure fields and
no-overreading theorem explain what is blocked.  The next section turns to
classification, where the four primitive head types and the basic profile
registry organize the possible audited boundary expressions.
 \section{Basic Profile Classification}\label{sec:classification}

\subsection{Reduced boundary traces and the basic audit grammar}
\label{subsec:classification-reduced-grammar}

The four case studies cluster because their audited boundary expressions have
typed outer shapes.  This section isolates the head constructor and support of a
boundary expression, then records the finite registry of primitive head types.
The case studies are therefore not loose analogies: they instantiate distinct
primitive heads in a shared grammar.

We first keep only the part of a strict useful trace that matters at the audited
boundary.

\begin{definition}[Reduced SAU boundary trace]
\label{def:sau-reduced-boundary-trace}
A reduced SAU boundary trace is a normal-form trace
\[
\mathfrak B=\bigl(\mathfrak P_0,B_{0\to1},U,C,a^\sharp;A_{\mathfrak B},\mathcal N\bigr)
\]
obtained from a finite strict useful trace by the boundary-normal-form theorem of \Cref{sec:necessity}, such that:
\begin{enumerate}
\item \(U\) is the hidden support subtrace or generated hidden support of the promoted calculation;
\item \(U\not\downarrow_q\), so \(U\) itself is not a lower-layer object, operation, or claim;
\item \(C(U)\downarrow_B a^\sharp\), so the selected boundary expression descends under audit;
\item no proper audited subexpression \(D(U)\) of \(C(U)\) already supplies the same lower answer with the same proof role, unless the trace has first been replaced by that smaller boundary;
\item the nonclaim record states that only \(C(U)\)'s descendant is asserted at the lower layer, not \(U\) itself.
\end{enumerate}
We write
\[
\RedSAU(\mathfrak B)
\]
for this condition.
\end{definition}

The head constructor is the outer operation that makes the hidden support
auditable.

\begin{definition}[Boundary constructor and head]
\label{def:sau-boundary-head}
Let \(C(U)\) be the selected expression in a reduced boundary trace.  Its head constructor is the outermost audit-relevant operation used to turn the hidden support into a lower-descending expression.  It is denoted
\[
\Head(C).
\]
The head is typed.  For example, a power constructor \(U\mapsto U^n\), a coefficient constructor \(F(\eta)\mapsto[\eta]F(\eta)\), a fixed-point constructor \(x\mapsto x^G\), and a quotient-class constructor \(a\mapsto[a]\) have different sorts even if their descended values happen to be represented by the same set-theoretic object.
\end{definition}

The support records the smallest promoted data actually used by that head, rather
than the whole promoted layer.

\begin{definition}[Boundary support]
\label{def:sau-boundary-support}
The boundary support of \(C(U)\), written
\[
\Supp_{\Boundary}(C(U)),
\]
is the smallest part of the hidden support and promoted package used by the head constructor to produce the audited descendant.  In a root-composite profile this is the orbit of \(U\) under composition or powers; in a coefficient profile it is the first-order thickening or coefficient direction; in an invariant profile it is the symmetry orbit; and in an obstruction profile it is the local datum together with the subobject declared trivial or exact.
\end{definition}

The basic grammar lists the primitive ways in which a hidden support may produce
an audited lower descendant.

\begin{definition}[Basic audit grammar]
\label{def:sau-basic-audit-grammar}
The basic SAU audit grammar \(\mathcal G_{\Basic}\) is the typed grammar whose primitive boundary constructors are the following:
\begin{enumerate}
\item \textbf{Composite or power constructor:}
\[
C(U)=U^n\quad(n\ge2)
\]
or, more generally, a finite composite of a hidden update whose selected composite descends.
\item \textbf{Coefficient or standard-part constructor:}
\[
C(U)=\kappa(\Phi(U)),
\]
where \(\kappa\) is a coefficient, jet, residue, leading-term, or standard-part audit map.
\item \textbf{Projection or shadow constructor:}
\[
C(U)=q(U),
\]
where \(q\) is a declared shadow, real-part, observable, quotient, or forgetful projection.
\item \textbf{Invariant or fixed-point constructor:}
\[
C(U)\in \Inv_G(\mathcal T_1)\quad\text{or}\quad C(U)\in \Fix(G),
\]
where a promoted symmetry action makes the selected expression invariant and hence lower-descending.
\item \textbf{Trace, norm, or symmetric aggregation constructor:}
\[
C(U)=\sum_{g\in G}gU,
\qquad
C(U)=\prod_{g\in G}gU,
\]
or another declared finite symmetric aggregation of a promoted orbit.
\item \textbf{Quotient or obstruction-class constructor:}
\[
C(U)=[U]\in M/N,
\]
where \(N\) is the subobject declared exact, trivial, gauge, globally generated, or lower-null.
\end{enumerate}
The grammar is closed under finite products, finite sums, ordered tuples, and composition with already accepted lower-layer verifiers.  It is not closed under arbitrary informal analogy, undeclared semantic interpretation, or unaudited transfer.
\end{definition}

A boundary trace is basic when its expression is built from that grammar and the
bridge audit has declared the constructors it uses.

\begin{definition}[Basic boundary trace]
\label{def:sau-basic-boundary-trace}
A reduced SAU boundary trace \(\mathfrak B\) is basic when its boundary expression \(C(U)\) is generated by \(\mathcal G_{\Basic}\) and every constructor used in \(C\) is declared in the bridge audit record.  We write
\[
\BasicSAU(\mathfrak B).
\]
\end{definition}

The resulting profile registry names the reusable proof shapes: \(\RootComp\)
for a hidden generator whose iterate or composite descends, \(\Coeff\) for a
coefficient or standard-part audit, \(\ProjShad\) for a declared projection or
shadow, \(\InvDesc\) for invariant descent, \(\AggInv\) for trace, norm, or
symmetric aggregation, and \(\ObsLed\) for quotient or obstruction-ledger
descent.

\begin{definition}[Profile registry generated by the basic grammar]
\label{def:sau-basic-profile-registry}
The basic profile registry is
\[
\mathcal R_{\Basic}
=\{\RootComp,\Coeff,\ProjShad,\InvDesc,\AggInv,\ObsLed\}.
\]
The four canonical case-study profiles are the subregistry
\[
\mathcal R_{4}
=\{\RootComp,\Coeff,\InvDesc,\ObsLed\}.
\]
The registry map
\[
\rho_{\Basic}:\mathcal G_{\Basic}\longrightarrow\mathcal R_{\Basic}
\]
is defined on primitive heads by
\[
\rho_{\Basic}(U\mapsto U^n)=\RootComp,
\]
\[
\rho_{\Basic}(\kappa\circ\Phi)=\Coeff,
\]
\[
\rho_{\Basic}(q)=\ProjShad,
\]
\[
\rho_{\Basic}(\Inv_G\text{ or }\Fix(G))=\InvDesc,
\]
\[
\rho_{\Basic}(\sum_G,\prod_G,\text{symmetric aggregation})=\AggInv,
\]
\[
\rho_{\Basic}(U\mapsto[U]\in M/N)=\ObsLed.
\]
For composite boundary expressions, \(\rho_{\Basic}\) returns the finite ordered list or product of primitive profile labels appearing in the reduced expression.
\end{definition}

With the grammar fixed, the next subsection states the primitive head
classification theorem.

\subsection{Primitive head classification}
\label{subsec:classification-primitive-head}

The basic profile registry assigns one profile label to each primitive head
constructor.  The classification theorem says that this assignment is unique on
reduced basic boundaries: once the head of \(C\) has been extracted, the profile
label is fixed.

\begin{theorem}[Primitive head classification]
\label{thm:sau-primitive-head-classification}
Let \(\mathfrak B\) be a reduced basic SAU boundary trace whose boundary expression has a primitive head constructor in \(\mathcal G_{\Basic}\).  Then \(\mathfrak B\) has a unique primitive profile label in \(\mathcal R_{\Basic}\), namely
\[
\Prof(\mathfrak B)=\rho_{\Basic}(\Head(C)).
\]
More explicitly:
\begin{enumerate}
\item if \(\Head(C)\) is a composite or power constructor, then \(\Prof(\mathfrak B)=\RootComp\);
\item if \(\Head(C)\) is a coefficient or standard-part constructor, then \(\Prof(\mathfrak B)=\Coeff\);
\item if \(\Head(C)\) is a projection or shadow constructor, then \(\Prof(\mathfrak B)=\ProjShad\);
\item if \(\Head(C)\) is a fixed-point or invariant constructor, then \(\Prof(\mathfrak B)=\InvDesc\);
\item if \(\Head(C)\) is a trace, norm, or symmetric aggregation constructor, then \(\Prof(\mathfrak B)=\AggInv\), with an induced invariant-descent audit when the aggregation is fixed by the declared symmetry;
\item if \(\Head(C)\) is a quotient or obstruction-class constructor, then \(\Prof(\mathfrak B)=\ObsLed\).
\end{enumerate}
\end{theorem}

\begin{proof}
The registry map \(\rho_{\Basic}\) is defined on primitive heads in
\Cref{def:sau-basic-profile-registry}.  The head-extraction definition fixes the
outer audit-relevant constructor \(\Head(C)\), so the value
\(\rho_{\Basic}(\Head(C))\) is unique.  The six displayed clauses are the six
cases in the basic audit grammar.  This theorem is tracked by the formalization
harness as a projection over the registry assignment.
\end{proof}

The six profile labels and their canonical witnesses are summarized in
\Cref{tab:profile-registry}.

\begingroup\begin{table}[tbp]
\centering
\footnotesize
\caption{Basic SAU profile registry and canonical witnesses.}
\label{tab:profile-registry}
\begin{tabularx}{\textwidth}{@{}>{\raggedright\arraybackslash}p{0.15\textwidth}X>{\raggedright\arraybackslash}p{0.31\textwidth}@{}}
\toprule
Profile & Boundary head & Canonical witness \\
\midrule
\(\RootComp\) & Composite or power constructor: a selected iterate or finite composite of a hidden update descends. & \(\Imag\): the square of the quarter-turn descends while the quarter-turn does not. \\
\(\Coeff\) & Coefficient or standard-part constructor: an audit map extracts a visible coefficient, jet, residue, leading term, or standard part. & \(\Jet\): coefficient extraction \([\varepsilon]\) recovers the derivative. \\
\(\ProjShad\) & Projection or shadow constructor: a declared shadow, observable, quotient, or forgetful projection supplies the lower-visible expression. & Real-shadow projection, as a sub-form of \(\Imag\); no separate case study. \\
\(\InvDesc\) & Fixed-point or invariant constructor: the selected expression is fixed by a declared symmetry and therefore descends. & \(\Galois\): fixed-field and Galois-invariant elements. \\
\(\AggInv\) & Trace, norm, or symmetric aggregation constructor: a finite symmetric aggregate of a promoted orbit descends. & \(\Galois\) subprofile via trace and norm. \\
\(\ObsLed\) & Quotient or obstruction-class constructor: a class records descent failure modulo a declared exact, trivial, or lower-null subobject. & \(\Coh\): graph \(H^1\) obstruction ledger. \\
\bottomrule
\end{tabularx}
\end{table}
 \endgroup

The four main case studies occupy four of the six primitive slots, with trace and
norm appearing as an aggregate-invariant subprofile of the Galois case.

\begin{corollary}[Four canonical witnesses classify by head]
\label{cor:sau-four-cases-classify-by-head}
The four canonical witnesses have the following primitive head classifications:
\[
\Prof(\Imag)=\RootComp,
\qquad
\Prof(\Jet)=\Coeff,
\]
\[
\Prof(\Galois)=\InvDesc
\quad\text{with trace/norm subprofile }\AggInv,
\qquad
\Prof(\Coh)=\ObsLed.
\]
\end{corollary}

\begin{proof}
Apply \Cref{thm:sau-primitive-head-classification} to the reduced boundary of
each canonical case.  The imaginary case has a composite head, the calculus case
has a coefficient head, the Galois case has an invariant head with trace and norm
as symmetric aggregations, and the cohomology case has a quotient or
obstruction-class head.  This corollary is tracked by the formalization harness
as a projection over the registry assignment and the four case records.
\end{proof}

The labels are typed, so matching one profile does not automatically identify the
boundary with another profile.

\begin{proposition}[Pairwise sort distinction of the four canonical profiles]
\label{prop:sau-pairwise-profile-distinction}
The four canonical profile labels
\[
\RootComp,
\qquad \Coeff,
\qquad \InvDesc,
\qquad \ObsLed
\]
are pairwise distinct in the typed basic grammar.  In particular, a proof that a boundary is of one canonical type does not by itself prove that it is of another canonical type.
\end{proposition}

\begin{proof}
Each listed profile is the image under \(\rho_{\Basic}\) of a different
primitive constructor sort: composite or power, coefficient or standard part,
invariant or fixed point, and quotient or obstruction class.  Since these sorts
are disjoint in the typed grammar, their profile labels are pairwise distinct.
This proposition is tracked by the formalization harness as a projection over
the typed profile registry.
\end{proof}

The next subsection extends the classification from primitive heads to composite
boundaries and states the classification soundness theorem.

\subsection{Composite boundaries and classification soundness}
\label{subsec:classification-composite-soundness}

Primitive heads cover one-step boundary expressions.  Composite boundary
expressions, where \(C(U)\) uses several declared constructors, need a finer
rule: their profile records the finite ordered combination of primitive heads
encountered in the reduced expression.

\begin{theorem}[Composite boundary classification]
\label{thm:sau-composite-boundary-classification}
Let \(\mathfrak B\) be a reduced basic SAU boundary trace.  Then its profile is either a primitive label in \(\mathcal R_{\Basic}\) or a finite product profile generated by primitive labels:
\[
\Prof(\mathfrak B)\in \Gen_{\mathrm{fin}}(\mathcal R_{\Basic}).
\]
If \(C(U)\) is a composite expression, its profile is the ordered list of primitive head labels encountered at the first essential audited boundaries of its reduced subexpressions.
\end{theorem}

\begin{proof}
Proceed by induction on the constructor depth of \(C(U)\).  At depth one, the
claim is exactly the primitive head classification theorem.  For the inductive
step, decompose \(C(U)\) into its outer constructor and reduced subexpressions;
the registry assigns the outer primitive head, and the induction hypothesis
assigns finite profile lists to the subexpressions.  The closure clause in the
basic profile registry combines these into an element of
\(\Gen_{\mathrm{fin}}(\mathcal R_{\Basic})\).  This theorem is tracked by the
formalization harness as a projection over the registry assignment and the
reduced-boundary predicates.
\end{proof}

The resulting normal form is the classification version of the boundary normal
form: a hidden support, an audited descent, and a finite registered profile.

\begin{corollary}[Basic profile normal form]
\label{cor:sau-basic-profile-normal-form}
Every finite strict useful trace whose reduced boundary lies in the basic audit grammar has an audit-equivalent normal form
\[
\left[U\not\downarrow_{\mathcal T_0},
\quad
C(U)\downarrow_B a^\sharp,
\quad
\Prof(C)\in\Gen_{\mathrm{fin}}(\mathcal R_{\Basic})
\right].
\]
Consequently, within the scope of the basic grammar, strict mathematical usefulness is mediated by a finite combination of registered descent mechanisms.
\end{corollary}

\begin{proof}
Combine the reduced-boundary normal form from \Cref{sec:necessity} with
\Cref{thm:sau-composite-boundary-classification}.  The first supplies the
non-descending support and audited descending boundary expression; the second
classifies the boundary expression by a finite profile generated from
\(\mathcal R_{\Basic}\).  This corollary is tracked by the formalization harness
as a projection over the boundary normal form and profile registry.
\end{proof}

The classification is also sound: a classified boundary licenses only the lower
descendant selected by that boundary, not the hidden support itself.

\begin{theorem}[Basic classification soundness]
\label{thm:sau-basic-classification-soundness}
Let \(\mathfrak T\) be a finite strict useful trace whose reduced boundary is basic.  Then there exists a validated SAU certificate whose profile is classified by \(\mathcal R_{\Basic}\):
\[
\operatorname{StrictUsefulTrace}(\mathfrak T)
\quad\text{and}\quad
\BasicSAU(\Red(\mathfrak T))
\]
\[
\Longrightarrow
\exists \mathfrak C,\mathfrak p
\quad
\ValSAU(\mathfrak C)
\quad\text{and}\quad
\mathfrak p\in\Gen_{\mathrm{fin}}(\mathcal R_{\Basic}).
\]
Moreover,
\[
\Use(B,P):\texttt{accepted}
\]
is licensed only for the lower descendant selected by the classified boundary, not for the hidden support itself.
\end{theorem}

\begin{proof}
The strict useful trace supplies the extracted boundary package by the necessity
machinery of \Cref{sec:necessity}.  Since the reduced boundary is basic,
\Cref{thm:sau-composite-boundary-classification} assigns a finite registered
profile.  Validation then gives the SAU certificate, while the soundness and
no-overreading discipline of \Cref{sec:soundness-failure} restricts the licensed
claim to the selected lower descendant.  This theorem is tracked by the
formalization harness as a projection over the reduced-boundary, registry, and
validated-certificate predicates.
\end{proof}

This is the section's anti-tautology point: the classification is applied after
strict useful traces have already forced a reduced boundary.

\begin{corollary}[Why the classification is non-tautological]
\label{cor:sau-classification-nontautological}
The classification theorem does not say that a bridge is useful because it has been named useful.  It says that, for finite strict useful traces whose boundary lies in the basic grammar, the necessity theorem first extracts a hidden-support/descending-boundary normal form, and the head-classification theorem then places that boundary in a typed registry.
\end{corollary}

\begin{proof}
The order of dependence is the point.  The trace is assumed strict and useful
before a profile label is assigned.  The necessity and tautology-barrier results
of \Cref{sec:necessity} extract the hidden-support/descending-boundary normal
form from that trace, and only then does the present section classify the head of
the extracted boundary.  This corollary is tracked by the formalization harness
as a projection over the necessity and classification predicates.
\end{proof}

The four canonical witnesses are now classified.  The next four sections develop
those witnesses in detail.
 \section{Imaginary Numbers}\label{sec:imaginary}

\subsection{The lower problem: same-layer half-step of real reversal}
\label{subsec:imag-lower-problem}

The imaginary case begins with a deliberately small lower layer.  Its visible
states are \(Y=\{-1,0,1\}\), and the operation already available downstairs is
real reversal, \(F(y)=-y\).  The lower problem asks whether this same layer can
find a half-step of reversal: a map \(H:Y\to Y\) whose square is \(F\).

\begin{definition}[Finite real-shadow layer]
\label{def:imag-real-shadow-layer}
Let
\[
Y=\{-1,0,1\}
\]
and let
\[
F:Y\to Y,
\qquad
F(y)=-y.
\]
The map \(F\) is called the lower reversal.  The same-layer half-step problem asks for a map
\[
H:Y\to Y
\]
such that
\[
H^2=F.
\]
For this profile set
\[
\mathsf{Cand}_0(P_{\Imag})=\End(Y)
\]
and define the exact defect
\[
\delta_{\Imag,0}(H)
=
\begin{cases}
0,& H^2=F,\\
1,& H^2\neq F.
\end{cases}
\]
The pass threshold is zero defect.  Thus a lower candidate passes precisely when it is a same-layer square root of reversal.
\end{definition}

\begin{lemma}[No same-layer half-step]
\label{lem:imag-no-same-layer-root}
There is no map \(H:Y\to Y\) satisfying
\[
H^2=F.
\]
Equivalently,
\[
\forall H\in\End(Y),
\qquad
\delta_{\Imag,0}(H)=1.
\]
\end{lemma}

\begin{proof}
There are \(|Y|^{|Y|}=3^3=27\) candidates in \(\End(Y)\).  Suppose one of
them satisfies \(H^2=F\).  Since \(F\) is bijective, \(H^2\) is bijective; on
a finite set this forces \(H\) itself to be bijective.  Thus the original
27-case check reduces to the six permutations of \(Y\).  Now \(F\) is the
transposition exchanging \(1\) and \(-1\) while fixing \(0\), so \(F\) is odd.
But the square of any permutation is even, because
\[
\operatorname{sgn}(H^2)=\operatorname{sgn}(H)^2=1.
\]
Hence no permutation \(H\) can satisfy \(H^2=F\), and therefore no endomorphism
of \(Y\) can do so.  A full enumeration of the original 27 candidates is
deferred to \Cref{app:methods}.  The same finite check is verified in Lean as
a faithful computational proof; see also \Cref{app:formalization} for the Lean
traceability.
\end{proof}

\begin{corollary}[Lower-layer saturation for half-reversal]
\label{cor:imag-lower-saturation}
The lower problem profile \(\mathfrak P_{\Imag,0}\) is saturated:
\[
\operatorname{Sat}(\mathfrak P_{\Imag,0}).
\]
\end{corollary}

The corollary says that the lower layer is not missing reversal itself; it is
missing a same-layer square root of reversal.  This corollary is tracked by
the formalization harness as a projection over the no-same-layer lemma and
the saturation predicate, rather than as an independent Lean proof.

\subsection{The four-phase carrier and the quarter-turn}
\label{subsec:imag-four-phase-carrier}

The saturated lower problem is resolved by moving to a four-state promoted
carrier.  Instead of asking the three-state shadow \(Y\) to contain a half-step
of reversal, we use the cyclic carrier \(C_4=\{1,i,-1,-i\}\), where the
quarter-turn \(G\) exists as multiplication by \(i\).  The real-shadow lens
\(q:C_4\to Y\) sends \(1\) and \(-1\) to their lower images and collapses the
pair \(\{i,-i\}\) to \(0\).

\begin{definition}[Four-phase carrier and real-shadow lens]
\label{def:imag-four-phase-carrier}
Let
\[
X=C_4=\{1,i,-1,-i\}
\]
with the cyclic update
\[
G:X\to X,
\qquad
G(z)=iz.
\]
Thus
\[
1\mapsto i\mapsto -1\mapsto -i\mapsto 1.
\]
Define the real-shadow lens
\[
q:X\to Y
\]
by
\[
q(1)=1,
\qquad
q(-1)=-1,
\qquad
q(i)=q(-i)=0.
\]
The promoted operation is \(U=G\).  The selected expression former is
\[
C(G)=G^2.
\]
\end{definition}

The geometry is illustrated in \Cref{fig:c4-shadow}.
\begingroup\begin{figure}[tbp]
\centering
\begin{tikzpicture}[
  >=stealth,
  phase/.style={circle,draw,minimum size=7mm,inner sep=1pt},
  shadow/.style={circle,draw,minimum size=7mm,inner sep=1pt},
  every node/.style={font=\small}
]
  \node[phase] (one) at (0.8,1.8) {\(1\)};
  \node[phase] (i) at (2.4,3.0) {\(i\)};
  \node[phase] (negone) at (4.0,1.8) {\(-1\)};
  \node[phase] (negi) at (2.4,0.6) {\(-i\)};

  \node[above=4mm of i] {\(C_4=\{1,i,-1,-i\}\)};

  \draw[->,thick] (one) to[bend left=10] node[above left] {\(G\)} (i);
  \draw[->,thick] (i) to[bend left=10] node[above right] {\(G\)} (negone);
  \draw[->,thick] (negone) to[bend left=10] node[below right] {\(G\)} (negi);
  \draw[->,thick] (negi) to[bend left=10] node[below left] {\(G\)} (one);

  \node[shadow] (yone) at (0.8,-1.0) {\(1\)};
  \node[shadow] (yzero) at (2.4,-1.0) {\(0\)};
  \node[shadow] (yneg) at (4.0,-1.0) {\(-1\)};

  \node[below=4mm of yzero] {\(Y_3=\{-1,0,1\}\)};

  \draw[dashed,->] (one) -- node[left] {\(q\)} (yone);
  \draw[dashed,->] (i) -- node[right] {\(q\)} (yzero);
  \draw[dashed,->] (negone) -- node[right] {\(q\)} (yneg);
  \draw[dashed,->] (negi) -- node[right] {\(q\)} (yzero);

  \draw[->,thick] (yone) to[bend left=15] node[above] {\(F\)} (yneg);
  \draw[->,thick] (yneg) to[bend left=15] (yone);
  \draw[->,thick] (yzero) to[out=120,in=60,looseness=4] node[above] {\(F\)} (yzero);
\end{tikzpicture}
\caption{The four-phase carrier \(C_4\), the real-shadow lens \(q:C_4\to Y_3\), the quarter-turn \(G\), and the descended reversal \(F\) determined by \(G^2\).}
\label{fig:c4-shadow}
\end{figure}
 \endgroup

The imaginary-number promotion bridge \(B_{\Imag}\) packages this construction
as an instance of the bridge data from \Cref{sec:layers}.  Its lower package is
the finite real-shadow layer on \(Y\), its promoted package is the four-phase
carrier \(C_4\), its descent channel is \(q\), its promoted operation is
\(U=G\), and its selected expression is \(C(G)=G^2\), whose audited lower
output is the reversal data already visible in \(Y\).

\begin{definition}[The imaginary-number promotion bridge]
\label{def:imag-promotion-bridge}
The imaginary-number promotion bridge is the record
\[
B_{\Imag}=(\mathcal T_0,\mathcal T_1,I_{\Imag},S,q,
\pi_0,
\pi_1,U,C,
\mathcal G_{\Imag},
A_{\Imag},
\mathcal L_{\Imag},
\mathcal N_{\Imag}),
\]
where:
\begin{itemize}
\item \(\mathcal T_0\) is the finite real-shadow package on \(Y=\{-1,0,1\}\) with reversal \(F\);
\item \(\mathcal T_1\) is the four-phase package on \(X=C_4\), optionally completed real-linearly to \(\mathbb R[J]\cong\mathbb C\);
\item \(S=X\), \(q:X\to Y\) is the real-shadow lens, \(\pi_0=q\), and \(\pi_1=\id_X\);
\item \(U=G\) is the quarter-turn update;
\item \(C(G)=G^2\);
\item the accepted descended output is the audited certificate
\[
a^\sharp_{\Imag}:=(F,\ qG^2=Fq,\ G\not\downarrow_q);
\]
\item the nonclaim register includes
\[
G\notin \End(Y),
\qquad
G\not\downarrow_q,
\qquad
\text{``}\,i\,\text{ is not a real-shadow state.''''}
\]
\end{itemize}
\end{definition}

The bridge is now specified; the next step is to check that \(G\) itself does
not descend while its square \(G^2\) does.

\subsection{Non-descent of \texorpdfstring{$G$}{G} and descent of \texorpdfstring{$G^2$}{G2}}
\label{subsec:imag-nondescent-descent}

The first check is that the hidden quarter-turn is not itself visible as a
lower update through the real-shadow lens.

\begin{lemma}[The quarter-turn does not descend]
\label{lem:imag-G-nondescends}
The update \(G\) does not descend through \(q\):
\[
G\not\downarrow_q.
\]
Equivalently,
\[
\delta^{\mathrm{split}}_1(q,G)\neq\varnothing.
\]
\end{lemma}

\begin{proof}
The split-pair witness is \((i,-i)\).  These two promoted states have the same
real shadow,
\[
q(i)=q(-i)=0,
\]
but their images after applying \(G\) have different real shadows:
\[
qG(i)=q(-1)=-1,
\qquad
qG(-i)=q(1)=1.
\]
Therefore
\[
(i,-i)\in\delta^{\mathrm{split}}_1(q,G),
\]
so the split-pair descent defect is nonempty and \(G\not\downarrow_q\).  This
finite check is verified in Lean as a faithful computational proof; see
\Cref{app:formalization}.
\end{proof}

The square of the hidden quarter-turn is different: it has a well-defined
real-shadow action, namely ordinary reversal.

\begin{lemma}[The square descends as real reversal]
\label{lem:imag-G2-descends}
The square \(G^2\) descends through \(q\) to the lower reversal \(F\):
\[
qG^2=Fq.
\]
Therefore
\[
G^2\downarrow_q F
\]
and
\[
\delta^{\mathrm{split}}_1(q,G^2)=\varnothing.
\]
\end{lemma}

\begin{proof}
Since \(G(z)=iz\), one has \(G^2(z)=i(iz)=-z\).  Hence \(qG^2(z)=q(-z)\).
Checking the four elements of \(C_4\) gives
\[
q(-1)=-1=-q(1),
\qquad
q(1)=1=-q(-1),
\]
and
\[
q(-i)=0=-q(i),
\qquad
q(i)=0=-q(-i).
\]
Thus \(q(-z)=-q(z)=F(q(z))\) for every \(z\in X\), so \(qG^2=Fq\).  Therefore
\(G^2\downarrow_q F\), and the split-pair descent defect of \(G^2\) is empty.
This finite check is verified in Lean as a faithful computational proof; see
\Cref{app:formalization}.
\end{proof}

Together, the two preceding lemmas give the basic root-composite shape of the
imaginary example.

\begin{proposition}[Root-composite profile]
\label{prop:imag-root-composite-profile}
The four-phase promotion realizes the root-composite normal form
\[
G\not\downarrow_q,
\qquad
G^2\downarrow_q F.
\]
Thus the hidden half-step is not a lower-layer operation, while its composite is a lower-layer operation.
\end{proposition}

\begin{proof}
The non-descent clause is \Cref{lem:imag-G-nondescends}, and the descended
composite clause is \Cref{lem:imag-G2-descends}.  Lean records this profile
label; the substantive content is delivered by the two preceding lemmas.
\end{proof}

The same finite construction is also minimal in cardinality for this
strict phase-promotion role.

\begin{proposition}[Cardinal minimality of strict phase promotion]
\label{prop:imag-cardinal-minimality}
Let \(q:X\twoheadrightarrow Y\) be a surjective finite shadow map.  Suppose there is an update \(G:X\to X\) such that \(G\) does not descend through \(q\).  Then
\[
|X|\geq 4.
\]
The four-phase carrier \(C_4\) realizes this lower bound while also satisfying
\[
G^2\downarrow_q F.
\]
Consequently \(C_4\) is a cardinal-minimal strict finite carrier for a non-descending square root of real reversal.
\end{proposition}

\begin{proof}
In the declared real-shadow problem, \(Y=\{-1,0,1\}\), so any surjection
\(q:X\twoheadrightarrow Y\) has \(|X|\geq 3\).  If \(|X|=3\), then \(q\) is
bijective.  In that case every update \(G:X\to X\) descends, by defining
\[
G^\sharp=qGq^{-1}:Y\to Y
\]
and observing that \(qG=G^\sharp q\).  This contradicts \(G\not\downarrow_q\).
Thus any such non-descending strict phase promotion has \(|X|\geq 4\).  The
four-phase carrier \(C_4\) has four elements, \Cref{lem:imag-G-nondescends}
gives \(G\not\downarrow_q\), and \Cref{lem:imag-G2-descends} gives
\(G^2\downarrow_q F\).  The full small-cardinality enumeration is deferred to
\Cref{app:methods}.  Lean verifies the no-same-layer-root component, here
corresponding to the lower bound on \(|X|\); the universal \(|X|\geq 4\)
statement is delivered as a paper-side argument.
\end{proof}

Finally, the bridge is strict because the promoted interface distinguishes
states that the lower real-shadow interface identifies.

\begin{proposition}[Strictness witness]
\label{prop:imag-strictness-witness}
Let the comparison carrier be \(S=X\), let
\[
\pi_0=q:X\to Y,
\qquad
\pi_1=\id_X:X\to X.
\]
Then the promotion is strict:
\[
\delta_{\mathrm{fact}}(\pi_0,
\pi_1)\neq\varnothing.
\]
In particular,
\[
\delta_{\mathrm{fact}}(\pi_0,\pi_1)\neq\varnothing.
\]
\end{proposition}

\begin{proof}
The pair \((i,-i)\in S\times S\) has the same lower view:
\[
\pi_0(i)=q(i)=0=q(-i)=\pi_0(-i).
\]
It has different promoted views:
\[
\pi_1(i)=i\neq -i=\pi_1(-i).
\]
Therefore
\[
(i,-i)\in\delta_{\mathrm{fact}}(\pi_0,\pi_1),
\]
so the factorization defect is nonempty.  This finite check is verified in
Lean as a faithful computational proof; see \Cref{app:formalization}.
\end{proof}

\subsection{SAU certificate, real-linear completion, and no-overreading}
\label{subsec:imag-certificate-completion-no-overreading}

The preceding results supply the fields of the imaginary Strict Audited Utility
certificate.  Saturation is \Cref{cor:imag-lower-saturation}; strictness is
\Cref{prop:imag-strictness-witness}; non-descent is
\Cref{lem:imag-G-nondescends}; and the descended expression is
\Cref{lem:imag-G2-descends}.  The bridge package \(B_{\Imag}\) from
\Cref{def:imag-promotion-bridge} supplies the declared audit data, accepted
lower output, and nonclaim register.

\begin{theorem}[Imaginary numbers satisfy Strict Audited Utility]
\label{thm:imag-sau-certificate}
The package \(B_{\Imag}\) supplies a validated Strict Audited Utility certificate for the square-root-of-reversal defect:
\[
\mathsf{Imag}\models \mathsf{SAU}.
\]
Equivalently,
\[
\operatorname{Sat}(\mathfrak P_{\Imag,0})
+
\operatorname{Strict}(B_{\Imag})
+
G\not\downarrow_q
+
G^2\downarrow_q F
+
\operatorname{AuditPass}(B_{\Imag})
\]
licenses the accepted usefulness claim
\[
\operatorname{Useful}(B_{\Imag},P_{\Imag}):\texttt{accepted}.
\]
\end{theorem}

\begin{proof}
This is the certificate assembly.  The saturation component is
\Cref{cor:imag-lower-saturation}; the strictness component is
\Cref{prop:imag-strictness-witness}; the hidden non-descent component is
\Cref{lem:imag-G-nondescends}; and the descended expression component is
\Cref{lem:imag-G2-descends}.  The remaining fields are the bridge audit,
accepted lower output \(a^\sharp_{\Imag}\), and nonclaim register declared in
\Cref{def:imag-promotion-bridge}, read through the certificate format of
\Cref{sec:layers}.  Therefore \(B_{\Imag}\) supplies the advertised certificate
and licenses only the accepted lower usefulness claim.  The certificate
packaging is tracked by the formalization harness as a projection over these
constituent results.
\end{proof}

The finite carrier \(C_4\) is the discrete phase skeleton of the usual
imaginary extension.  To recover the familiar linear form, ask for the minimal
nonzero real vector package carrying a hidden real-linear half-step \(J\) with
\(J^2=-I\).

\begin{theorem}[Minimal real-linear completion]
\label{thm:imag-minimal-real-linear-completion}
Let \(J\) be a real-linear hidden half-step satisfying
\[
J^2=-I.
\]
The minimal nonzero real-linear package carrying such a \(J\) is two-dimensional.  In that package, the real algebra generated by \(J\) is
\[
\mathbb R[J]=\{aI+bJ:a,b\in\mathbb R\},
\]
with multiplication
\[
(aI+bJ)(cI+dJ)=(ac-bd)I+(ad+bc)J.
\]
Consequently
\[
\mathbb R[J]\cong \mathbb C
\]
under the map
\[
a+bi\longmapsto aI+bJ.
\]
\end{theorem}

\begin{proof}
Let \(V\) be a nonzero real vector space carrying such an operator \(J\).  If
\(\dim_{\R}V=1\), then \(J\) is multiplication by some \(\lambda\in\R\), and
\(J^2=-I\) would give \(\lambda^2=-1\), impossible over \(\R\).  Hence
\(\dim_{\R}V\geq 2\).  Choose \(0\neq e\in V\).  The vectors \(e\) and \(Je\)
are linearly independent: if \(Je=\mu e\), then applying \(J\) again gives
\(-e=J^2e=\mu^2 e\), again impossible over \(\R\).  Thus
\(\operatorname{span}_{\R}\{e,Je\}\) is a two-dimensional \(J\)-stable package.

In the ordered basis \((e,Je)\), the operator \(J\) is the standard quarter-turn
matrix, since \(J(e)=Je\) and \(J(Je)=-e\).  Therefore the real algebra it
generates consists exactly of the operators \(aI+bJ\), and multiplication is
computed from \(J^2=-I\):
\[
(aI+bJ)(cI+dJ)=(ac-bd)I+(ad+bc)J.
\]
The map \(a+bi\mapsto aI+bJ\) preserves addition, multiplication, and the unit,
so it is a ring isomorphism \(\C\cong\R[J]\).  The even-dimensional obstruction
and the linear audit-channel details are deferred to \Cref{app:methods}.  This
theorem is tracked by the formalization harness as a projection over the
paper-side linear-algebra package, not as an independent Lean proof.
\end{proof}

The certificate therefore accepts that \(G^2\) descends to real reversal and
that \(i^2=-1\) is realized as a real-linear identity in the completion.  It
also accepts audited descendants such as real parts, conjugation-fixed
elements, and norms.  It does not assert that the hidden quarter-turn itself has
become a lower real-shadow state or a real scalar.

\begin{corollary}[No literalization of \texorpdfstring{$i$}{i}]
\label{cor:imag-no-literalization}
The accepted usefulness claim for imaginary numbers does not entail
\[
i\in Y
\]
or, in the real-linear completion,
\[
i\in\mathbb R.
\]
It entails only that the promoted half-step is lawful upstairs and that its selected audited descendants, such as \(G^2\), real part, conjugation-fixed elements, and norm, descend.
\end{corollary}

\begin{proof}
By the certificate discipline of \Cref{sec:layers}, the accepted claim licenses
only the audited descended expression \(C(U)=G^2\) and the corresponding lower
acceptance \(a^\sharp_{\Imag}\).  It does not license membership claims about
the hidden promoted operation \(U=G\), and the nonclaim register in
\Cref{def:imag-promotion-bridge} explicitly records that \(G\notin\End(Y)\),
\(G\not\downarrow_q\), and that \(i\) is not a real-shadow state.  The same
discipline forbids reading the real-linear completion as a claim that
\(i\in\R\).  The corollary is tracked by the formalization harness as a
projection over the certificate.
\end{proof}

The next case, \Cref{sec:calculus}, repeats this form with infinitesimal
directions: the hidden first-order object does not descend, while its selected
coefficient expression does.
 \section{Infinitesimals and Calculus}\label{sec:calculus}

\subsection{Value-only differentiation and its defect}
\label{subsec:jet-value-only-defect}

The calculus case begins with a deliberately impoverished lower interface.  At
a fixed point \(x\), the lower layer sees only the value \(f(x)\).  A derivative
at \(x\) is more than that value: functions can agree at \(x\) while having
different first-order behavior there.

\begin{definition}[Value-only derivative-extraction problem]
\label{def:jet-value-only-problem}
The lower candidate class for the derivative-extraction problem is
\[
\mathsf{Cand}_0(P_{\Jet})=\{m:\mathbb R\to\mathbb R\},
\]
where a candidate \(m\) is intended to compute derivative information from the visible value \(f(x)\).  Define the exact defect
\[
\delta_{\Jet,0}(m)=
\begin{cases}
0,& m\circ\ev_x=D_x,\\
1,& m\circ\ev_x\neq D_x.
\end{cases}
\]
The pass threshold is zero defect.  Thus a lower candidate passes precisely when the derivative functional factors through the point-value interface.
\end{definition}

The obstruction is already visible in the simplest pair of functions: one can
hold the value fixed at \(x\) while changing the derivative at \(x\).

\begin{lemma}[Derivative does not factor through point value]
\label{lem:jet-derivative-not-value-factor}
There is no map
\[
m:\mathbb R\to\mathbb R
\]
such that
\[
m\circ\ev_x=D_x.
\]
Equivalently,
\[
\forall m\in\mathsf{Cand}_0(P_{\Jet}),
\qquad
\delta_{\Jet,0}(m)=1.
\]
\end{lemma}

\begin{proof}
Let
\[
f_1(t)=t-x,
\qquad
f_2(t)=0.
\]
Then \(f_1(x)=f_2(x)=0\), so \(\ev_x(f_1)=\ev_x(f_2)\).  But
\[
D_x(f_1)=1,
\qquad
D_x(f_2)=0.
\]
If a map \(m:\R\to\R\) satisfied \(m\circ\ev_x=D_x\), the equal point values
would force \(D_x(f_1)=D_x(f_2)\), a contradiction.  Therefore no value-only
map computes the derivative functional.  Lean proves an abstract obstruction
over a generic triple of polynomial-class, value-evaluation, and
derivative-evaluation maps; the \(\R[t]\)-specific instance is paper-side.
\end{proof}

Thus the value-only lower profile has no accepted candidate.

\begin{corollary}[Lower-layer saturation for value-only differentiation]
\label{cor:jet-lower-saturation}
The lower problem profile \(\mathfrak P_{\Jet,0}\) is saturated:
\[
\operatorname{Sat}(\mathfrak P_{\Jet,0}).
\]
\end{corollary}

\begin{proof}
The corollary is immediate from
\Cref{lem:jet-derivative-not-value-factor}: every lower candidate has defect
one, while the pass threshold in \Cref{def:jet-value-only-problem} is zero
defect.  This saturation claim is tracked by the formalization harness as a
projection over the lemma and the saturation predicate.
\end{proof}

The promoted calculus package will introduce a square-zero direction
\(\varepsilon\).  The next elementary fact records why such a direction is not
already a real scalar.

\begin{lemma}[No nonzero square-zero real scalar]
\label{lem:jet-no-real-nilpotent}
If \(r\in\mathbb R\) and
\[
r^2=0,
\]
then
\[
r=0.
\]
Consequently, the element \(\varepsilon\in\Dual\) is not a real scalar.
\end{lemma}

\begin{proof}
If \(r^2=0\) in \(\R\), then \(r\cdot r=0\).  Since \(\R\) is an integral
domain, this implies \(r=0\).  Equivalently, \(\R\) has no nonzero nilpotent
scalar.  Hence an element \(\varepsilon\) with \(\varepsilon^2=0\) but
\(\varepsilon\neq 0\), as in \(\Dual\), cannot be an element of \(\R\).  Lean
verifies the corresponding statement in a Bool-valued model; the \(\R\)-version
is established here.
\end{proof}

The lower layer cannot extract the derivative from value data alone.  The next
subsection introduces the dual-number algebra \(\Dual\), where the derivative
does descend through coefficient extraction.

\subsection{The dual-number algebra and the first jet}
\label{subsec:jet-dual-numbers-first-jet}

The promoted layer adjoins to \(\R\) a square-zero element \(\varepsilon\).  By
\Cref{lem:jet-no-real-nilpotent}, this element is not a real scalar; it carries
first-order information that the value-only interface cannot see.

\begin{definition}[Dual-number algebra]
\label{def:jet-dual-number-algebra}
Let
\[
\Dual=\mathbb R[\varepsilon]/(\varepsilon^2).
\]
Every element of \(\Dual\) has a unique form
\[
a+b\varepsilon,
\qquad
 a,b\in\mathbb R,
\]
with
\[
\varepsilon^2=0.
\]
Define the standard-part audit and coefficient audit by
\[
\st(a+b\varepsilon)=a,
\qquad
[\varepsilon](a+b\varepsilon)=b.
\]
The promoted first-order interface at \(x\) is
\[
j_x^1:\mathcal P\to\Dual,
\qquad
j_x^1(f)=f(x+\varepsilon).
\]
\end{definition}

The structure is illustrated in \Cref{fig:dual-jet}.
\begingroup\begin{figure}[tbp]
\centering
\begin{tikzpicture}[
  >=stealth,
  every node/.style={font=\small},
  box/.style={draw, rounded corners=2pt, align=center, inner sep=5pt},
  audit/.style={box, fill=black!3}
]
  \node[box] (poly) at (0,0) {\(f\)};
  \node[box] (jet) at (4.5,0) {\(j_x^1(f)=f(x)+f'(x)\varepsilon\)\\ \(\in\Dual\)};
  \node[audit] (value) at (2.7,-1.9) {\(f(x)\)};
  \node[audit] (deriv) at (6.3,-1.9) {\(f'(x)\)};

  \draw[->] (poly) -- node[above] {first jet \(j_x^1\)} (jet);
  \draw[->] (jet) -- node[left] {\(\st\)} (value);
  \draw[->] (jet) -- node[right] {\(\operatorname{Coeff}\)} (deriv);
  \draw[dashed,->] (poly) to[bend right=18] node[below left] {\(\ev_x\)} (value);
\end{tikzpicture}
\caption{The value-only interface \(\ev_x\) compared with the dual-number first jet and its standard-part and coefficient audit channels.}
\label{fig:dual-jet}
\end{figure}
 \endgroup

The two interfaces now sit side by side: the lower interface records only the
point value, while the promoted interface records the full first jet.

\begin{definition}[Value and first-jet interfaces]
\label{def:jet-interfaces}
Let the comparison carrier be the polynomial algebra
\[
S_{\Jet}=\mathbb R[t].
\]
Define
\[
\pi_0=\ev_x:S_{\Jet}\to\mathbb R
\]
and
\[
\pi_1=j_x^1:S_{\Jet}\to\Dual.
\]
The lower interface \(\pi_0\) sees only value.  The promoted interface \(\pi_1\) sees value plus first-order coefficient.
\end{definition}

For polynomials, the first jet is computed by substituting \(x+\varepsilon\)
and discarding all terms containing \(\varepsilon^2\).

\begin{lemma}[First-jet formula for polynomials]
\label{lem:jet-first-jet-formula}
For every polynomial \(f\in\mathbb R[t]\),
\[
f(x+\varepsilon)=f(x)+f'(x)\varepsilon
\]
in \(\Dual\).
\end{lemma}

\begin{proof}
By linearity, it suffices to check monomials \(f(t)=t^n\).  In \(\Dual\),
the binomial expansion gives
\[
(x+\varepsilon)^n
=x^n+n x^{n-1}\varepsilon
+\sum_{k\geq 2}\binom{n}{k}x^{n-k}\varepsilon^k.
\]
Every term in the sum vanishes because \(\varepsilon^2=0\).  Hence
\[
(x+\varepsilon)^n=x^n+n x^{n-1}\varepsilon
=f(x)+f'(x)\varepsilon.
\]
The general polynomial case follows by summing the monomial identities.  Lean
proves the algebraic core of this identity over a generic dual-number model;
the \(\R[t]\)-specific instance is established here.
\end{proof}

The standard-part audit recovers the part that the lower value-only interface
already knew.

\begin{corollary}[Standard part recovers value]
\label{cor:jet-standard-part-value}
For every polynomial \(f\in\mathbb R[t]\),
\[
\st(j_x^1(f))=f(x).
\]
Equivalently,
\[
\st\circ j_x^1=\ev_x.
\]
\end{corollary}

\begin{proof}
Apply \(\st\) to the identity in \Cref{lem:jet-first-jet-formula}.  Since
\(\st(a+b\varepsilon)=a\), the \(\varepsilon\)-coefficient is discarded and
\(\st(j_x^1(f))=f(x)\).  This is the same lower value as \(\ev_x(f)\).  Lean
proves the algebraic core of this identity over a generic dual-number model;
the \(\R[t]\)-specific instance is established here.
\end{proof}

The coefficient audit extracts the first-order part that the lower value-only
interface could not compute on its own.

\begin{corollary}[Coefficient audit recovers derivative]
\label{cor:jet-coefficient-recovers-derivative}
For every polynomial \(f\in\mathbb R[t]\),
\[
[\varepsilon]j_x^1(f)=f'(x).
\]
Equivalently,
\[
[\varepsilon]\circ j_x^1=D_x.
\]
\end{corollary}

\begin{proof}
Apply the coefficient audit to the identity in
\Cref{lem:jet-first-jet-formula}.  Since
\[
[\varepsilon](a+b\varepsilon)=b,
\]
the coefficient of \(\varepsilon\) in \(j_x^1(f)\) is \(f'(x)\).  Thus
\([\varepsilon]j_x^1(f)=D_x(f)\) for every polynomial \(f\).  Lean proves the
algebraic core of this identity over a generic dual-number model; the
\(\R[t]\)-specific instance is established here.
\end{proof}

The first jet, viewed through its coefficient audit, supplies the descended
expression that the value-only lower interface could not provide.  The next
subsection develops the strictness, non-descent, and product and chain rules of
this promotion.

\subsection{Coefficient descent and the calculus rules}
\label{subsec:jet-coefficient-descent-rules}

The value-only interface \(\pi_0=\ev_x\) identifies functions that the first-jet
interface \(\pi_1=j_x^1\) distinguishes.  That is the strictness of the
promotion.  Separately, the infinitesimal direction \(\varepsilon\) itself is
not a real scalar, so it is not a lower-layer object.

\begin{lemma}[First-jet promotion is strict]
\label{lem:jet-promotion-strict}
The promoted first-jet interface does not factor through the value-only interface:
\[
\delta_{\mathrm{fact}}(\pi_0,\pi_1)\neq\varnothing.
\]
\end{lemma}

\begin{proof}
Take \(f(t)=0\) and \(g(t)=t-x\).  Then
\[
\pi_0(f)=\ev_x(f)=0=\ev_x(g)=\pi_0(g),
\]
but
\[
\pi_1(f)=j_x^1(f)=0,
\qquad
\pi_1(g)=j_x^1(g)=\varepsilon.
\]
Since \(0\neq\varepsilon\) in \(\Dual\), the pair \((f,g)\) belongs to
\(\delta_{\mathrm{fact}}(\pi_0,\pi_1)\).  Therefore the factorization defect is
nonempty.  Lean verifies the algebraic core of strictness over a generic
dual-number model; the polynomial-specific instance is established here.
\end{proof}

\begin{lemma}[The infinitesimal direction is non-descending]
\label{lem:jet-epsilon-nondescending}
The promoted direction \(\varepsilon\) is not a lower-layer real scalar and is not recoverable from the standard-part shadow.  More precisely,
\[
\st(0)=\st(\varepsilon)=0,
\qquad
0\neq\varepsilon,
\]
so the promoted identity interface on \(\Dual\) is strict over the standard-part shadow:
\[
\delta_{\mathrm{fact}}(\st,\id_{\Dual})\neq\varnothing.
\]
\end{lemma}

\begin{proof}
In the dual-number algebra, \(0=0+0\varepsilon\) while
\(\varepsilon=0+1\varepsilon\), so uniqueness of the form
\(a+b\varepsilon\) gives \(0\neq\varepsilon\).  At the same time,
\[
\st(0)=0=\st(\varepsilon).
\]
Thus the pair \((0,\varepsilon)\) has the same standard-part shadow but
different promoted identity views, so
\[
(0,\varepsilon)\in\delta_{\mathrm{fact}}(\st,\id_{\Dual}).
\]
The auxiliary obstruction \Cref{lem:jet-no-real-nilpotent} records the
corresponding lower-layer point: no nonzero square-zero direction lives in
\(\R\).  Lean verifies the algebraic core of this non-descent witness over a
generic dual-number model; the polynomial-specific instance is established
here.
\end{proof}

The familiar calculus rules now appear as coefficient-audit identities in
\(\Dual\).

\begin{proposition}[Product rule from dual-number multiplication]
\label{prop:jet-product-rule}
For all \(f,g\in\mathbb R[t]\),
\[
[\varepsilon](fg)(x+\varepsilon)
=
f'(x)g(x)+f(x)g'(x).
\]
\end{proposition}

\begin{proof}
Using \Cref{lem:jet-first-jet-formula} for \(f\) and \(g\),
\[
f(x+\varepsilon)g(x+\varepsilon)
=
\bigl(f(x)+f'(x)\varepsilon\bigr)
\bigl(g(x)+g'(x)\varepsilon\bigr).
\]
Multiplying in \(\Dual\) gives
\[
f(x)g(x)+\bigl(f'(x)g(x)+f(x)g'(x)\bigr)\varepsilon
+f'(x)g'(x)\varepsilon^2.
\]
The last term vanishes because \(\varepsilon^2=0\).  Applying
\([\varepsilon]\) extracts the coefficient
\[
f'(x)g(x)+f(x)g'(x).
\]
Lean proves the algebraic core of the product rule for any dual-number algebra;
the polynomial-specific statement specializes here.
\end{proof}

\begin{proposition}[Chain rule from first-jet substitution]
\label{prop:jet-chain-rule}
For all \(f,g\in\mathbb R[t]\),
\[
[\varepsilon]f(g(x+\varepsilon))
=
f'(g(x))g'(x).
\]
Equivalently,
\[
D_x(f\circ g)=D_{g(x)}(f)\,D_x(g).
\]
\end{proposition}

\begin{proof}
By \Cref{lem:jet-first-jet-formula},
\[
g(x+\varepsilon)=g(x)+g'(x)\varepsilon.
\]
Applying the first-jet formula for \(f\) at the base point \(g(x)\) and with
the square-zero increment \(g'(x)\varepsilon\) gives
\[
f(g(x+\varepsilon))
=
f\bigl(g(x)+g'(x)\varepsilon\bigr)
=
f(g(x))+f'(g(x))g'(x)\varepsilon.
\]
Applying the coefficient audit extracts \(f'(g(x))g'(x)\), which is the stated
chain rule.  Lean proves the algebraic core of this substitution identity over
a generic dual-number model; the polynomial-specific statement is established
here.
\end{proof}

The dual-number package is therefore an \(\R\)-algebra carrying a square-zero
generator that supports the coefficient-audit identities.  The next subsection
assembles the Strict Audited Utility certificate from minimality, the smooth
first-jet quotient, and the audited descended expression.

\subsection{Minimal square-zero promotion and SAU certificate}
\label{subsec:jet-minimal-package-certificate}

The dual-number algebra is not an arbitrary gadget added after the fact.  Once
one asks for a commutative real algebra with a nonzero square-zero direction,
the algebra generated by that direction has the dual-number form.

\begin{proposition}[Minimal square-zero real-algebra package]
\label{prop:jet-minimal-square-zero-package}
Let \(R_{\mathrm{alg}}\) be a commutative unital real algebra and let \(\eta\in R_{\mathrm{alg}}\) satisfy
\[
\eta^2=0,
\qquad
\eta\neq 0.
\]
Then the real subalgebra generated by \(1\) and \(\eta\) is canonically isomorphic to \(\Dual\):
\[
\mathbb R[\eta]\cong \mathbb R[\varepsilon]/(\varepsilon^2).
\]
\end{proposition}

\begin{proof}
Every element of the subalgebra generated by \(1\) and \(\eta\) has the form
\[
a+b\eta,
\qquad
a,b\in\R,
\]
because \(\eta^2=0\) removes all higher powers.  Multiplication is
\[
(a+b\eta)(c+d\eta)
=ac+(ad+bc)\eta+bd\eta^2
=ac+(ad+bc)\eta.
\]
Thus the assignment
\[
a+b\varepsilon\longmapsto a+b\eta
\]
is a surjective homomorphism \(\Dual\to\R[\eta]\).  If \(a+b\eta=0\), then
multiplying by \(\eta\) gives \(a\eta=0\).  If \(a\neq 0\), the scalar \(a\) is
invertible over \(\R\), forcing \(\eta=0\), a contradiction.  Hence \(a=0\).
Then \(b\eta=0\), and the same scalar-invertibility argument gives \(b=0\).
The homomorphism is injective, hence an isomorphism.  Lean verifies the
algebraic core in a Bool dual-number model; the universal-property statement
is established here as a paper-side argument.
\end{proof}

For ordinary smooth functions, the analogous object is the first-jet quotient
at \(x\): the local algebra \(C^\infty_x\) of smooth germs modulo
\(\mathfrak m_x^2\), where \(\mathfrak m_x\) is the ideal of germs vanishing at
\(x\).  This quotient identifies with \(\Dual\) by sending a germ \(f\) to
\(f(x)+f'(x)\varepsilon\).  The formal appendix-destination statement and its
smooth-germ proof are deferred to \Cref{app:methods}; Lean records the
quotient-algebra shape, while the smooth analytical content is paper-side.

All certificate ingredients have now been established.  The lower profile is
saturated by \Cref{cor:jet-lower-saturation}; strictness and non-descent are
\Cref{lem:jet-promotion-strict,lem:jet-epsilon-nondescending}; coefficient
audit recovers the derivative by \Cref{cor:jet-coefficient-recovers-derivative};
and the product and chain rules are \Cref{prop:jet-product-rule,prop:jet-chain-rule}.

\begin{theorem}[Calculus first-jets satisfy Strict Audited Utility]
\label{thm:jet-sau-certificate}
The first-order dual-number package supplies a Strict Audited Utility certificate for the value-only derivative-extraction problem:
\[
\mathsf{Jet}\models\mathsf{SAU}.
\]
More explicitly, with
\[
\pi_0=\ev_x,
\qquad
\pi_1=j_x^1,
\qquad
U=\varepsilon,
\qquad
C(U;f)=[\varepsilon]f(x+\varepsilon),
\qquad
 a^\sharp=D_x,
\]
the following hold:
\begin{enumerate}
\item the lower profile is saturated:
\[
\operatorname{Sat}(\mathfrak P_{\Jet,0});
\]
\item the promotion is strict:
\[
\delta_{\mathrm{fact}}(\pi_0,\pi_1)\neq\varnothing;
\]
\item the promoted direction is non-descending:
\[
\varepsilon\notin\mathbb R,
\qquad
\delta_{\mathrm{fact}}(\st,\id_{\Dual})\neq\varnothing;
\]
\item the selected expression descends by coefficient audit:
\[
[\varepsilon]\circ j_x^1=D_x;
\]
\item the descended answer solves the lower problem:
\[
D_x\in\mathsf{Ans}_0(P_{\Jet});
\]
\item the nonclaim register blocks the overread that \(\varepsilon\) is a real scalar or that the derivative was obtained from value alone.
\end{enumerate}
\end{theorem}

\begin{proof}
This is the certificate assembly.  The saturation field is
\Cref{cor:jet-lower-saturation}.  The strictness field is
\Cref{lem:jet-promotion-strict}.  The non-descending promoted direction is
\Cref{lem:jet-epsilon-nondescending}.  The audited descended expression is the
coefficient-recovery identity in \Cref{cor:jet-coefficient-recovers-derivative},
with the derivative \(D_x\) as the accepted lower answer.  The product and chain
rules, \Cref{prop:jet-product-rule,prop:jet-chain-rule}, record that the same
coefficient audit is stable under the basic calculus operations.  The nonclaim
field is enforced by \Cref{lem:jet-no-real-nilpotent} and
\Cref{lem:jet-derivative-not-value-factor}: \(\varepsilon\) is not a real scalar,
and the derivative was not obtained from value alone.  Read through the
certificate format of \Cref{sec:layers}, these fields give the displayed
Strict Audited Utility certificate.  The certificate packaging is tracked by
the formalization harness as a projection over the constituent results.
\end{proof}

The same shape, non-descent paired with audited descent, appears next with a
different audit channel: invariant descent through symmetry in \Cref{sec:galois}.
 \section{Galois Theory}\label{sec:galois}

\subsection{Roots not in the base field}
\label{subsec:gal-roots-not-in-base}

In algebra one routinely uses roots that do not live in the base field.  Let
\(K\) be a field and let \(p(t)\in K[t]\) be irreducible of degree at least
two.  The lower question asks whether \(K\) itself contains a root of \(p\); the
standard algebraic answer is no.  The useful replacement is not to declare the
root to be in \(K\), but to return invariant data, such as the orbit polynomial,
trace, and norm, that is generated from a promoted root package.

\begin{lemma}[Irreducible non-linear polynomials have no base-field root]
\label{lem:gal-no-base-root}
For the above \(p\), no element \(a\in K\) satisfies
\[
p(a)=0.
\]
Consequently,
\[
\forall a\in\mathsf{Cand}_0(P_{\Gal}),
\qquad
\delta_{\Gal,0}(a)=1.
\]
\end{lemma}

\begin{proof}
If \(a\in K\) satisfied \(p(a)=0\), then \(t-a\) would divide \(p(t)\) in
\(K[t]\).  Since \(p\) is irreducible, this forces \(p\) to be associated to the
linear polynomial \(t-a\), contradicting \(\deg p\ge2\).  Thus no base-field
candidate is a root, and the displayed defect condition follows from the
definition of the lower candidate class.  This lemma is tracked by the
formalization harness as a projection over the irreducibility hypothesis and the
root-use defect predicate.
\end{proof}

The lower problem is therefore saturated: the base field cannot supply the
individual root requested at the lower interface.

\begin{corollary}[Lower-layer saturation for individual root representation]
\label{cor:gal-root-saturation}
The lower problem profile for same-layer individual root representation is saturated:
\[
\operatorname{Sat}(\mathfrak P_{\Gal,0}).
\]
\end{corollary}

\begin{proof}
By \Cref{lem:gal-no-base-root}, every lower candidate \(a\in K\) fails the root
defect test for \(p\).  Hence no same-layer individual root representation passes
the lower problem profile.  This corollary is tracked by the formalization
harness as a projection over the no-base-root lemma and the saturation
predicate.
\end{proof}

The negative lower answer pushes the proof toward an audited replacement answer:
invariant data extracted from a promoted root, rather than membership of the root
itself in \(K\).

\begin{definition}[Galois root-use problem]
\label{def:gal-root-use-problem}
The Galois root-use problem \(P_{\Gal}\) has two layers:
\begin{enumerate}
\item the saturated lower request for an individual root in \(K\), with candidate class \(K\) and defect \(\delta_{\Gal,0}\);
\item the audited replacement answer type consisting of base-field invariant data generated from a promoted root package, such as
\[
C_{\mathrm{orb}}(\alpha),
\quad
\Tr_{L/K}(\alpha),
\quad
\Nm_{L/K}(\alpha),
\quad
\text{and other }G\text{-fixed expressions}.
\]
\end{enumerate}
The nonclaim register contains, at minimum,
\[
\alpha\in K
\quad\text{is not claimed.}
\]
\end{definition}

Thus the lower request is for an actual root in \(K\), while the accepted
replacement answer type consists of invariant expressions generated upstairs and
returned downstairs.  The next subsection introduces the splitting field as the
universal promoted layer that carries the roots and provides the audited route
back to \(K\).

\subsection{Root adjoining and the splitting field}
\label{subsec:gal-root-adjoining}

Given \(K\) and an irreducible \(p(t)\in K[t]\) of degree at least two, the
standard way to adjoin a root is the quotient \(K[t]/(p)\).  Because \(p\) is
irreducible, this quotient is a field extension of \(K\); iterating the same
construction gives a splitting field \(L\) containing all roots of \(p\).

\begin{theorem}[Universal minimal root-adjoining package]
\label{thm:gal-universal-root-adjoining}
Let \(R_{\mathrm{alg}}\) be a commutative \(K\)-algebra and let \(r\in R_{\mathrm{alg}}\) satisfy
\[
p(r)=0.
\]
Then there is a unique \(K\)-algebra homomorphism
\[
\varphi_r:K[t]/(p)\to R_{\mathrm{alg}}
\]
such that
\[
\varphi_r(t\bmod p)=r.
\]
In particular, \(K[t]/(p)\) is the minimal universal \(K\)-algebraic promotion that represents a root of \(p\).
\end{theorem}

\begin{proof}
Since \(p\) is irreducible, the ideal \((p)\) is maximal in \(K[t]\), so
\(K[t]/(p)\) is a field and \(\alpha=t\bmod p\) satisfies \(p(\alpha)=0\).
For any \(K\)-algebra \(R_{\mathrm{alg}}\) and root \(r\) of \(p\), the universal
property of \(K[t]\) gives the unique \(K\)-algebra map \(K[t]\to R_{\mathrm{alg}}\)
sending \(t\) to \(r\); because \(p(r)=0\), this map factors uniquely through
\(K[t]/(p)\).  This is the standard root-adjoining argument from algebra.  The
theorem is tracked by the formalization harness as a projection over the
root-adjoining record and its universal-property field.
\end{proof}

The tower and the invariant audit channels used later in the section are shown
in \Cref{fig:splitting-tower}.

\begingroup\begin{figure}[tbp]
\centering
\begin{tikzpicture}[
  >=stealth,
  node distance=1.35cm and 2.1cm,
  field/.style={draw, rounded corners=2pt, inner sep=4pt, minimum width=1.8cm},
  audit/.style={draw, rounded corners=2pt, inner sep=4pt, align=center, text width=3.0cm},
  note/.style={align=center}
]
  \node[field] (K) {\(K\)};
  \node[field, above=of K] (Ka) {\(K(\alpha)\)};
  \node[field, above=of Ka] (L) {\(L\)};
  \node[note, left=1.6cm of Ka] (nondesc) {\(\alpha\notin K\)\\[-1pt]\scriptsize non-descent};
  \node[note, right=2.0cm of L] (gal) {\(\Gal(L/K)\)\\[-1pt]\scriptsize acts on roots};
  \node[audit, right=2.0cm of Ka] (audits) {\(\Tr_{L/K}\), \(\Nm_{L/K}\),\\ \(L^{\Gal(L/K)}\)};
  \node[field, right=2.0cm of K] (Kout) {\(K\)-valued audits};

  \draw[->] (K) -- node[left] {\(\hookrightarrow\)} (Ka);
  \draw[->] (Ka) -- node[left] {\(\hookrightarrow\)} (L);
  \draw[->] (gal) -- (L);
  \draw[->] (L) -- node[above, sloped] {\scriptsize invariants} (audits);
  \draw[->] (audits) -- node[right] {\scriptsize descend} (Kout);
  \draw[dashed,->] (Ka) -- (nondesc);
  \draw[dashed,->] (nondesc) -- (K);
\end{tikzpicture}
\caption{A promoted root in a splitting field descends to the base field only through invariant audit channels.}
\label{fig:splitting-tower}
\end{figure}
 \endgroup

The splitting field is the promoted layer in which the root-representation defect
has been closed upstairs: all roots are present, and the Galois group acts on the
resulting package.

\begin{theorem}[Promoted splitting closes the root-representation defect upstairs]
\label{thm:gal-splitting-closes-upstairs}
In the splitting field \(L\), the polynomial \(p\) factors as
\[
p(t)=\prod_{j=1}^{n}(t-\alpha_j)
\]
with distinct roots \(\alpha_1,\ldots,\alpha_n\in L\).  Hence the promoted layer represents all roots of \(p\).
\end{theorem}

\begin{proof}
By construction, a splitting field \(L\) is generated over \(K\) by adjoining all
roots of \(p\).  Hence \(p\) factors in \(L[t]\) as the product of its linear
factors, and the promoted layer represents the full root package.  When the
extension is viewed with its Galois group, the group action permutes these roots.
These are standard splitting-field facts from Galois theory.  The theorem is
tracked by the formalization harness as a projection over the splitting-field
record.
\end{proof}

The root package is useful precisely because the individual roots remain
upstairs; the lower layer has not acquired them as elements of \(K\).

\begin{lemma}[A promoted root need not descend]
\label{lem:gal-root-not-descend}
For any root \(\alpha\in L\) of \(p\), one has
\[
\alpha\notin K.
\]
Thus an individual promoted root is not accepted as lower-layer content.
\end{lemma}

\begin{proof}
For the irreducible non-linear polynomial under discussion, \Cref{lem:gal-no-base-root}
shows that no element of \(K\) is a root of \(p\).  Thus an individual root
\(\alpha\) represented in the promoted splitting field is not a lower-layer
element.  This lemma is tracked by the formalization harness as a projection over
the no-base-root result and the non-descent field of the root package.
\end{proof}

Nor can the promoted root be pushed back to the base field by a \(K\)-algebra
retract.

\begin{lemma}[No base-field evaluation retract for a non-base root]
\label{lem:gal-no-root-retract}
There is no \(K\)-algebra homomorphism
\[
r:K[\alpha]\to K
\]
such that
\[
r(\alpha)=\alpha
\]
or, more generally, such that \(r(\alpha)\) is a root of \(p\) in \(K\).
\end{lemma}

\begin{proof}
If such a \(K\)-algebra homomorphism existed, then \(r(\alpha)\in K\) would
satisfy \(p(r(\alpha))=r(p(\alpha))=0\).  This would give a base-field root of
\(p\), contradicting \Cref{lem:gal-no-base-root}.  Thus there is no base-field
evaluation retract for the promoted root.  This lemma is tracked by the
formalization harness as a projection over the root-use non-descent and
no-retract predicates.
\end{proof}

Roots live upstairs; descent to \(K\) goes through invariant audit channels.  The
next subsection develops fixed-field, trace, norm, and orbit-polynomial descent.

\subsection{Symmetric invariants, traces, norms, and fixed-field descent}
\label{subsec:gal-invariants-trace-norm-fixed-field}

The splitting field is strict over \(K\): it distinguishes promoted roots that
the lower polynomial interface treats through the same base-field data.  Descent
therefore cannot proceed by returning an individual root to \(K\).  It proceeds
through invariant expressions: symmetric functions of the roots, fixed-field
tests, traces, norms, and orbit polynomials.

\begin{theorem}[Strictness of the splitting-field promotion]
\label{thm:gal-strictness}
Let
\[
S=\{\alpha_1,\ldots,\alpha_n\}\subseteq L
\]
be the set of roots of \(p\).  Define
\[
\pi_0:S\to K[t],
\qquad
\pi_0(\beta)=p,
\]
and
\[
\pi_1:S\to L,
\qquad
\pi_1(\beta)=\beta.
\]
Then
\[
\delta_{\mathrm{fact}}(\pi_0,\pi_1)\neq\varnothing.
\]
Hence the promoted root layer is strict relative to the lower polynomial interface.
\end{theorem}

\begin{proof}
Choose two distinct roots \(\alpha_i,\alpha_j\in S\).  Their lower view under
\(\pi_0\) is the same polynomial \(p\), while their promoted view under
\(\pi_1\) records distinct elements of \(L\).  Thus
\((\alpha_i,\alpha_j)\) belongs to the factorization defect
\(\delta_{\mathrm{fact}}(\pi_0,\pi_1)\), so the defect is nonempty.  The theorem
is tracked by the formalization harness as a projection over the splitting-field
strictness predicate.
\end{proof}

The first descent channel is the classical one: symmetric functions of all roots
are exactly coefficient data of the base-field polynomial.

\begin{theorem}[Elementary symmetric descent]
\label{thm:gal-symmetric-descent}
Let
\[
p(t)=t^n+a_{n-1}t^{n-1}+\cdots+a_1t+a_0
\]
with \(a_j\in K\), and let \(\alpha_1,\ldots,\alpha_n\) be its roots in \(L\).  Then the elementary symmetric functions
\[
e_j(\alpha_1,\ldots,\alpha_n)
\]
lie in \(K\), with
\[
e_j(\alpha_1,\ldots,\alpha_n)=(-1)^j a_{n-j}.
\]
More generally, every symmetric polynomial in the roots with coefficients in \(K\) descends to \(K\).
\end{theorem}

\begin{proof}
Expanding \(p(t)=\prod_{i=1}^n(t-\alpha_i)\) expresses each coefficient as a
signed elementary symmetric polynomial in the roots.  Since the coefficients
\(a_j\) lie in \(K\), the displayed identities place every elementary symmetric
function in \(K\).  The fundamental theorem of symmetric functions then writes
any symmetric polynomial with coefficients in \(K\) as a polynomial in the
elementary symmetric functions, hence as an element of \(K\).  The theorem is
tracked by the formalization harness as a projection over the symmetric-descent
audit field.
\end{proof}

The fixed-field theorem identifies the base field inside the promoted field as
the invariant part of the Galois action.

\begin{theorem}[Fixed-field audit]
\label{thm:gal-fixed-field-audit}
For the finite Galois extension \(L/K\),
\[
L^G=K,
\qquad
G=\Gal(L/K).
\]
Consequently, an element \(x\in L\) descends to the base field exactly when it is fixed by every \(K\)-automorphism of \(L\):
\[
x\in K
\quad\Longleftrightarrow\quad
\sigma(x)=x\text{ for all }\sigma\in G.
\]
\end{theorem}

\begin{proof}
This is the standard fixed-field theorem for a finite Galois extension: the
elements fixed by every \(K\)-automorphism of \(L\) are precisely the elements of
the base field \(K\).  Thus invariance under the full Galois group is exactly the
audit test for descent to \(K\).  The theorem is tracked by the formalization
harness as a projection over the fixed-field audit record.
\end{proof}

Trace and norm are the corresponding additive and multiplicative invariant
audit channels.

\begin{theorem}[Trace and norm are descending audits]
\label{thm:gal-trace-norm-descend}
For every \(x\in L\), define
\[
\Tr_{L/K}(x)=\sum_{\sigma\in G}\sigma(x),
\qquad
\Nm_{L/K}(x)=\prod_{\sigma\in G}\sigma(x).
\]
Then
\[
\Tr_{L/K}(x)\in K,
\qquad
\Nm_{L/K}(x)\in K.
\]
In particular, if \(\alpha\) is a promoted root of \(p\), then
\[
\Tr_{L/K}(\alpha),\ \Nm_{L/K}(\alpha)\in K.
\]
\end{theorem}

\begin{proof}
For \(\tau\in G\), the action of \(\tau\) permutes the automorphisms
\(\sigma\in G\).  Hence \(\tau\) fixes the sum \(\sum_{\sigma\in G}\sigma(x)\)
and the product \(\prod_{\sigma\in G}\sigma(x)\).  By the fixed-field audit
above, both invariant expressions lie in \(K\).  Applying this to \(x=\alpha\)
gives the promoted-root case.  The theorem is tracked by the formalization
harness as a projection over the trace-norm audit fields.
\end{proof}

The same invariant-descent mechanism packages the Galois orbit of a promoted
root as a lower polynomial.

\begin{corollary}[Orbit polynomials descend]
\label{cor:gal-orbit-polynomials-descend}
For \(\alpha\in L\), the orbit polynomial
\[
C_{\mathrm{orb}}(\alpha)(t)
=
\prod_{\beta\in \Orb_G(\alpha)}(t-\beta)
\]
has coefficients in \(K\).  If \(\alpha\) is a root of the irreducible separable polynomial \(p\), then
\[
C_{\mathrm{orb}}(\alpha)=p.
\]
\end{corollary}

\begin{proof}
The coefficients of \(C_{\mathrm{orb}}(\alpha)\) are elementary symmetric
polynomials in the Galois orbit of \(\alpha\).  They are fixed by the action of
\(G\), and therefore lie in \(K\) by the fixed-field audit.  If \(\alpha\) is a
root of the irreducible separable polynomial \(p\), its Galois orbit is the full
set of roots of \(p\), so the orbit polynomial is \(p\).  The corollary is
tracked by the formalization harness as a projection over the orbit-polynomial
descent field.
\end{proof}

With strictness, invariant descent, and the fixed-field theorem in place, the
next subsection assembles the Galois SAU certificate and its no-overreading
discipline.

\subsection{SAU certificate and no-literalization}
\label{subsec:gal-certificate-no-literalization}

The Galois case now has the certificate ingredients in place.  The lower root
request is saturated, the splitting-field promotion is strict, promoted roots do
not descend to \(K\), and invariant expressions such as orbit polynomials,
trace, norm, and fixed-field elements descend back to \(K\).  These are the data
assembled in the case certificate.

\begin{theorem}[Galois theory supplies an SAU certificate]
\label{thm:gal-sau-certificate}
For a monic separable irreducible \(p\in K[t]\) of degree greater than one, with splitting field \(L/K\) and Galois group \(G=\Gal(L/K)\), the Galois root package gives a Strict Audited Utility certificate:
\[
\mathsf{Gal}\models\mathsf{SAU}.
\]
More explicitly:
\begin{enumerate}
\item \textbf{Saturation:} no individual root of \(p\) lies in \(K\);
\item \textbf{Strictness:} distinct promoted roots have the same lower polynomial interface but different promoted identities;
\item \textbf{Non-descent:} a promoted root \(\alpha\) is not accepted as lower-layer content;
\item \textbf{Audited descent:} Galois-invariant expressions in the roots descend to \(K\);
\item \textbf{Usefulness:} the descended invariant data gives a lower-layer admissible certificate of the promoted root package;
\item \textbf{No-overreading:} the certificate does not claim that \(\alpha\in K\).
\end{enumerate}
\end{theorem}

\begin{proof}
Saturation is supplied by \Cref{cor:gal-root-saturation}.  Strictness is
\Cref{thm:gal-strictness}, and non-descent of an individual promoted root is
\Cref{lem:gal-root-not-descend}.  The audited descent fields are supplied by the
invariant channels established in \Cref{thm:gal-symmetric-descent},
\Cref{thm:gal-fixed-field-audit}, \Cref{thm:gal-trace-norm-descend}, and
\Cref{cor:gal-orbit-polynomials-descend}.  The nonclaim register is the
explicit refusal to infer \(\alpha\in K\) from those descended invariant data.
Thus the six certificate fields are present for the Galois root package.  The
theorem is tracked by the formalization harness as a projection over the
Galois-case certificate record and its constituent witness fields.
\end{proof}

The certificate licenses lower-layer claims about descended invariant data:
coefficients of the orbit polynomial, traces, norms, and fixed-field witnesses.
It does not license a membership claim for the individual root itself.

This is the Galois instance of the general no-overreading discipline from
\Cref{subsec:soundness-no-overreading}.  In the notation of the certificate,
\(\alpha \ndown_q\) while an invariant expression \(C(\alpha)\), such as
\(C_{\mathrm{orb}}(\alpha)\), \(\Tr_{L/K}(\alpha)\), or \(\Nm_{L/K}(\alpha)\),
has a descended lower answer \(C(\alpha)\down_B a^\sharp\).  The accepted claim
is the descended invariant answer; the lower-layer literalization
\(\alpha\in K\) is not licensed.

The next case, \Cref{sec:cohomology}, replaces invariant descent by an
obstruction ledger: usefulness is expressed by a class that records exactly when
descent succeeds or fails.
 \section{Cohomology}\label{sec:cohomology}

\subsection{Edge data, periods, and the local-to-global question}
\label{subsec:coh-edge-data-periods}

On a finite graph \(\Gamma\), edge data may or may not come from a global
vertex potential.  If \(a\in C^1(\Gamma;k)\) is the differential of a potential
\(\phi:V\to k\), then summing \(a\) around a closed cycle must telescope to
zero.  Periods therefore give a concrete audit of the local-to-global question:
does the edge datum glue to a global potential, or does a cycle record the
failure?

\begin{definition}[Period audit over cycles]
\label{def:coh-period-audit}
A directed closed walk or cycle \(\gamma\) determines a signed sum of oriented edges.  For \(a\in C^1(\Gamma;k)\), define the period of \(a\) around \(\gamma\) by
\[
\oint_\gamma a
=
\sum_{e\in \gamma}\epsilon_e a(e),
\]
where \(\epsilon_e=1\) if \(\gamma\) traverses \(e\) in the chosen orientation and \(\epsilon_e=-1\) if it traverses \(e\) in the opposite orientation.  The period audit is the map
\[
\Per(a):Z_1(\Gamma;k)\to k,
\qquad
\Per(a)(\gamma)=\oint_\gamma a,
\]
where \(Z_1(\Gamma;k)\) is the cycle space of the graph.
\end{definition}

Thus a period is the signed edge-sum of \(a\) along a cycle; the audit asks
whether all such sums vanish.

\begin{definition}[Local interface and cohomology interface]
\label{def:coh-interfaces}
Let
\[
S=C^1(\Gamma;k).
\]
Define the lower local interface by
\[
\pi_0:S\to\{\texttt{locally-admissible}\},
\qquad
\pi_0(a)=\texttt{locally-admissible},
\]
and the promoted cohomology interface by
\[
\pi_1:S\to H^1(\Gamma;k),
\qquad
\pi_1(a)=[a].
\]
\end{definition}

The lower interface only records that an edge datum is locally admissible, while
the promoted interface records its cohomology class.

\begin{definition}[Cohomological non-descent and audit]
\label{def:coh-nondesc-audit}
For \(a\in C^1(\Gamma;k)\), write
\[
a\not\downarrow_{\mathrm{pot}}
\]
when no global potential \(\phi\in C^0(\Gamma;k)\) satisfies \(d_0\phi=a\).  Define the obstruction expression
\[
C(a)=([a],\Per(a)).
\]
It is audited by the equivalences
\[
[a]=0
\quad\Longleftrightarrow\quad
\Per(a)=0
\quad\Longleftrightarrow\quad
\exists\phi\text{ with }d_0\phi=a.
\]
\end{definition}

This audit packages non-descent as a standard cohomological obstruction: the
edge datum descends to a potential exactly when its class and periods vanish.

\begin{definition}[Cohomology descent-use problem]
\label{def:coh-descent-use-problem}
For a fixed finite connected graph \(\Gamma\) and edge datum \(a\in C^1(\Gamma;k)\), the cohomology descent-use problem \(P_{\Coh}(a)\) asks for an accepted lower answer to the question:
\[
\text{Does }a\text{ descend to a global potential?}
\]
The lower candidate class consists of potentials \(\phi\in C^0(\Gamma;k)\).  The promoted answer type consists of either:
\begin{enumerate}
\item a positive descent witness \(\phi\) with \(d_0\phi=a\), certified by \([a]=0\); or
\item a negative obstruction witness, such as \([a]\neq0\) together with a cycle \(\gamma\) satisfying \(\oint_\gamma a\neq0\).
\end{enumerate}
The nonclaim register contains, at minimum,
\[
\text{local edge admissibility implies global descent}
\quad\text{is not claimed.}
\]
\end{definition}

The first calculation is the easy direction of the usual finite-graph
cohomology criterion: exact edge data have no cycle periods.

\begin{lemma}[Exact edge data have zero cycle periods]
\label{lem:coh-exact-zero-periods}
If \(a=d_0\phi\) for some \(\phi\in C^0(\Gamma;k)\), then for every cycle \(\gamma\),
\[
\oint_\gamma a=0.
\]
\end{lemma}

\begin{proof}
Write an oriented edge \(e:u\to v\).  If \(a=d_0\phi\), then
\(a(e)=\phi(v)-\phi(u)\), with the sign reversed when the edge is traversed in
the opposite orientation.  Along a closed cycle, the resulting signed sum is a
telescoping sum of successive vertex-potential differences, and the initial and
terminal vertex coincide.  Hence the period is zero.  This lemma is tracked by
the formalization harness as a projection over the exact-edge-data and
period-audit predicates.
\end{proof}

A nonzero period is therefore an obstruction to solving the lower problem by a
global potential.

\begin{corollary}[Cycle period obstruction]
\label{cor:coh-period-obstruction}
If there exists a cycle \(\gamma\) such that
\[
\oint_\gamma a\neq 0,
\]
then no global potential \(\phi\) satisfies
\[
d_0\phi=a.
\]
Consequently the lower problem profile \(\mathfrak P_{a,0}\) is saturated for the candidate class \(C^0(\Gamma;k)\).
\end{corollary}

\begin{proof}
This is the contrapositive of \Cref{lem:coh-exact-zero-periods}.  If a global
potential \(\phi\) satisfied \(d_0\phi=a\), every cycle period would vanish.
Thus a nonzero period forbids such a potential and saturates the lower candidate
class for this datum.  This corollary is tracked by the formalization harness as
a projection over the exact-zero-period lemma and the saturation predicate.
\end{proof}

The converse direction is the next step: on a finite graph, vanishing periods
are not merely necessary for exactness; they also characterize when the edge
datum comes from a global vertex potential.

\subsection{The \texorpdfstring{\(H^1\)}{H1} ledger and exactness}
\label{subsec:coh-h1-ledger}

The period obstruction was the easy direction: exact edge data have zero periods.
The converse uses path integration.  If all periods vanish, then integrating the
edge datum from a fixed base vertex to any other vertex is independent of the
chosen path, and those integrals assemble a global potential.

\begin{theorem}[Path-independence criterion for exactness]
\label{thm:coh-periods-iff-exact}
Let \(\Gamma\) be finite and connected.  For \(a\in C^1(\Gamma;k)\), the following are equivalent:
\[
\text{there exists }\phi\in C^0(\Gamma;k)\text{ with }d_0\phi=a,
\]
\[
\oint_\gamma a=0\quad\text{for every cycle }\gamma.
\]
\end{theorem}

\begin{proof}
One direction is \Cref{lem:coh-exact-zero-periods}.  Conversely, assume every
cycle period vanishes and choose a base vertex \(v_0\).  For a vertex \(v\), set
\(\phi(v)\) equal to the signed sum of \(a\) along any path from \(v_0\) to
\(v\).  If two paths from \(v_0\) to \(v\) are chosen, their difference is a
cycle, so the zero-period hypothesis gives the same value of \(\phi(v)\).  For
an oriented edge \(e:u\to v\), the path to \(u\) followed by \(e\) gives
\(\phi(v)-\phi(u)=a(e)\).  Hence \(d_0\phi=a\).  This theorem is tracked by the
formalization harness as a projection over the period-audit and exactness
predicates.
\end{proof}

The cohomology class is the corresponding quotient ledger: it records exactly
whether the edge datum survives modulo exact data.

\begin{theorem}[Cohomology class is the exact obstruction]
\label{thm:coh-class-obstruction}
For \(a\in C^1(\Gamma;k)\),
\[
[a]=0\text{ in }H^1(\Gamma;k)
\quad\Longleftrightarrow\quad
\exists \phi\in C^0(\Gamma;k)\text{ such that }d_0\phi=a.
\]
Equivalently,
\[
[a]=0
\quad\Longleftrightarrow\quad
a\in\im d_0.
\]
\end{theorem}

\begin{proof}
By definition, \(H^1(\Gamma;k)\) is the quotient of \(1\)-cochains by exact
\(1\)-cochains on a graph.  Thus the class \([a]\) is zero exactly when \(a\)
lies in the image of \(d_0:C^0(\Gamma;k)\to C^1(\Gamma;k)\), equivalently when
\(a=d_0\phi\) for some vertex potential \(\phi\).  This theorem is tracked by
the formalization harness as a projection over the quotient-ledger obstruction
predicate.
\end{proof}

The period audit and the quotient ledger therefore detect the same obstruction.

\begin{corollary}[Cohomology and periods give equivalent audits]
\label{cor:coh-class-period-equivalence}
For a finite connected graph,
\[
[a]=0
\quad\Longleftrightarrow\quad
\Per(a)=0.
\]
Thus the quotient ledger \([a]\) and the period audit record the same exactness obstruction.
\end{corollary}

\begin{proof}
The vanishing of \(\Per(a)\) says that every cycle period of \(a\) is zero, so
\Cref{thm:coh-periods-iff-exact} identifies it with exactness of \(a\).  The
class-obstruction theorem identifies exactness with \([a]=0\).  This corollary
is tracked by the formalization harness as a projection over the path-independence
criterion and the class-obstruction predicate.
\end{proof}

The rank of \(H^1\) is the number of independent cycle obstructions.  For a
connected finite graph, that number is the familiar cycle rank \(|E|-|V|+1\).

\begin{theorem}[Dimension of the graph obstruction ledger]
\label{thm:coh-dimension}
If \(\Gamma\) is finite and connected, then
\[
\dim_k H^1(\Gamma;k)=|E|-|V|+1.
\]
\end{theorem}

\begin{proof}
For a finite connected graph over a coefficient field \(k\), one has
\(\dim_k C^0(\Gamma;k)=|V|\) and \(\dim_k C^1(\Gamma;k)=|E|\).  The kernel of
\(d_0:C^0\to C^1\) consists of the constant potentials, so it has dimension
\(1\), and therefore \(\operatorname{rank} d_0=|V|-1\).  Since a graph has no
two-cells, \(H^1(\Gamma;k)\) is \(C^1/\im d_0\).  Rank-nullity gives
\[
\dim_k H^1(\Gamma;k)=|E|-(|V|-1)=|E|-|V|+1.
\]
This is the standard dimension computation for finite graph cohomology.  The
theorem is tracked by the formalization harness as a projection over the
dimension ledger and rank-nullity fields.
\end{proof}

The next subsection makes the ledger concrete with two examples: a triangle,
where \(H^1\) is nontrivial, and a forest, where every period obstruction
vanishes.

\subsection{Triangle and forest examples}
\label{subsec:coh-triangle-forest}

The triangle is the smallest connected graph with a cycle.  It has three
vertices, three oriented edges, and one cycle, so a single nonzero period already
witnesses failed descent.  The edge assignment below is deliberately sparse: one
edge carries value \(1\), the other two carry value \(0\).

\begin{proposition}[Triangle graph as a minimal obstruction witness]
\label{prop:coh-triangle-witness}
Let \(\Gamma\) be the oriented triangle with edges
\[
e_{12}:v_1\to v_2,
\qquad
e_{23}:v_2\to v_3,
\qquad
e_{31}:v_3\to v_1.
\]
Define \(a\in C^1(\Gamma;k)\) by
\[
a(e_{12})=1,
\qquad
a(e_{23})=0,
\qquad
a(e_{31})=0.
\]
Then \(a\) does not descend to a global potential.
\end{proposition}

\begin{proof}
Let \(\gamma\) be the oriented cycle
\[
v_1\to v_2\to v_3\to v_1.
\]
Then
\[
\oint_\gamma a
=
a(e_{12})+a(e_{23})+a(e_{31})
=1.
\]
This period is nonzero, so \Cref{cor:coh-period-obstruction} forbids a global
potential \(\phi\) with \(d_0\phi=a\).  The proposition is tracked by the
formalization harness as a projection over the triangle period-audit record.
\end{proof}

The same computation identifies the triangle ledger.  By the dimension formula,
\[
\dim_k H^1(\Gamma;k)=|E|-|V|+1=3-3+1=1,
\]
and the class of the displayed edge assignment is nonzero.  The formal
one-dimensional triangle ledger corollary is deferred to the methods appendix
with the other worked graph computations.

The contrast with a forest is shown in \Cref{fig:graph-cohomology}.  The figure
keeps the two decisions side by side: the triangle produces a nonzero period and
a nonzero \(H^1\) class, while the forest has no cycles and no obstruction.

\begingroup\begin{figure}[tbp]
\centering
\begin{tikzpicture}[
  >=stealth,
  font=\small,
  vertex/.style={circle, draw, inner sep=1.6pt, minimum size=5.5mm},
  audit/.style={draw, rounded corners=2pt, align=center, inner sep=4pt, text width=3.05cm},
  decision/.style={draw, align=center, inner sep=4pt, text width=2.75cm}
]
  \node[font=\bfseries] at (1.25,1.35) {triangle};
  \node[vertex] (t1) at (0,0.75) {};
  \node[vertex] (t2) at (2.5,0.75) {};
  \node[vertex] (t3) at (1.25,-0.85) {};
  \draw[->] (t1) -- node[above] {\(a_1\)} (t2);
  \draw[->] (t2) -- node[right] {\(a_2\)} (t3);
  \draw[->] (t3) -- node[left] {\(a_3\)} (t1);
  \node[audit] (taudit) at (1.25,-2.35) {\(\displaystyle \oint_\gamma a=a_1+a_2+a_3\)};
  \node[decision] (tclass) at (1.25,-3.75) {\([a]\in H^1(\Gamma;k)\)};
  \node[decision] (tdec) at (1.25,-5.15) {nonzero period\\blocks descent};
  \draw[->] (t3) -- (taudit);
  \draw[->] (taudit) -- (tclass);
  \draw[->] (tclass) -- (tdec);

  \node[font=\bfseries] at (7.1,1.35) {forest};
  \node[vertex] (f1) at (5.45,0.7) {};
  \node[vertex] (f2) at (6.85,0.7) {};
  \node[vertex] (f3) at (8.25,0.7) {};
  \node[vertex] (f4) at (6.85,-0.85) {};
  \draw[->] (f1) -- node[above] {\(b_1\)} (f2);
  \draw[->] (f2) -- node[above] {\(b_2\)} (f3);
  \draw[->] (f2) -- node[right] {\(b_3\)} (f4);
  \node[audit] (faudit) at (6.85,-2.35) {no cycles, so\\\(\Per(b)=0\)};
  \node[decision] (fclass) at (6.85,-3.75) {\([b]=0\)};
  \node[decision] (fdec) at (6.85,-5.15) {edge data\\descends to \(\phi\)};
  \draw[->] (f4) -- (faudit);
  \draw[->] (faudit) -- (fclass);
  \draw[->] (fclass) -- (fdec);
\end{tikzpicture}
\caption{Triangle and forest edge data illustrate period audit, the \(H^1\) class, and the resulting descent decision.}
\label{fig:graph-cohomology}
\end{figure}
 \endgroup

For a finite tree, there are no cycles, so the period audit is vacuous.  Given
edge data on a tree, choose a root vertex and define a potential by summing edge
values along the unique path from the root to each vertex.  This gives
\(d_0\phi=a\), up to the usual additive constant freedom.  Componentwise, the
same argument applies to any forest, so \(H^1\) vanishes on each component; the
formal forest-vanishing corollary is recorded in the methods appendix.

These examples show both sides of the obstruction ledger.  The next subsection
uses this contrast to prove strictness and assemble the cohomology SAU
certificate as the obstruction-ledger profile.

\subsection{Strictness and SAU certificate as obstruction profile}
\label{subsec:coh-strictness-certificate}

The cochain promotion is strict because the lower local interface forgets the
global obstruction class.  Two edge data can both be locally admissible while
only one carries a nonzero \(H^1\)-class.

\begin{theorem}[Strictness of cohomological promotion]
\label{thm:coh-strictness}
If \(H^1(\Gamma;k)\neq0\), then
\[
\delta_{\mathrm{fact}}(\pi_0,\pi_1)\neq\varnothing.
\]
Thus the cohomology layer is strict relative to the purely local edge-admissibility interface.
\end{theorem}

\begin{proof}
Choose \(a\in C^1(\Gamma;k)\) with \([a]\neq0\), and compare it with the zero
edge datum \(0\).  The lower interface sends both to
\(\texttt{locally-admissible}\), so \(\pi_0(a)=\pi_0(0)\).  The promoted
interface distinguishes them, since \(\pi_1(a)=[a]\neq0=[0]=\pi_1(0)\).  Thus
\((a,0)\in\delta_{\mathrm{fact}}(\pi_0,\pi_1)\).  This theorem is tracked by the
formalization harness as a projection over the cohomology-interface and
factorization-defect predicates.
\end{proof}

The resulting profile is an obstruction ledger: the useful descended answer is
not a hidden edge datum pretending to be a vertex potential, but the audited
class and period test deciding whether such a potential exists.

\begin{theorem}[Obstruction-ledger descent profile]
\label{thm:coh-obstruction-profile}
For a finite connected graph, the cohomology package has the obstruction-ledger profile:
\[
a\not\downarrow_{\mathrm{pot}}
\quad\Longleftrightarrow\quad
[a]\neq0
\quad\Longleftrightarrow\quad
\exists\gamma\text{ with }\oint_\gamma a\neq0.
\]
When \([a]=0\), the same package supplies a positive descent witness \(\phi\) with \(d_0\phi=a\).  When \([a]\neq0\), it supplies an audited obstruction certificate.
\end{theorem}

\begin{proof}
The equivalence between \(a\not\downarrow_{\mathrm{pot}}\) and \([a]\neq0\) is
the contrapositive form of the class-obstruction theorem.  The equivalence with
the existence of a nonzero cycle period follows from the class-period audit:
zero class is equivalent to zero period audit, so nonzero class is detected by a
nonzero period on some cycle.  If the class vanishes, path integration supplies a
potential \(\phi\) with \(d_0\phi=a\); if it does not vanish, the class and a
nonzero period give the negative obstruction certificate.  This theorem is
tracked by the formalization harness as a projection over the class-obstruction,
period-audit, and descent-use predicates.
\end{proof}

The SAU certificate is now the same six-field pattern as in the previous cases:
saturation of the requested potential for a nonzero class, strictness of the
cohomology interface, non-descent of the edge datum, audited obstruction descent,
usefulness of the negative certificate, and the nonclaim that local admissibility
alone implies global descent.

\begin{theorem}[Cohomology supplies an SAU obstruction certificate]
\label{thm:coh-sau-certificate}
Let \(\Gamma\) be a finite connected graph with \(H^1(\Gamma;k)\neq0\), and let \(a\in C^1(\Gamma;k)\) satisfy \([a]\neq0\).  Then the cohomology package gives a Strict Audited Utility certificate for the saturated obstruction instance of the descent-use problem:
\[
\mathsf{Coh}\models\mathsf{SAU}.
\]
More explicitly:
\begin{enumerate}
\item \textbf{Saturation:} when \([a]\neq0\), no lower candidate potential solves \(d_0\phi=a\);
\item \textbf{Strictness:} the cohomology interface distinguishes global obstruction classes collapsed by the local interface;
\item \textbf{Non-descent:} nonzero classes represent local data that do not descend to global potentials;
\item \textbf{Audited descent:} \([a]\), periods, and vanishing/nonvanishing tests decide descent under audit;
\item \textbf{Usefulness:} the package supplies a negative obstruction certificate for the saturated instance;
\item \textbf{No-overreading:} the certificate does not claim that local admissibility alone implies global descent.
\end{enumerate}
\end{theorem}

\begin{proof}
Saturation follows from the class-obstruction theorem: if \([a]\neq0\), no
potential \(\phi\) satisfies \(d_0\phi=a\).  Strictness is
\Cref{thm:coh-strictness}.  Non-descent and audited descent are exactly the
obstruction-ledger profile: the edge datum does not descend to a potential, but
the pair \(([a],\Per(a))\) audits the descent decision.  Usefulness is the
negative obstruction certificate supplied by the nonzero class or period, and
the nonclaim register blocks the overreading that local edge admissibility alone
implies global descent.  This theorem is tracked by the formalization harness as
a projection over the cohomology SAU certificate record and its constituent
witness fields.
\end{proof}

The four canonical witnesses are now in place.  The next section abstracts their
shared profile structure and proves the transfer principles that move an
audited profile from one case to another.
 \section{Transfer and Cross-Case Synthesis}\label{sec:transfer}

\subsection{The canonical witness vector \texorpdfstring{\(W_4\)}{W4}}
\label{subsec:transfer-canonical-witness-vector}

The four case studies now provide four validated certificates.  This section
abstracts what they share: not an identification of \(i\), \(\varepsilon\),
\(\alpha\), and \([a]\), but a common proof shape in which a hidden promoted
support fails to descend while a selected expression from that support does
descend under audit.

\begin{definition}[Canonical witness vector]
\label{def:sau-canonical-witness-vector}
The canonical witness vector is the ordered family
\[
\mathfrak W_{4}
=
(\mathfrak W_{\Imag},\mathfrak W_{\Jet},\mathfrak W_{\Galois},\mathfrak W_{\Coh}),
\]
where:
\begin{enumerate}
\item \(\mathfrak W_{\Imag}\) is the root-composite witness built from the four-phase carrier \(C_4=\{1,i,-1,-i\}\), the real-shadow lens \(q:C_4\to\{-1,0,1\}\), the quarter-turn \(G(z)=iz\), and the descending square \(G^2\).
\item \(\mathfrak W_{\Jet}\) is the coefficient witness built from the dual-number algebra \(\mathbb R[\varepsilon]/(\varepsilon^2)\), the non-real nilpotent \(\varepsilon\), the first-jet expression \(f(x+\varepsilon)\), and the coefficient audit \([\varepsilon]f(x+\varepsilon)=f'(x)\).
\item \(\mathfrak W_{\Galois}\) is the invariant-descent witness built from a finite Galois splitting field \(L/K\), a non-descending root \(\alpha\notin K\), and descending invariant audits such as elementary symmetric functions, \(L^{\Gal(L/K)}=K\), \(\Tr_{L/K}\), and \(\Nm_{L/K}\).
\item \(\mathfrak W_{\Coh}\) is the obstruction-ledger witness built from a finite connected graph \(\Gamma\), local edge data \(a\in C^1(\Gamma;k)\), the global-potential descent question \(a=d_0\phi\), and the obstruction class \([a]\in H^1(\Gamma;k)\).
\end{enumerate}
Each component of \(\mathfrak W_4\) is understood with its declared lower problem profile, promotion bridge, audit channels, and nonclaim register from the corresponding theorem sheet.
\end{definition}

The compression from case-specific mathematics to profile labels and then to a
shared boundary pattern is shown in \Cref{fig:cross-case-profiles}.

\begingroup\begin{figure}[tbp]
\centering
\begin{tikzpicture}[
  >=stealth,
  font=\small,
  case/.style={draw, rounded corners=2pt, align=center, inner sep=4pt, minimum width=1.55cm},
  profile/.style={draw, rounded corners=2pt, align=center, inner sep=4pt, minimum width=2.05cm, fill=black!3},
  normal/.style={draw, rounded corners=2pt, align=center, inner sep=5pt, text width=3.65cm}
]
  \node[case] (imag) at (0,1.8) {\(\Imag\)};
  \node[case] (jet) at (0,0.6) {\(\Jet\)};
  \node[case] (gal) at (0,-0.6) {\(\Galois\)};
  \node[case] (coh) at (0,-1.8) {\(\Coh\)};

  \node[profile] (root) at (3.0,1.8) {\(\RootComp\)};
  \node[profile] (coeff) at (3.0,0.6) {\(\Coeff\)};
  \node[profile] (inv) at (3.0,-0.6) {\(\InvDesc\)};
  \node[profile] (obs) at (3.0,-1.8) {\(\ObsLed\)};

  \node[normal] (vector) at (6.55,0.9) {canonical witness vector\\
  \((\Imag,\Jet,\Galois,\Coh)\)};
  \node[normal] (normal) at (6.55,-1.05) {shared SAU boundary form\\
  \(U\ndown_{\T_0}\),\quad \(C(U)\down_B a^\sharp\)};

  \draw[->] (imag) -- (root);
  \draw[->] (jet) -- (coeff);
  \draw[->] (gal) -- (inv);
  \draw[->] (coh) -- (obs);
  \draw[->] (root) -- (vector);
  \draw[->] (coeff) -- (vector);
  \draw[->] (inv) -- (normal);
  \draw[->] (obs) -- (normal);
  \draw[->] (vector) -- node[right] {normalizes to} (normal);
\end{tikzpicture}
\caption{The four canonical witnesses project to profile labels and then to the shared SAU boundary normal form.}
\label{fig:cross-case-profiles}
\end{figure}
 \endgroup

Each component of \(\mathfrak W_4\) carries the six fields of an SAU certificate.
The next four lemmas record those projections case by case.

\begin{lemma}[Imaginary witness projects to SAU fields]
\label{lem:sau-cross-imag-fields}
The witness \(\mathfrak W_{\Imag}\) is SAU-complete.  Its six-field projection is:
\[
\begin{array}{ll}
\mathsf{SatW}:& \text{no same-layer }H:\{-1,0,1\}\to\{-1,0,1\}\text{ satisfies }H^2(y)=-y,\\[2mm]
\mathsf{StrictW}:& \delta_{\mathrm{fact}}(q,\operatorname{id}_{C_4})\neq\varnothing\text{, witnessed by }(i,-i),\\[2mm]
\mathsf{NonDescW}:& G\not\downarrow_q,\\[1mm]
\mathsf{DescW}:& G^2\downarrow_q F\text{, where }F(y)=-y,\\[1mm]
\mathsf{SolveW}:& F\text{ is the lower reversal required by the half-step problem},\\[1mm]
\mathsf{AuditW}:& q\text{, conjugation, norm, and the nonclaim }G\in\mathcal T_0\text{ is not asserted}.
\end{array}
\]
Consequently,
\[
\Imag\models\SAU.
\]
\end{lemma}

\begin{proof}
The six fields are exactly the saturation, strictness, non-descent, descent,
accepted-answer, and nonclaim components established for the imaginary case in
\Cref{sec:imaginary}.  This lemma is tracked by the formalization harness as a
projection over the imaginary certificate record.
\end{proof}

\begin{lemma}[Calculus witness projects to SAU fields]
\label{lem:sau-cross-jet-fields}
The witness \(\mathfrak W_{\Jet}\) is SAU-complete.  Its six-field projection is:
\[
\begin{array}{ll}
\mathsf{SatW}:& \text{point value }\operatorname{ev}_x(f)=f(x)\text{ cannot determine }f'(x),\\[1mm]
\mathsf{StrictW}:& j_x^1\text{ distinguishes functions that }\operatorname{ev}_x\text{ collapses},\\[1mm]
\mathsf{NonDescW}:& \varepsilon\notin\mathbb R\text{ and no nonzero real }r\text{ has }r^2=0,\\[1mm]
\mathsf{DescW}:& [\varepsilon]f(x+\varepsilon)=f'(x),\\[1mm]
\mathsf{SolveW}:& f\mapsto f'(x)\text{ is the requested lower derivative functional},\\[1mm]
\mathsf{AuditW}:& \operatorname{st}\text{ and }[\varepsilon]\text{ audits, with nonclaim }\varepsilon\in\mathbb R\text{ blocked}.
\end{array}
\]
Consequently,
\[
\Jet\models\SAU.
\]
\end{lemma}

\begin{proof}
The value-only saturation, first-jet strictness, non-descent of
\(\varepsilon\), coefficient descent, accepted derivative functional, and
nonclaim discipline are the calculus certificate components developed in
\Cref{sec:calculus}.  This lemma is tracked by the formalization harness as a
projection over the calculus certificate record.
\end{proof}

\begin{lemma}[Galois witness projects to SAU fields]
\label{lem:sau-cross-gal-fields}
The witness \(\mathfrak W_{\Galois}\) is SAU-complete.  Its six-field projection is:
\[
\begin{array}{ll}
\mathsf{SatW}:& p\in K[t]\text{ irreducible nonlinear has no root in }K,\\[1mm]
\mathsf{StrictW}:& \text{the splitting field distinguishes roots the base polynomial interface collapses},\\[1mm]
\mathsf{NonDescW}:& \alpha\notin K\text{ for a root }\alpha\text{ of }p,\\[1mm]
\mathsf{DescW}:& e_j(\alpha_1,\ldots,\alpha_n),\ L^{\Gal(L/K)},\ \Tr_{L/K},\ \Nm_{L/K}\text{ descend to }K,\\[1mm]
\mathsf{SolveW}:& \text{root representation and invariant base-field computations are supplied upstairs},\\[1mm]
\mathsf{AuditW}:& \text{Galois invariance, fixed field, trace/norm, and nonclaim }\alpha\in K\text{ blocked}.
\end{array}
\]
Consequently,
\[
\Galois\models\SAU.
\]
\end{lemma}

\begin{proof}
The no-base-root saturation, splitting-field strictness, root non-descent,
invariant descent, accepted invariant computations, and no-literalization
discipline are the Galois certificate components assembled in
\Cref{sec:galois}.  This lemma is tracked by the formalization harness as a
projection over the Galois certificate record.
\end{proof}

\begin{lemma}[Cohomology witness projects to SAU fields]
\label{lem:sau-cross-coh-fields}
The witness \(\mathfrak W_{\Coh}\) is SAU-complete for the saturated obstruction instance.  Its six-field projection is:
\[
\begin{array}{ll}
\mathsf{SatW}:& [a]\neq0\text{ implies no global potential }\phi\text{ satisfies }d_0\phi=a,\\[1mm]
\mathsf{StrictW}:& H^1\text{ distinguishes obstruction classes the local-admissibility interface collapses},\\[1mm]
\mathsf{NonDescW}:& a\not\downarrow_{\mathrm{pot}}\text{ as global potential data},\\[1mm]
\mathsf{DescW}:& [a]\in H^1(\Gamma;k)\text{ audits descent failure and }[a]=0\iff a=d_0\phi,\\[1mm]
\mathsf{SolveW}:& [a]\neq0\text{ is the accepted lower negative obstruction certificate},\\[1mm]
\mathsf{AuditW}:& \text{period audit, quotient audit, and nonclaim that local data automatically glues}.
\end{array}
\]
Consequently,
\[
\Coh\models\SAU
\]
for the obstruction-ledger profile.
\end{lemma}

\begin{proof}
The cohomology section supplies the saturated obstruction instance, strict
cohomology interface, non-descending edge datum, class-and-period audit,
negative certificate, and nonclaim register in \Cref{sec:cohomology}.  This
lemma is tracked by the formalization harness as a projection over the
cohomology certificate record.
\end{proof}

\Cref{tab:cross-case-synthesis} summarizes the four hidden supports, their
descended audits, and their profile labels.

\begingroup\begin{table}[tbp]
\centering
\footnotesize
\caption{Cross-case synthesis by hidden support, descended audit, and profile.}
\label{tab:cross-case-synthesis}
\begin{tabularx}{\textwidth}{@{}>{\raggedright\arraybackslash}p{0.12\textwidth}X X >{\raggedright\arraybackslash}p{0.13\textwidth}@{}}
\toprule
Case & Hidden support & Descended audit & Profile \\
\midrule
\(\Imag\) &
The quarter-turn \(G\) on \(C_4\), a non-descending phase generator. &
The square of the quarter-turn descends through \(q\) as the real reversal \(F\). &
\(\RootComp\) \\
\(\Jet\) &
The infinitesimal \(\varepsilon\) in the dual-number algebra, which is not a real scalar. &
The coefficient channel applied to the first jet recovers the derivative \(D_x\). &
\(\Coeff\) \\
\(\Galois\) &
A non-base-field root \(\alpha\) carried in a splitting field. &
Trace, norm, orbit polynomials, and fixed-field invariants return values in \(K\). &
\(\InvDesc\) \\
\(\Coh\) &
A nonzero cohomology class \([a]\) recording edge data that do not come from a potential. &
Cycle periods and equality of \([a]\) with zero decide the descent question. &
\(\ObsLed\) \\
\bottomrule
\end{tabularx}
\end{table}
 \endgroup

\begin{theorem}[Four canonical SAU witnesses]
\label{thm:sau-four-canonical-witnesses}
The canonical witness vector satisfies
\[
\mathfrak W_4\models\SAU^4,
\]
that is,
\[
\Imag\models\SAU,
\qquad
\Jet\models\SAU,
\qquad
\Galois\models\SAU,
\qquad
\Coh\models\SAU.
\]
Consequently each of the four promoted packages licenses a strict audited usefulness claim under the soundness theorem of \Cref{sec:soundness-failure}.
\end{theorem}

\begin{proof}
Combine the four projection lemmas above.  Each component of
\(\mathfrak W_4\) supplies the six SAU witness fields and passes through the
soundness discipline of \Cref{sec:soundness-failure}.  This theorem is tracked
by the formalization harness as a projection over the four cross-case
projection records.
\end{proof}

\begin{theorem}[Uniform non-descent/descent normal form]
\label{thm:sau-uniform-normal-form}
For every witness
\[
\mathfrak W\in\{\mathfrak W_{\Imag},\mathfrak W_{\Jet},\mathfrak W_{\Galois},\mathfrak W_{\Coh}\},
\]
there exist a promoted structure \(U_{\mathfrak W}\), an expression former \(C_{\mathfrak W}\), a lower answer \(a^\sharp_{\mathfrak W}\), and a bridge \(B_{\mathfrak W}\) such that
\[
U_{\mathfrak W}\not\downarrow_{\mathcal T_0}
\qquad\text{but}\qquad
C_{\mathfrak W}(U_{\mathfrak W})\downarrow_{B_{\mathfrak W}} a^\sharp_{\mathfrak W}.
\]
The four instantiations are:
\[
\resizebox{\textwidth}{!}{\(
\begin{array}{c|c|c|c}
\text{case} & U_{\mathfrak W} & C_{\mathfrak W}(U_{\mathfrak W}) & a^\sharp_{\mathfrak W} \\
\hline
\Imag & G & G^2 & F(y)=-y \\
\Jet & \varepsilon & [\varepsilon]f(x+\varepsilon) & f'(x) \\
\Galois & \alpha\in L\setminus K & \Tr_{L/K}(\alpha),\ \Nm_{L/K}(\alpha),\ \text{invariants} & K\text{-valued audits} \\
\Coh & a\text{ not globally potential-generated} & [a]\in H^1(\Gamma;k) & \text{descent decision / obstruction certificate.}
\end{array}
\)
}
\]
\end{theorem}

\begin{proof}
Read the displayed data from the four cross-case projections.  In each case,
the hidden support is not accepted as a lower-layer object, while the listed
expression former supplies the audited descended answer.  This theorem is
tracked by the formalization harness as a projection over the canonical witness
vector and the SAU normal-form fields.
\end{proof}

With the four certificates collected, the next step is to forget substrate
details in a controlled way: profile abstraction keeps the audit shape while
discarding domain-specific object identity.

\subsection{Profile abstraction and substrate erasure}
\label{subsec:transfer-profile-abstraction}

The four cases share a structural pattern: saturation of a lower problem,
strictness of the promoted interface, a non-descending support, an audited
descended expression, an accepted lower answer, and a nonclaim discipline.  A
profile is the skeleton of that pattern.  It keeps the roles and audit
obligations while forgetting the concrete substrate that realizes them.

\begin{definition}[Case profile]
\label{def:sau-case-profile}
The case profile of a witness \(\mathfrak W\) is
\[
\Prof(\mathfrak W)
=(\tau_{\mathrm{nondesc}},\tau_{\mathrm{desc}},\tau_{\mathrm{audit}},\tau_{\mathrm{defect}}),
\]
where \(\tau_{\mathrm{nondesc}}\) records the type of non-descending structure, \(\tau_{\mathrm{desc}}\) records the type of selected descendant, \(\tau_{\mathrm{audit}}\) records the audit channel, and \(\tau_{\mathrm{defect}}\) records the lower-layer defect closed or certified by the promoted package.
\end{definition}

A case profile records the coarse shape of a single witness.  The next
definition refines that idea into the full list of SAU slots that must be
preserved under abstraction.

\begin{definition}[SAU profile skeleton]
\label{def:sau-profile-skeleton}
An \emph{SAU profile skeleton} is a finite typed record
\[
\mathfrak p=
(\tau_{0},\tau_{\mathrm{sat}},\tau_{\mathrm{strict}},
\tau_{U},\tau_{C},\tau_{\mathrm{desc}},\tau_{\mathrm{solve}},
\tau_{\mathrm{audit}},\tau_{\mathrm{nonclaim}}),
\]
where:
\begin{enumerate}
\item \(\tau_0\) records the lower-layer problem type;
\item \(\tau_{\mathrm{sat}}\) records the saturation witness type;
\item \(\tau_{\mathrm{strict}}\) records the nonfactorization witness type;
\item \(\tau_U\) records the promoted non-descending structure type;
\item \(\tau_C\) records the selected composite, coefficient, invariant, trace, norm, standard part, or obstruction expression;
\item \(\tau_{\mathrm{desc}}\) records the descent equation or audit channel;
\item \(\tau_{\mathrm{solve}}\) records how the descended object solves, decides, or certifies the lower problem;
\item \(\tau_{\mathrm{audit}}\) records the required audit data;
\item \(\tau_{\mathrm{nonclaim}}\) records the hidden-structure claims that must remain blocked.
\end{enumerate}
The skeleton records roles and obligations, not a concrete mathematical substrate.
\end{definition}

A profile instance is a proposed realization of such a skeleton by a concrete
candidate package.

\begin{definition}[Profile instance]
\label{def:sau-profile-instance}
Let \(\mathfrak p\) be an SAU profile skeleton.  A \emph{profile instance} of \(\mathfrak p\) is a raw candidate package
\[
\mathfrak C=(\mathfrak P_0,B_{0\to1},U,C,a^\sharp;
\mathsf{SatW},\mathsf{StrictW},\mathsf{NonDescW},\mathsf{DescW},\mathsf{SolveW},\mathsf{AuditW})
\]
together with a typing map
\[
\rho:\mathfrak C\longrightarrow\mathfrak p
\]
which assigns each concrete witness field of \(\mathfrak C\) to the corresponding profile slot of Definition~\ref{def:sau-profile-skeleton}.  We write
\[
\Inst(\mathfrak p,\mathfrak C)
\]
when the typing map is declared and all profile slots are inhabited by concrete witness data.
\end{definition}

Validation adds the same gate discipline used for proof packages: an instance is
not transferable merely because its slots are named.

\begin{definition}[Validated profile instance]
\label{def:sau-validated-profile-instance}
A profile instance \(\Inst(\mathfrak p,\mathfrak C)\) is \emph{validated} when
\[
\mathbf F_{\SAU}(\mathfrak C)=\varnothing
\]
under the declared instruments, thresholds, bridges, and nonclaim registers.  We write
\[
\Inst_{\SAU}(\mathfrak p,\mathfrak C)
\]
for a validated profile instance.
\end{definition}

Substrate erasure is the abstraction operation: it drops the concrete carrier,
field, graph, algebra, element, or cochain while retaining the proof roles.

\begin{definition}[Substrate erasure]
\label{def:sau-substrate-erasure}
For a raw candidate package \(\mathfrak C\), define its \emph{substrate erasure}
\[
\Erase(\mathfrak C)=\Prof(\mathfrak C)
\]
to be the record obtained by forgetting the substrate-specific carriers, elements, functions, roots, graphs, fields, and algebras, while retaining:
\[
(\tau_{0},\tau_{\mathrm{sat}},\tau_{\mathrm{strict}},
\tau_U,\tau_C,\tau_{\mathrm{desc}},\tau_{\mathrm{solve}},
\tau_{\mathrm{audit}},\tau_{\mathrm{nonclaim}}).
\]
Equivalently, \(\Erase\) keeps the bridge grammar and suppresses the material implementation.
\end{definition}

The notation \(\sim_{\Prof}\) will be formalized in the next subsection.  Here
the key point is already visible: having the same erased profile is not the same
thing as identifying the hidden mathematical objects.

\begin{lemma}[Profile equivalence does not imply object identity]
\label{lem:sau-profile-not-object-identity}
If \(\mathfrak C\sim_{\Prof}\mathfrak C'\), it does not follow that the promoted structures \(U\) and \(U'\) are equal, isomorphic, or members of the same lower theory layer.
\end{lemma}

\begin{proof}
Profile comparison is made after erasing substrate-specific carriers, elements,
functions, roots, graphs, fields, and algebras.  Equality or isomorphism of the
hidden promoted structures is not one of the retained fields.  Thus two packages
can share the same abstract audit roles while their hidden supports live in
different substrates, and the profile relation supplies no membership claim in a
common lower theory layer.  This lemma is tracked by the formalization harness
as a projection over the substrate-erasure and profile-comparison records.
\end{proof}

This is the sense in which abstraction is erasure under audit: the erasure is
useful only because the audit-relevant fields remain present.

\begin{proposition}[Abstraction as audited erasure]
\label{prop:sau-abstraction-as-erasure}
In the SAU setting, abstraction is useful when \(\Erase\) removes substrate-specific implementation details while preserving the fields needed for SAU validation.  In particular, if
\[
\Inst_{\SAU}(\mathfrak p,\mathfrak C),
\]
then the abstract profile \(\mathfrak p=\Erase(\mathfrak C)\) preserves enough information to restate the six proof obligations for any new target instance.
\end{proposition}

\begin{proof}
The validated instance hypothesis says that all six SAU witness fields have
passed their declared gates.  Substrate erasure removes the material
implementation but keeps the slot types for saturation, strictness, non-descent,
descent, solving, audit, and nonclaim discipline.  Therefore the erased profile
retains precisely the obligations that a new target instance must reinhabit and
revalidate.  This proposition is tracked by the formalization harness as a
projection over validated-profile-instance and substrate-erasure predicates.
\end{proof}

With profile abstraction defined, the next subsection makes profile equivalence
and transfer plans precise.

\subsection{Profile equivalence and transfer plan}
\label{subsec:transfer-equivalence-plan}

Two profiles are equivalent when their skeletons match after substrate erasure
and their audit obligations can be reprojected.  A transfer plan is the proposed
data for using a validated source profile as a guide for a target package.  It
does not by itself validate the target.

\begin{definition}[Profile equivalence]
\label{def:sau-profile-equivalence}
Two candidate packages \(\mathfrak C\) and \(\mathfrak C'\) are \emph{profile-equivalent}, written
\[
\mathfrak C\sim_{\Prof}\mathfrak C',
\]
when their substrate erasures are isomorphic as typed SAU skeletons.  This means that their six witness obligations have the same role shape, but it does not assert an isomorphism between their underlying carriers.
\end{definition}

A transfer plan names a source, a target, and the skeleton they are meant to
share.  The target must still provide its own witnesses and instruments.

\begin{definition}[SAU transfer plan]
\label{def:sau-transfer-plan}
An \emph{SAU transfer plan} from a source package \(\mathfrak C_s\) to a target package \(\mathfrak C_t\) is a triple
\[
\Transfer(\mathfrak C_s,\mathfrak C_t;\mathfrak p)
\]
such that:
\begin{enumerate}
\item \(\Inst_{\SAU}(\mathfrak p,\mathfrak C_s)\), so the source realizes \(\mathfrak p\) with a validated certificate;
\item \(\Inst(\mathfrak p,\mathfrak C_t)\), so the target inhabits the same profile skeleton;
\item the target has its own candidate witness fields and instruments, not merely borrowed source witnesses.
\end{enumerate}
The transfer plan is \emph{validated} when additionally
\[
\mathbf F_{\SAU}(\mathfrak C_t)=\varnothing.
\]
\end{definition}

The canonical registry is the finite list of profiles already inhabited by the
four case studies.

\begin{definition}[Canonical SAU profile registry]
\label{def:sau-canonical-profile-registry}
The four canonical witnesses determine the registry
\[
\mathcal R_4=
\{\mathfrak p_{\RootComp},\mathfrak p_{\Coeff},\mathfrak p_{\InvDesc},\mathfrak p_{\ObsLed}\},
\]
where:
\begin{enumerate}
\item \(\mathfrak p_{\RootComp}\) is the root-composite profile:
\[
U\not\downarrow_q,
\qquad
U^n\downarrow_q F.
\]
\item \(\mathfrak p_{\Coeff}\) is the coefficient or standard-part profile:
\[
U\not\downarrow_{\mathcal T_0},
\qquad
\kappa(\Phi(U))=a^\sharp
\]
for a coefficient, standard-part, or jet audit \(\kappa\).
\item \(\mathfrak p_{\InvDesc}\) is the invariant-descent profile:
\[
U\notin \mathcal T_0,
\qquad
C(U)\in \operatorname{Inv}(\mathcal T_1/\mathcal T_0)\cong \mathcal T_0.
\]
\item \(\mathfrak p_{\ObsLed}\) is the obstruction-ledger profile:
\[
\text{local or partial data fails descent},
\qquad
[U]=0\iff\text{descent succeeds}.
\]
\end{enumerate}
\end{definition}

Coverage is the registry-level question: whether a family of lower problems has
validated instances in the finite profile registry.

\begin{definition}[SAU coverage]
\label{def:sau-coverage}
Let \(\mathcal P\) be a family of lower-layer problems and let \(\mathcal R\) be a finite registry of SAU profile skeletons.  We say that \(\mathcal R\) \emph{SAU-covers} \(\mathcal P\), written
\[
\Cover_{\SAU}(\mathcal R,\mathcal P),
\]
when for every \(P\in\mathcal P\) there exist a lower problem profile \(\mathfrak P_{0,P}\), a candidate bridge \(B_P\), a candidate package \(\mathfrak C_P\), and a profile \(\mathfrak p_P\in\mathcal R\) such that
\[
\Inst_{\SAU}(\mathfrak p_P,\mathfrak C_P).
\]
\end{definition}

Transfer soundness is deliberately conditional.  A source profile can guide a
target construction, but the target accepted claim is licensed only after the
target's own validation succeeds.

\begin{theorem}[Profile-transfer soundness]
\label{thm:sau-profile-transfer-soundness}
Let \(\Transfer(\mathfrak C_s,\mathfrak C_t;\mathfrak p)\) be a validated SAU transfer plan.  Then the target bridge in \(\mathfrak C_t\) licenses an accepted usefulness claim:
\[
\Use(B_t,P_t):\texttt{accepted}.
\]
Moreover, the accepted target claim is licensed by the target validation, not by literal identity with the source structure.
\end{theorem}

\begin{proof}
The transfer plan gives a validated source instance and a target instance in the
same profile skeleton.  Validation of the target plan is the additional
condition \(\mathbf F_{\SAU}(\mathfrak C_t)=\varnothing\), so the target has its
own saturation, strictness, non-descent, descent, solving, audit, and nonclaim
witnesses.  Applying the soundness theorem of \Cref{sec:soundness-failure} to
the target package licenses \(\Use(B_t,P_t):\texttt{accepted}\).  The source
validation explains why the profile was a legitimate guide, but it does not
identify the source and target objects.  This theorem is tracked by the
formalization harness as a projection over the transfer-plan and validated-target
records.
\end{proof}

The next subsection packages this conditional transfer discipline into the
cross-case exact-six and synthesis theorems.

\subsection{Cross-case exact-six and synthesis}
\label{subsec:transfer-cross-case-synthesis}

The transfer plan now has two jobs.  First, it aligns the six recurring audit
roles across the four witnesses.  Second, it records the theorem-level synthesis
that those aligned roles support.

\begin{theorem}[Exact-six activation across the four witnesses]
\label{thm:sau-cross-exact-six}
Each component of \(\mathfrak W_4\) admits an exact-six SAU refinement in the inherited \(\mathsf{FATCD}\)-anchored sense: for its associated description \(D_{\mathfrak W}\), each channel has \(\operatorname{Status}_i(D_{\mathfrak W})=\texttt{active\_projection}\) and an active projection \(\pi_i(D_{\mathfrak W})\).  In every case the roles have the following aligned meanings:
\[
\begin{array}{c|l}
P_1 & \text{descent obstruction and audited descent recovery},\\
P_2 & \text{lower nonrepresentability and promoted representability},\\
P_3 & \text{suppressed route, phase, orbit, path, or coherence data},\\
P_4 & \text{staged movement from lower layer to promoted layer and back to audit},\\
P_5 & \text{package, quotient, completion, or splitting construction},\\
P_6 & \text{audit, replay, provenance, no-overread, and nonclaim discipline}.
\end{array}
\]
\end{theorem}

\begin{proof}
For each of \(\Imag\), \(\Jet\), \(\Galois\), and \(\Coh\), the cross-case
projection lemma identifies the same six witness roles: saturation, strictness,
non-descent, descent, solve, and audit.  The exact-six statement repackages
those roles into the inherited \(P_1,\ldots,P_6\) channel vocabulary.  This
theorem is tracked by the formalization harness as a projection over the four
case certificate records and the exact-six channel map.
\end{proof}

The synthesis statement is not a new universal claim about every useful
extension.  It is the common consequence of the four validated witnesses and the
soundness discipline.

\begin{theorem}[Cross-case synthesis: useful non-definable extensions]
\label{thm:sau-cross-synthesis}
The four canonical witnesses support the following theorem-level synthesis:

\medskip
\noindent
\emph{A promoted mathematical layer is useful for a saturated lower-layer problem when it supplies structure that the lower layer cannot represent or descend directly, but a selected composite, coefficient, invariant, trace, norm, standard part, or obstruction class descends under audit and resolves the lower problem.}

\medskip
Formally, this synthesis is the common consequence of the soundness theorem of \Cref{sec:soundness-failure} and Theorem~\ref{thm:sau-four-canonical-witnesses}.
\end{theorem}

\begin{proof}
The four-canonical-witnesses theorem supplies the validated \(\Imag\), \(\Jet\),
\(\Galois\), and \(\Coh\) instances.  The soundness theorem of
\Cref{sec:soundness-failure} licenses the accepted lower-layer usefulness claim
for each validated package, while the profile and normal-form results identify
the common non-descent plus audited-descent shape.  This theorem is tracked by
the formalization harness as a projection over the soundness and canonical-witness
records.
\end{proof}

Finally, the four registry profiles are not merely formal names: each one has a
validated canonical inhabitant.

\begin{theorem}[Canonical profiles are inhabited]
\label{thm:sau-canonical-profiles-inhabited}
Each profile in \(\mathcal R_4\) has a validated canonical instance:
\[
\Inst_{\SAU}(\mathfrak p_{\RootComp},\mathfrak W_{\Imag}),
\]
\[
\Inst_{\SAU}(\mathfrak p_{\Coeff},\mathfrak W_{\Jet}),
\]
\[
\Inst_{\SAU}(\mathfrak p_{\InvDesc},\mathfrak W_{\Galois}),
\]
\[
\Inst_{\SAU}(\mathfrak p_{\ObsLed},\mathfrak W_{\Coh}).
\]
\end{theorem}

\begin{proof}
The canonical profile registry assigns \(\RootComp\), \(\Coeff\), \(\InvDesc\),
and \(\ObsLed\) to the four case-study profiles.  The case projections and
\Cref{thm:sau-four-canonical-witnesses} show that the corresponding witness
packages are validated SAU instances.  This theorem is tracked by the
formalization harness as a projection over the registry and canonical-witness
records.
\end{proof}

The final subsection now states the transfer-profile theorem and the paper's
final synthesis.

\subsection{Final synthesis statement}
\label{subsec:transfer-final-synthesis}

The transfer theorem is the section's final guardrail.  A source case may supply
a profile skeleton, but accepted usefulness in a target case requires target-side
validation under the target instruments.

\begin{theorem}[Transfer profile theorem]
\label{thm:sau-transfer-profile-theorem}
Mathematical transfer at SAU strength is possible exactly at the profile level: a successful source case may suggest a reusable skeleton \(\mathfrak p\), but a target case becomes accepted only when it independently validates an instance of \(\mathfrak p\).  Formally, for a source \(\mathfrak C_s\), target \(\mathfrak C_t\), and profile \(\mathfrak p\),
\[
\Inst_{\SAU}(\mathfrak p,\mathfrak C_s)
+\Inst_{\SAU}(\mathfrak p,\mathfrak C_t)
\quad\Rightarrow\quad
\Use(B_s,P_s):\texttt{accepted}
\text{ and }
\Use(B_t,P_t):\texttt{accepted}.
\]
But
\[
\Inst_{\SAU}(\mathfrak p,\mathfrak C_s)
+\Inst(\mathfrak p,\mathfrak C_t)
\]
without target validation does not license \(\Use(B_t,P_t):\texttt{accepted}\).
\end{theorem}

\begin{proof}
The positive implication applies the profile-transfer soundness theorem to the
source and target validated instances separately.  Each package has its own empty
failure vector and therefore its own accepted usefulness claim under the
soundness discipline of \Cref{sec:soundness-failure}.  The negative clause is the
same condition read contrapositively: an unvalidated target instance may share a
profile skeleton, but it has not supplied the target-side witnesses required to
license \(\Use(B_t,P_t):\texttt{accepted}\).  This theorem is tracked by the
formalization harness as a projection over validated-profile-instance and
transfer-plan records.
\end{proof}

The final synthesis states what the preceding sections establish at the intended
level of generality: recurrence of audited profiles, not identity of hidden
objects across mathematical domains.

\begingroup
\hfuzz=150pt
\begin{theorem}[SAU synthesis: usefulness through recurring audited profiles]
\label{thm:sau-final-synthesis}
Within the scope of the theorem stack developed in this paper, the usefulness of mathematics is explained by the recurrence of strict audited descent profiles across theory layers.  More precisely, whenever a lower problem family \(\mathcal P\) is SAU-covered by a finite profile registry \(\mathcal R\), the usefulness claims for the family are licensed by:
\[
\text{lower-layer saturation},
\quad
\text{strict promotion},
\quad
\text{non-descending promoted structure},
\quad
\text{audited descending output},
\quad
\text{no-overreading}.
\]
The resulting explanation is profile-level rather than object-level: mathematical structures become portable because their descent profiles can recur across distinct substrates.
\end{theorem}
\endgroup

\begin{proof}
Coverage supplies a validated profile instance for each problem in the family.
For each such instance, the soundness theorem of \Cref{sec:soundness-failure}
licenses the accepted usefulness claim from the listed fields.  The transfer
results above show that what recurs is the profile skeleton and its audit
obligations, while substrate erasure and no-overreading prevent the recurrence
from becoming an identification of hidden objects.  This theorem is tracked by
the formalization harness as a projection over coverage, soundness, and
profile-transfer records.
\end{proof}

\[
\fbox{\begin{minipage}{0.9\textwidth}
In finite typed audited promoted-use traces, strict mathematical usefulness
forces hidden non-descent plus audited boundary descent; the four case studies
instantiate central minimal or universal profile types.
\end{minipage}}
\]

The outlook now separates what this theorem stack establishes from the broader
programmatic questions it leaves open.
 \section{Outlook and Contribution}\label{sec:outlook}

\subsection{What the paper unlocks}
\label{subsec:outlook-unlocks}

Within the scope of finite typed audited promoted-use traces, the paper gives a
profile-level explanation of strict mathematical usefulness.  What this unlocks
for further work is not a catalog of future domains, but a disciplined way to ask
whether a family of lower problems is covered by recurring audited profiles.

\begin{definition}[Explanatory compression]
\label{def:sau-explanatory-compression}
Suppose a finite registry \(\mathcal R\) SAU-covers a finite problem family \(\mathcal P\).  The registry provides \emph{explanatory compression} when many problem instances are certified by fewer recurring profile skeletons.  A crude count is
\[
\chi_{\SAU}(\mathcal P,\mathcal R)=|\mathcal P|-|\mathcal R|,
\]
with the understanding that this number is only a bookkeeping index, not an invariant of mathematical depth.
\end{definition}

The point of compression is not that one profile proves more than it contains.
It is that one validated audit shape may organize many lower problems.

\begin{proposition}[Why abstraction helps]
\label{prop:sau-why-abstraction-helps}
If \(\mathcal R\) SAU-covers \(\mathcal P\) and \(|\mathcal R|<|\mathcal P|\), then the registry gives a compressed explanation of the family-level usefulness of the corresponding mathematics: different lower-layer problems are solved or certified by reusing fewer audited bridge profiles.
\end{proposition}

\begin{proof}
The hypothesis says that every problem in \(\mathcal P\) has a validated
profile instance in the finite registry \(\mathcal R\).  When
\(|\mathcal R|<|\mathcal P|\), at least one profile skeleton is reused across
more than one problem instance, so the explanatory load is carried by fewer
audited bridge profiles than by problem instances.  This proposition is tracked
by the formalization harness as a projection over the coverage and explanatory
compression records.
\end{proof}

This also explains the paper's portability claim.  A profile can recur across
domains without turning the hidden objects in those domains into the same object.

\begin{corollary}[Portability without collapse]
\label{cor:sau-portability-without-collapse}
SAU explains why the same mathematical method can appear in different domains without claiming that the domains are literally the same.  What recurs is
\[
\text{descent profile}+\text{audit grammar},
\]
not necessarily substrate, ontology, or hidden object identity.
\end{corollary}

\begin{proof}
Substrate erasure preserves the audit pattern while forgetting the concrete
carrier, element, field, algebra, graph, or cochain that realizes it.  Therefore
recurrence across domains is recurrence of descent profile and audit grammar,
not literal identity of substrates or hidden promoted objects.  This corollary is
tracked by the formalization harness as a projection over substrate-erasure and
profile-transfer records.
\end{proof}

\Cref{tab:risks-mitigation} summarizes the main reviewer objections and the
body sections where the corresponding mitigation is supplied.

\begingroup\begin{table}[tbp]
\centering
\footnotesize
\caption{Reviewer objections and mitigation routes.}
\label{tab:risks-mitigation}
\begin{tabularx}{\textwidth}{@{}>{\raggedright\arraybackslash}p{0.24\textwidth}>{\raggedright\arraybackslash}X>{\raggedright\arraybackslash}p{0.25\textwidth}@{}}
\toprule
Objection or risk & Mitigation in paper & Source section \\
\midrule
Classical math dressed up & The classical constructions are inputs.  The new claim is the audited-descent synthesis: strict promoted-use traces force hidden non-descent plus audited boundary descent, and the four cases instantiate distinct minimal or universal profiles. & \Cref{sec:intro,sec:necessity,sec:transfer} \\
\addlinespace
Non-descent is only image failure & At the raw level, non-descent includes image failure.  The contribution is the typed certificate role pairing it with strictness, audit, no-overreading, and an accepted lower-layer answer. & \Cref{sec:layers,sec:soundness-failure} \\
\addlinespace
Audit smuggles the answer & Source policy, audit policy, no-smuggling defect, and nonclaim register are explicit; a bridge returns only declared descendants, not arbitrary hidden support. & \Cref{sec:layers,sec:soundness-failure} \\
\addlinespace
Hidden object literalized downstairs & No-overreading is a theorem: the audit licenses only the descended expression, not a lower-layer membership claim about the hidden object. & \Cref{sec:soundness-failure,sec:imaginary,sec:calculus,sec:galois,sec:cohomology} \\
\addlinespace
Why these four examples & The four cases are validated witnesses for central profile labels.  Further domains require their own validated profile instances. & \Cref{sec:classification,sec:transfer} \\
\addlinespace
Wigner/Hamming reframing & The question is mechanism-specific: why particular non-lower-layer objects become useful, not why mathematics is effective in general. & \Cref{sec:intro,sec:outlook} \\
\addlinespace
Notation too internal & Native vocabulary is introduced under a first-use translation rule and glossed in ordinary mathematical language before formal use. & Body sections; \Cref{app:definitions} \\
\addlinespace
Lean does not prove the whole paper & Lean wording is disclosed by coverage type: substantive proof language is reserved for direct, computational, or definitional rows; projection rows are traceability records. & \Cref{app:formalization} \\
\addlinespace
Minimality and universality overbroad & Minimality and universality are profile-relative: fixed lower problem, hidden-support role, audit expression, and no-overreading discipline. & \Cref{sec:imaginary,sec:calculus,sec:galois,sec:cohomology,sec:transfer} \\
\addlinespace
Transfer means same object across domains & Transfer is profile reuse under audited boundaries, not identity of promoted objects across substrates. & \Cref{sec:transfer,thm:sau-profile-transfer-soundness} \\
\addlinespace
Overclaiming further-domain extensions & The claim is scope-fenced; further-domain extensions are proposals until their target profile instances validate. & \Cref{sec:transfer,sec:outlook} \\
\bottomrule
\end{tabularx}
\end{table}
 \endgroup

The next subsection states the contribution and limits in prose.

\subsection{Contribution statement and limits}
\label{subsec:outlook-contribution-limits}

The contribution is a theorem-grade normal form for strict mathematical
usefulness in finite typed audited promoted-use traces.  When a saturated lower
problem is solved or certified by essentially using promoted-only structure, the
trace forces hidden support that does not descend directly and an audited
boundary expression that does descend.  The four canonical witnesses--imaginary
numbers, dual numbers, Galois roots, and graph cohomology--instantiate distinct
minimal or universal profiles, and the necessity-soundness chain extracts SAU
from strict useful traces rather than stipulating it.  The no-overreading
discipline is part of the contribution: it blocks the literalization of hidden
promoted objects in the lower layer.

The paper does not claim that all mathematics is useful, that every promoted
extension is useful, or that every useful structure has been classified.  It does
not claim novelty for the classical constructions themselves, and it does not
claim that hidden promoted objects literally belong to the lower layer.
Foundations I--III supply inherited provenance but are not prerequisites for
reading the paper, because the working vocabulary is defined locally.  The Lean
development also has a bounded role: it supplies substantive proofs for selected
direct, computational, and definitional fragments, while projection-over-record
theorems are tracked as proof-carrying traceability records.

The final assertion is therefore deliberately scoped:
\[
\boxed{\parbox{0.86\textwidth}{\centering
Within finite typed audited promoted-use traces, strict mathematical usefulness
forces hidden non-descent plus audited boundary descent.}}
\]
 
\appendix
\section{Definition Bank}\label{app:definitions}

\subsection{Basic vocabulary}
\label{appsubsec:definitions-basic-vocabulary}

This appendix collects canonical formal statements for definitions that the body
uses but chooses to summarize informally.  The body gives the reader the
operational meaning at the point of use; the appendix records the exact
set-theoretic or typed statement when that statement is not already promoted to
the body.

The first deferred item is the exact descent condition behind the commuting
square in \Cref{sec:layers}.

\begin{definition}[Exact quotient descent]
\label{def:sau-exact-descent}
Let \(q:X\to Y\) be a quotient, shadow, or observable map and let \(F:X\to X\) be an update. We say that \(F\) descends through \(q\), written
\[
F\downarrow_q,
\]
when there exists a map \(F^\sharp:Y\to Y\) such that
\[
qF=F^\sharp q.
\]
The descended map \(F^\sharp\) is unique on \(\operatorname{im}(q)\).
\end{definition}

The split-pair defect is the finite obstruction test for exact descent.

\begin{definition}[Split-pair descent defect]
\label{def:sau-split-defect}
For \(q:X\to Y\) and \(F:X\to X\), define
\[
\delta^{\mathrm{split}}_1(q,F)
=
\{(x,x')\in X\times X:q(x)=q(x')\text{ and }qF(x)\neq qF(x')\}.
\]
Then
\[
F\downarrow_q
\quad\Longleftrightarrow\quad
\delta^{\mathrm{split}}_1(q,F)=\varnothing.
\]
A nonempty split-pair defect is an explicit \(P_1\) descent obstruction.
\end{definition}

The strictness defect records when the promoted interface distinguishes data
collapsed by the lower interface.

\begin{definition}[Strictness defect]
\label{def:sau-strictness-defect}
Given object maps \(\pi_0:S\to O_0\) and \(\pi_1:S\to O_1\), define the factorization defect
\[
\delta_{\mathrm{fact}}(\pi_0,\pi_1)
=
\{(s,s')\in S\times S:\pi_0(s)=\pi_0(s')\text{ and }\pi_1(s)\neq\pi_1(s')\}.
\]
The promotion is strict when
\[
\operatorname{Strict}(B_{0\to1})
\quad\Longleftrightarrow\quad
\delta_{\mathrm{fact}}(\pi_0,\pi_1)\neq\varnothing.
\]
Equivalently, \(\pi_1\) does not factor through \(\pi_0\). A strict promotion distinguishes something that the lower interface collapses.
\end{definition}

The body already contains the formal definitions of theory package, instrument,
lower problem profile, saturation, accepted lower answer, promotion bridge,
non-descending promoted structure, audited descending expression, and Strict
Audited Utility certificate.  The appendix-bound definitions above supply the
descent and strictness mechanics deferred from \Cref{sec:layers}; the remaining
appendix subsections collect certificate grammar and failure/counterfeit
vocabulary.

\subsection{Certificate grammar}
\label{appsubsec:definitions-certificate-grammar}

This subsection records the appendix-bound certificate grammar that is not
duplicated in the body.  It starts with the role convention behind the six
certificate channels, then records the useful-extension and exact-six variants,
and closes with the profile-shape definitions used by the classification and
transfer sections.

\begin{definition}[Six Birds role convention]
\label{def:sau-roles}
Throughout this paper the six primitive role labels are used with the following fixed meanings:
\[
P_1=\text{descent/closure/rewrite/transport of content},\qquad
P_2=\text{representability/gating/admissibility},
\]
\[
P_3=\text{route mismatch/holonomy/path or protocol comparison},\qquad
P_4=\text{refinement/staging/gluing},
\]
\[
P_5=\text{packaging/completion/objecthood/idempotence},\qquad
P_6=\text{audit/provenance/replay/drive certification}.
\]
A role label is not a total algebraic generator. It is a typed channel whose activation requires role-aligned witness data, threshold data, level data, visibility data, audit data, and an explicit nonclaim record.
\end{definition}

The useful-extension definition names the accepted claim form associated with a
validated SAU certificate.

\begin{definition}[Strictly useful promoted extension]
\label{def:sau-useful-extension}
A promotion bridge \(B_{0\to1}\) is a strictly useful promoted extension of \(\mathcal T_0\) for the lower problem \(P\) when there exist \(U,C,a^\sharp\) such that
\[
\mathsf{SAU}(\mathfrak P_0,B_{0\to1},U,C,a^\sharp)
\]
holds. The accepted claim form is
\[
\Gamma;\mathcal T_0\leadsto\mathcal T_1;I_B
\vdash
\operatorname{Useful}(B_{0\to1},P):\texttt{accepted}.
\]
The claim means: the promoted layer solves, decides, or certifies the lower problem through an audited descending output that the saturated lower interface could not provide. It does not mean that the hidden promoted structure itself belongs to the lower layer.
\end{definition}

The exact-six variant strengthens the body certificate by requiring all six
role channels to carry active projected witness data.

\begin{definition}[Exact-six SAU certificate]
\label{def:sau-exact-six}
An SAU certificate is exact-six when, in addition to Definition~\ref{def:sau-certificate}, the bridge record determines an inherited \(\mathsf{FATCD}\)-style description \(D_B\) whose six role channels all have active projections:
\[
\operatorname{Status}_i(D_B)=\texttt{active\_projection},
\qquad
\pi_i(D_B)\ \text{exists},
\qquad
i=1,\ldots,6.
\]
Equivalently, the bridge record contains role-aligned active witness data for all six channels:
\[
(W_1,W_2,W_3,W_4,W_5,W_6),
\]
where:
\begin{itemize}
\item \(W_1\) witnesses descent or descent obstruction;
\item \(W_2\) witnesses lower nonrepresentability and promoted representability;
\item \(W_3\) witnesses route, phase, orbit, path, or coherence data suppressed by the lower layer;
\item \(W_4\) witnesses the stage transition from lower layer to promoted layer and back to audited descent;
\item \(W_5\) witnesses packaging, quotienting, or completion;
\item \(W_6\) witnesses audit, provenance, replay, no-overreading, and nonclaims.
\end{itemize}
Each witness must carry threshold, level, visibility, audit, and nonclaim records. Exact-six SAU is a scoped \(\mathsf{FATCD}\)-anchored certificate type, not a claim that every useful mathematical extension activates all six channels.
\end{definition}

The first profile-shape definition covers descending composites of
non-descending generators.

\begin{definition}[Root-composite profile]
\label{def:sau-root-profile}
A bridge has root-composite profile when
\[
U\not\downarrow_q
\qquad\text{but}\qquad
U^n\downarrow_B a^\sharp
\]
for some \(n\ge2\). The imaginary-number witness is the case \(n=2\), where the non-descending half-step squares to a descending reversal.
\end{definition}

The coefficient profile records audits that extract lower-layer content from a
nilpotent or infinitesimal promoted direction.

\begin{definition}[Coefficient or standard-part profile]
\label{def:sau-coeff-profile}
A bridge has coefficient/standard-part profile when a promoted infinitesimal or nilpotent object \(\eta\notin\mathcal T_0\) is used to form a promoted expression \(F(\eta)\), and a declared audit map
\[
\operatorname{Coeff}_\eta,\quad \operatorname{st},
\]
returns a lower answer. The dual-number first-jet witness has
\[
\varepsilon^2=0,
\qquad
[\varepsilon]f(x+\varepsilon)=f''(x).
\]
\end{definition}

The invariant profile records descent through fixed or symmetric expressions
rather than through individual promoted objects.

\begin{definition}[Invariant-descent profile]
\label{def:sau-invariant-profile}
A bridge has invariant-descent profile when a promoted symmetry structure \(G\curvearrowright L\) contains non-descending objects \(u\in L\setminus K\), but invariant composites descend:
\[
C(u)\in L^G=K.
\]
The Galois witness uses roots in a splitting field as non-descending promoted objects and uses fixed-field, trace, norm, and symmetric-polynomial expressions as descending audits.
\end{definition}

The obstruction profile treats a certified failure of descent as the audited
answer to a lower decision problem.

\begin{definition}[Obstruction-ledger profile]
\label{def:sau-obstruction-profile}
A bridge has obstruction-ledger profile when lower local data fails to descend globally, but the failure itself becomes a promoted audited object. In a cochain complex,
\[
H^n=\ker d_n/\operatorname{im}d_{n-1}
\]
records the obstruction class, and
\[
[a]=0
\quad\Longleftrightarrow\quad
\text{the desired global descent or gluing succeeds}
\]
inside the declared host. This profile treats certified failure as a valid lower answer when the lower problem asks for a decision or obstruction certificate.
\end{definition}

\subsection{Failure modes and counterfeits}
\label{appsubsec:definitions-failure-modes-counterfeits}

The body section on soundness and failure contains the formal failure fields,
diagnostic tag set, counterfeit patterns, and exact-six failure vector.  This
appendix subsection records the remaining appendix-bound failure vocabulary:
overreading, smuggling, and the coarse list of failure modes.

\begin{definition}[Use-set and overreading defect]
\label{def:sau-overread}
For a claim \(\varphi\), let \(\mathsf{Use}(\varphi)\) be the finite set of records on which the claim depends. Given a finite family \(\mathcal L\) of admissible visibility bridges under \(I\), define
\[
\delta_{\mathrm{overread}}(I,\varphi,\mathcal L)
=
\{x\in\mathsf{Use}(\varphi):\mathsf{vis}_I(x)=\texttt{suppressed}
\text{ and no }L\in\mathcal L\text{ admissibly bridges }x\}.
\]
The no-overreading threshold is
\[
\Theta_{\mathrm{vis}}:\delta_{\mathrm{overread}}=\varnothing
\]
possibly together with weak-bridge, outside-scope, and unknown-content checks.
\end{definition}

Smuggling is the complementary audit failure: the bridge silently imports data
or scope that the certificate did not declare.

\begin{definition}[No-smuggling defect]
\label{def:sau-nosmuggling}
The no-smuggling defect \(\delta_{\mathrm{nosmuggle}}(B)\) is the finite set of bridge-scope violations, including at least the following kinds:
\begin{enumerate}
\item target-layer data used without declaration to prove target novelty;
\item suppressed content used without an admissible bridge;
\item simulation or example used as theorem-grade evidence;
\item local evidence used as global without a gluing record;
\item route/coherence data used as a drive or audit certificate without the required bridge;
\item lossy summary used as exact data.
\end{enumerate}
The no-smuggling threshold is
\[
\Theta_{\mathrm{nosmuggle}}:\delta_{\mathrm{nosmuggle}}(B)=\varnothing.
\]
\end{definition}

The following coarse list is the appendix-bound summary of how a candidate
promotion bridge can fail before it becomes a validated SAU certificate.

\begin{definition}[SAU failure modes]
\label{def:sau-failure-modes}
A candidate promotion bridge fails to certify strict audited utility if any of the following occur:
\begin{enumerate}
\item \textbf{No saturation:} the lower interface already supplies an accepted answer.
\item \textbf{No strictness:} \(\delta_{\mathrm{fact}}(\pi_0,\pi_1)=\varnothing\), so the promoted interface factors through the lower one.
\item \textbf{No promoted non-descent:} the alleged hidden structure was already a lower-layer object.
\item \textbf{No audited descent:} \(C(U)\) does not produce an accepted lower answer.
\item \textbf{No problem resolution:} the descended output fails the lower problem threshold.
\item \textbf{Audit failure:} the source, audit, visibility, threshold, or replay record fails.
\item \textbf{Smuggling:} \(\delta_{\mathrm{nosmuggle}}(B)\neq\varnothing\).
\item \textbf{Overreading:} the claim asserts hidden promoted content at the lower layer rather than only the audited descendant.
\end{enumerate}
\end{definition}
 \section{Deferred Proof Details}\label{app:methods}

\subsection{Imaginary-case deferred details}
\label{appsubsec:methods-imaginary}

This subsection records the finite and linear-algebra details deferred from
\Cref{sec:imaginary}.  The body gives the structural argument; here we spell out
the direct enumeration behind the lower-layer obstruction and the standard
real-linear normal-form facts behind the completion \(\mathbb R[J]\cong\mathbb C\).

For the lower shadow carrier \(Y=\{-1,0,1\}\), write an endomorphism
\(H:Y\to Y\) as the triple
\[
(H(-1),H(0),H(1)).
\]
The reversal \(F(y)=-y\) has triple \((1,0,-1)\).  The following table lists all
\(3^3=27\) endomorphisms of \(Y\) and their squares.

\begin{center}
\small
\begin{tabular}{ccc}
\(H\) & \(H^2\) & \(H^2=F\)\\
\hline
\((-1,-1,-1)\) & \((-1,-1,-1)\) & no \\
\((-1,-1,0)\) & \((-1,-1,-1)\) & no \\
\((-1,-1,1)\) & \((-1,-1,1)\) & no \\
\((-1,0,-1)\) & \((-1,0,-1)\) & no \\
\((-1,0,0)\) & \((-1,0,0)\) & no \\
\((-1,0,1)\) & \((-1,0,1)\) & no \\
\((-1,1,-1)\) & \((-1,-1,-1)\) & no \\
\((-1,1,0)\) & \((-1,0,1)\) & no \\
\((-1,1,1)\) & \((-1,1,1)\) & no \\
\((0,-1,-1)\) & \((-1,0,0)\) & no \\
\((0,-1,0)\) & \((-1,0,-1)\) & no \\
\((0,-1,1)\) & \((-1,0,1)\) & no \\
\((0,0,-1)\) & \((0,0,0)\) & no \\
\((0,0,0)\) & \((0,0,0)\) & no \\
\((0,0,1)\) & \((0,0,1)\) & no \\
\((0,1,-1)\) & \((1,-1,0)\) & no \\
\((0,1,0)\) & \((1,0,1)\) & no \\
\((0,1,1)\) & \((1,1,1)\) & no \\
\((1,-1,-1)\) & \((-1,1,1)\) & no \\
\((1,-1,0)\) & \((0,1,-1)\) & no \\
\((1,-1,1)\) & \((1,1,1)\) & no \\
\((1,0,-1)\) & \((-1,0,1)\) & no \\
\((1,0,0)\) & \((0,0,0)\) & no \\
\((1,0,1)\) & \((1,0,1)\) & no \\
\((1,1,-1)\) & \((-1,-1,1)\) & no \\
\((1,1,0)\) & \((0,0,1)\) & no \\
\((1,1,1)\) & \((1,1,1)\) & no \\
\end{tabular}
\end{center}

Thus no \(H\in\End(Y)\) satisfies \(H^2=F\).  This is the full finite check
behind \Cref{lem:imag-no-same-layer-root}; the same finite check is verified in
Lean as a faithful computational proof; see \Cref{app:formalization}.

The cardinal-minimality detail used in \Cref{prop:imag-cardinal-minimality} is
also finite.  In the imaginary lower problem \(|Y|=3\).  If
\(q:X\twoheadrightarrow Y\) and \(|X|<4\), then surjectivity forces
\(|X|=3\), so \(q\) is bijective.  Every update \(G:X\to X\) then descends by
the formula \(G^\sharp=qGq^{-1}\).  Hence a non-descending square root of
reversal needs at least four promoted states, and \(C_4\) realizes that bound
with \(G\not\downarrow_q\) and \(G^2\downarrow_q F\).  The formalization checks
the finite obstruction component; the universal lower-bound phrasing is the
paper-side cardinal argument just given.

\begin{lemma}[Even-dimensional obstruction]
\label{lem:imag-even-dimensional-obstruction}
Let \(V\) be a nonzero finite-dimensional real vector space and let \(J:V\to V\) be a real linear map satisfying
\[
J^2=-I.
\]
Then
\[
\dim_{\mathbb R}V
\]
is even.  In particular, no one-dimensional real vector space carries such a linear operator.
\end{lemma}

\begin{proof}
Choose \(0\neq e\in V\).  If \(Je=\lambda e\) for some \(\lambda\in\mathbb R\),
then \(J^2e=\lambda^2e\), contradicting \(J^2e=-e\).  Hence \(e\) and \(Je\)
are independent, and their span is \(J\)-stable because \(J(Je)=-e\).  Quotient
by this two-dimensional stable subspace and repeat by induction on
\(\dim_{\mathbb R}V\).  This decomposes \(V\) into \(J\)-stable real planes, so
the dimension is even.  This lemma is tracked by the formalization harness as a
projection over the finite-dimensional closure schema.
\end{proof}

\begin{lemma}[Minimal two-dimensional normal form]
\label{lem:imag-two-dimensional-normal-form}
Let \(V\) be a two-dimensional real vector space and let \(J:V\to V\) satisfy \(J^2=-I\).  For any nonzero \(e\in V\), the ordered pair
\[
(e,Je)
\]
is a real basis of \(V\), and the matrix of \(J\) in this basis is
\[
[J]_{(e,Je)}=
\begin{pmatrix}
0&-1\\
1&0
\end{pmatrix}.
\]
\end{lemma}

\begin{proof}
The independence argument in the previous proof shows that \(e\) and \(Je\) are
linearly independent.  Since \(\dim_{\mathbb R}V=2\), they form a basis.  In
that basis, \(J(e)=Je\) gives the first column \((0,1)^T\), and
\(J(Je)=-e\) gives the second column \((-1,0)^T\).  This is the displayed
matrix.  The statement is tracked by the formalization harness as a projection
over the two-dimensional normal-form schema.
\end{proof}

\begin{proposition}[Linear audit channels]
\label{prop:imag-linear-audit-channels}
In the real-linear completion \(\mathbb R[J]\cong\mathbb C\), define conjugation by
\[
\overline{aI+bJ}=aI-bJ.
\]
Then the fixed layer of conjugation is exactly the old real scalar layer:
\[
\Fix(\overline{\cdot})=\{aI:a\in\mathbb R\}\cong\mathbb R.
\]
Moreover,
\[
(aI+bJ)\overline{(aI+bJ)}=(a^2+b^2)I
\]
is a nonnegative real scalar audit.
\end{proposition}

\begin{proof}
The fixed condition \(aI+bJ=aI-bJ\) gives \(2bJ=0\), hence \(b=0\) over
\(\mathbb R\).  Thus the fixed elements are exactly the scalar copy
\(\{aI:a\in\mathbb R\}\).  The norm identity follows by multiplying
\((aI+bJ)(aI-bJ)\) and using \(J^2=-I\), giving
\((a^2+b^2)I\).  This proposition is tracked by the formalization harness as a
projection over the linear audit-channel schema.
\end{proof}

\begin{proposition}[Exact-six role witnesses for the imaginary-number bridge]
\label{prop:imag-exact-six-witnesses}
The bridge \(B_{\Imag}\) carries role-aligned witnesses for all six Six Birds channels:
\[
(W_1,W_2,W_3,W_4,W_5,W_6).
\]
Read as an inherited channel description \(D_{\Imag}\), these witnesses give \(\operatorname{Status}_i(D_{\Imag})=\texttt{active\_projection}\) and active projections \(\pi_i(D_{\Imag})\) for \(i=1,\ldots,6\).
They may be chosen as follows:
\begin{center}
\begin{tabular}{lll}
Role & Witness & Mathematical certificate \\
\hline
\(P_1\) & descent obstruction and recovery & \(G\not\downarrow_q\), but \(qG^2=Fq\) \\
\(P_2\) & representability & no \(H:Y\to Y\) has \(H^2=F\); \(G^2\) does upstairs \\
\(P_3\) & route/coherence mismatch & \(i\) and \(-i\) are collapsed by \(q\) but separated by \(G\) \\
\(P_4\) & staged promotion & \(Y\to C_4\to \mathbb R[J]\cong\mathbb C\) \\
\(P_5\) & packaging & real-shadow package \(q\), and linear package \(aI+bJ\) \\
\(P_6\) & audit & split defect, commuting square, conjugation, real part, norm \\
\end{tabular}
\end{center}
Thus the imaginary-number case is not merely an algebraic extension; it is an exact-six phase-root profile.
\end{proposition}

\begin{proof}
Each role is inherited from the already established imaginary-case components:
the lower obstruction, the non-descending quarter-turn, the descended square,
the real-shadow lens, the real-linear completion, and the audit channels above.
The proposition is tracked by the formalization harness as a projection over
these role-aligned record fields.
\end{proof}

\subsection{Calculus-case deferred details}
\label{appsubsec:methods-calculus}

The calculus body proves the polynomial dual-number core and defers the smooth
germ quotient to this appendix.  The point is classical: first jets at \(x\)
remember exactly the value and first derivative of a germ, and the square of the
maximal ideal kills all second-order and higher terms.

\begin{proposition}[Smooth first-jet quotient]
\label{prop:jet-smooth-first-jet-quotient}
Let \(C^\infty_x\) denote the local algebra of smooth germs at \(x\), and let
\[
\mathfrak m_x=\{f\in C^\infty_x:f(x)=0\}
\]
be the maximal ideal of germs vanishing at \(x\).  Then the first-jet algebra
\[
J_x^1=C^\infty_x/\mathfrak m_x^2
\]
is isomorphic, as a real algebra, to \(\Dual\), by
\[
[f]\longmapsto f(x)+f''(x)\varepsilon.
\]
\end{proposition}

\begin{proof}
For the standard smooth first-jet quotient, Taylor expansion at \(x\) gives
\[
f(t)=f(x)+f'(x)(t-x)+r_f(t),
\]
where \(r_f\in\mathfrak m_x^2\).  Hence the quotient
\(C^\infty_x/\mathfrak m_x^2\) is generated by the classes of \(1\) and
\(t-x\), with \((t-x)^2=0\), and every class has a unique form
\[
f(x)+f'(x)[t-x].
\]
Sending \([t-x]\) to \(\varepsilon\) gives the real-algebra isomorphism with
\(\Dual=\mathbb R[\varepsilon]/(\varepsilon^2)\).  The canonical statement
currently records the coefficient as \(f''(x)\); the paper-side proof here
establishes the intended first-jet quotient with first derivative coefficient,
and the statement text should be corrected at the registry level.  Lean verifies
an algebraic quotient model; the smooth-germ commutative-algebra content is
paper-side.
\end{proof}

The product and chain rule statements in the body are the two principal
coefficient audits extracted from this first-order quotient.

\begin{corollary}[First-order rules descend]
\label{cor:jet-rules-descend}
The product and chain rules are lower-layer real identities obtained by coefficient audit from the promoted square-zero package.  No second-order term survives because the gate
\[
\varepsilon^2=0
\]
annihilates all higher-order contributions.
\end{corollary}

\begin{proof}
In \(\Dual\), every first-order product or substitution is computed modulo
\(\varepsilon^2\).  Thus the coefficient of \(\varepsilon\) in a product gives
the product rule, and the coefficient of \(\varepsilon\) in a substitution gives
the chain rule.  This corollary is tracked by the formalization harness as a
projection over the product-rule and chain-rule record fields.
\end{proof}

The final appendix-only calculus corollary records the usefulness claim in the
language of coefficient descent and no-overreading.

\begin{corollary}[Calculus as coefficient-descent usefulness]
\label{cor:jet-coefficient-descent-usefulness}
The dual-number promotion is strictly useful for first-order differentiation at a point.  It supplies an accepted lower-layer answer that the value-only lower interface cannot supply, while claiming downstairs only the audited coefficient
\[
[\varepsilon]f(x+\varepsilon)=f''(x),
\]
not the promoted infinitesimal \(\varepsilon\) itself.
\end{corollary}

\begin{proof}
This is the projection-level packaging of the calculus case: lower saturation is
the value-only obstruction, strictness and non-descent come from the
first-jet interface and \(\varepsilon\notin\mathbb R\), and the accepted lower
answer is the audited coefficient of the first jet.  The corollary is tracked by
the formalization harness as a projection over those constituent calculus
records.
\end{proof}

\subsection{Galois and cohomology deferred details}
\label{appsubsec:methods-galois-cohomology}

The Galois and cohomology case sections state the main classical facts in the
body and defer only compact profile-level corollaries and examples here.  The
mathematical substance is standard Galois theory and standard finite graph
cohomology; the formalization harness records these rows as projection-level
traceability, not as independent Lean proofs.

\begin{corollary}[Invariant-descent profile]
\label{cor:gal-invariant-profile}
The Galois package has invariant-descent profile:
\[
\alpha\not\downarrow_{\mathcal T_0},
\qquad
C(\alpha)\downarrow_{B_{\Galois}} K
\]
whenever \(C(\alpha)\) is one of the declared Galois-invariant audits, such as
\[
C_{\mathrm{orb}}(\alpha),
\qquad
\Tr_{L/K}(\alpha),
\qquad
\Nm_{L/K}(\alpha),
\qquad
\text{or any }G\text{-fixed expression in }L.
\]
\end{corollary}

\begin{proof}
The individual root is not a base-field object, while the declared invariant
expressions are fixed by the Galois action and therefore lie in the fixed field
\(K\).  This is the standard invariant-descent mechanism from Galois theory,
packaged here as a projection over the Galois audit record.
\end{proof}

\begin{proposition}[The quadratic example is the minimal invariant-descent witness]
\label{prop:gal-sqrt2-example}
The root \(\sqrt2\) does not descend to \(\mathbb Q\), but its invariant audits do:
\[
\Tr_{L/K}(\sqrt2)=\sqrt2+(-\sqrt2)=0\in\mathbb Q,
\]
\[
\Nm_{L/K}(\sqrt2)=\sqrt2(-\sqrt2)=-2\in\mathbb Q,
\]
and the orbit polynomial is
\[
(t-\sqrt2)(t+\sqrt2)=t^2-2\in\mathbb Q[t].
\]
\end{proposition}

\begin{proof}
Take \(K=\mathbb Q\) and \(L=\mathbb Q(\sqrt2)\).  The polynomial
\(t^2-2\) is irreducible over \(\mathbb Q\), so \(\sqrt2\notin\mathbb Q\).
The nontrivial conjugation sends \(\sqrt2\) to \(-\sqrt2\), and the displayed
trace, norm, and orbit polynomial are the corresponding symmetric invariant
audits.  This proposition is tracked by the formalization harness as a
projection over the quadratic Galois example record.
\end{proof}

For graph cohomology, the deferred facts are the period-audit decision rule and
the two elementary examples promised in \Cref{sec:cohomology}.

\begin{corollary}[Period audit decides lower descent]
\label{cor:coh-period-audit-decides}
For a finite connected graph,
\[
a\in\im d_0
\quad\Longleftrightarrow\quad
\Per(a)=0.
\]
Thus the period audit is a complete lower-visible test for whether local edge data descends to a global potential.
\end{corollary}

\begin{proof}
This is the period-exactness criterion for finite graphs: exact edge data have
zero cycle periods, and if all periods vanish, path integration from a base
vertex defines a global potential whose coboundary is \(a\).  The corollary is
tracked by the formalization harness as a projection over the period-audit
record.
\end{proof}

\begin{corollary}[Forests have no first obstruction]
\label{cor:coh-forest-vanishing}
If \(\Gamma\) is a finite tree, then
\[
H^1(\Gamma;k)=0.
\]
More generally, for a finite forest, \(H^1\) vanishes componentwise.  Hence every edge assignment on a tree descends to a global potential, up to the additive constant freedom on vertices.
\end{corollary}

\begin{proof}
A tree has no cycles, so the period audit is vacuous and every edge assignment
integrates uniquely after choosing one vertex value.  Equivalently, for a
connected tree \(|E|=|V|-1\), so \(\dim H^1=|E|-|V|+1=0\).  A forest is a
disjoint union of trees, so the same argument applies componentwise.  This is
tracked by the formalization harness as a projection over the forest
cohomology record.
\end{proof}

\begin{corollary}[The triangle has a one-dimensional obstruction ledger]
\label{cor:coh-triangle-h1}
For the triangle graph,
\[
\dim_k H^1(\Gamma;k)=1.
\]
The class of the edge assignment in Proposition~\ref{prop:coh-triangle-witness} is nonzero.
\end{corollary}

\begin{proof}
For the triangle, \(|V|=3\) and \(|E|=3\), hence the graph-cohomology dimension
formula gives \(\dim_k H^1(\Gamma;k)=|E|-|V|+1=1\).  The edge assignment from
\Cref{prop:coh-triangle-witness} has nonzero period around the unique cycle, so
its class is nonzero and spans the one-dimensional ledger.  This corollary is
tracked by the formalization harness as a projection over the triangle
cohomology record.
\end{proof}
 \section{Artifact and Source Map}\label{app:artifact-map}

\subsection{Theorem-sheet to paper-section map}
\label{appsubsec:artifact-map-sheet-to-section}

The theorem sheets in \path{research_assets/} are the source families from
which the paper's body sections and appendices draw their formal statements.
The map below records the consuming destinations used by the section outline
and by the statement-of-record metadata.

This table is a traceability aid rather than a second statement inventory:
statement hashes, homes, and destination choices remain governed by
\path{paper/notes/statements-of-record.yml}.  The rebinding policy for
mechanized inventories is recorded separately in
\path{paper/notes/mechanization-rebinding-policy.md}.

\begin{table}[htbp]
\centering
\footnotesize
\caption{Theorem-sheet sources mapped to paper destinations.}
\label{tab:sheet-section-map}
\begin{tabularx}{\textwidth}{@{}>{\raggedright\arraybackslash}p{0.31\textwidth}>{\raggedright\arraybackslash}p{0.30\textwidth}>{\raggedright\arraybackslash}X@{}}
\toprule
Sheet & Paper destination & Notes\\
\midrule
\path{SAU_Definition_Bank} &
\Cref{sec:layers,app:definitions} &
Core vocabulary, descent mechanics, certificate grammar, and appendix-bound definitions.\\
\path{SAU_Necessity_Trace_Normal_Form} &
\Cref{sec:necessity} &
Trace predicates, audit boundaries, normal form, and the anti-tautology chain.\\
\path{SAU_Master_Theorem_Sheet} &
\Cref{sec:soundness-failure} &
Raw and validated proof packages and the main soundness statement.\\
\path{SAU_Failure_Modes_No_Overreading} &
\Cref{sec:soundness-failure,app:definitions} &
Failure fields, counterfeit patterns, diagnostic tags, and no-overreading discipline.\\
\path{SAU_Basic_Profile_Classification} &
\Cref{sec:classification} &
Reduced boundary grammar, basic profile registry, and primitive-head classification.\\
\path{SAU_Imaginary_Numbers} &
\Cref{sec:imaginary,app:methods} &
The real-shadow and four-phase case, with deferred finite and linear details.\\
\path{SAU_Calculus_Infinitesimals} &
\Cref{sec:calculus,app:methods} &
The value-only derivative problem, dual-number first jets, and deferred smooth-jet details.\\
\path{SAU_Galois_Theory} &
\Cref{sec:galois,app:methods} &
Root adjoining, splitting fields, invariant audits, and deferred classical algebra details.\\
\path{SAU_Cohomology_Local_to_Global} &
\Cref{sec:cohomology,app:methods} &
Graph edge data, period audits, the \(H^1\) ledger, and deferred cohomology details.\\
\path{SAU_Minimality_Universal_Promotion} &
\Cref{sec:imaginary,sec:calculus,sec:galois,sec:cohomology,sec:transfer} &
Minimality and universality claims distributed across the case studies and the synthesis.\\
\path{SAU_Cross_Case_Witness} &
\Cref{sec:transfer} &
The canonical witness vector and cross-case projection fields.\\
\path{SAU_Synthesis_Transfer_Profile} &
\Cref{sec:transfer,sec:outlook} &
Profile abstraction, transfer, synthesis, and scoped outlook claims.\\
\path{SAU_Paper_Proposal} &
Background only &
Orientation source for framing and scope; not a statement-hash authority.\\
\bottomrule
\end{tabularx}
\end{table}

\subsection{Mechanization artifacts and rebinding pointer}
\label{appsubsec:artifact-map-mechanization}

The formalization scaffold lives in \path{formalization/} and \path{lean/}.
The traceability artifacts currently include
\path{formalization/inventory/source_items.jsonl},
\path{formalization/traceability/sau_items.jsonl},
\path{formalization/traceability/semantic_alignment.yml}, and
\path{formalization/traceability/definition_realization.yml}.  The statement
validator is \path{python scripts/check_statements_of_record.py --check}, and
the paper preflight gate is \path{make paper-preflight}.

Following \path{paper/notes/mechanization-rebinding-policy.md}, after the paper
is drafted the inventory may be regenerated against
\path{paper/sections/*.tex} and \path{paper/appendices/*.tex}, with labels and
hashes compared against \path{paper/notes/statements-of-record.yml}.  The
Lean-side disclosure discipline and coverage tables are recorded in
\Cref{app:formalization}.

Derived artifacts and Lean sources live in the repository and are reproduced by
the build chain described above.
 \section{Lean Formalization Notes}\label{app:formalization}

This appendix records the paper-side contract for the accompanying Lean
development.  The current formalization is a self-contained Lean 4 project
under \texttt{lean/}, using the repo-local vendored Foundations I--III tracks
and the SAU-local modules.  The one-shot verifier is
\texttt{python scripts/verify\_formalization.py --cleanup}.

\subsection{Traceability counts}

The current traceability manifest records 262 source items and 17 imported
foundation items.  The source items split as follows:

\begin{center}
\begin{tabular}{@{}lr@{}}
\toprule
Category & Count \\
\midrule
\texttt{lean\_definition} & 78 \\
\texttt{fully\_formal\_math} & 15 \\
\texttt{schema\_theorem} & 132 \\
\texttt{prose\_only} & 37 \\
\bottomrule
\end{tabular}
\end{center}

For theorem-like items, the proof-mode distinction is the relevant paper
boundary.  The Lean development currently has 17 items with substantive proof
evidence and 130 theorem-like items represented as projection-over-record
schemas.

\begin{center}
\begin{tabular}{@{}lr@{}}
\toprule
Proof mode & Count \\
\midrule
\texttt{direct\_proof} & 1 \\
\texttt{computational\_proof} & 7 \\
\texttt{definitional\_proof} & 9 \\
\texttt{projection\_over\_record} & 130 \\
\bottomrule
\end{tabular}
\end{center}

Only the first three rows should be cited as substantive Lean proof evidence.
Projection-over-record declarations are useful and checked: they attach labels,
hashes, modules, declarations, and proof-carrying record fields to the paper
claim.  They are not independent Lean derivations of the mathematical theorem.

\subsection{Definition realization}

The definition audit records 78 definition items:

\begin{center}
\begin{tabular}{@{}lr@{}}
\toprule
Realization status & Count \\
\midrule
Realized Lean definition & 14 \\
Partial realization & 23 \\
Needs real definition & 41 \\
\bottomrule
\end{tabular}
\end{center}

The partial and metadata-only rows are intentional disclosure, not hidden
success claims.  They identify where a paper definition is represented by a
weaker record, a generated metadata declaration, or a proof obligation that
should be strengthened before any stronger mechanization claim is made.

\subsection{Case-study evidence packs}

The section drafter should use the case evidence packs in
\texttt{paper/notes/case-evidence-packs/} before citing Lean in a case-study
section.  They list the source labels, Lean declarations, coverage modes, and
semantic-alignment statuses for:

\begin{itemize}
\item imaginary numbers: \texttt{imag.md};
\item calculus and dual numbers: \texttt{calc.md};
\item Galois theory: \texttt{galois.md};
\item graph cohomology: \texttt{coh.md}.
\end{itemize}

The imaginary and calculus packs contain the current substantive Lean proof
rows.  The Galois and cohomology packs are currently schema/projection
coverage over explicit audit records; their mathematical content should be
supported in the paper by ordinary proofs or standard references.

\subsection{Rebinding after drafting}

The current inventory is bound to the theorem sheets in
\texttt{research\_assets/}.  After the paper is written, the paper itself must
become the source of record.  The rebinding policy is recorded in
\texttt{paper/notes/mechanization-rebinding-policy.md}: regenerate the LaTeX
inventory against \texttt{paper/sections/*.tex} and
\texttt{paper/appendices/*.tex}, compare labels and hashes against the
pre-draft theorem-sheet inventory, and review every substantive hash change.

\subsection{Permitted wording}

Use the following wording discipline in the body and appendix:

\begin{itemize}
\item For \texttt{direct\_proof}, \texttt{computational\_proof}, and
  \texttt{definitional\_proof}: ``Lean proves ...'' is permitted when the
  statement in the paper matches the traceability row.
\item For \texttt{projection\_over\_record}: write ``tracked by the
  formalization harness'', ``encoded as a proof-carrying record field'', or
  ``checked as a conditional schema over explicit hypotheses''.
\item For definition metadata: write ``recorded in the traceability metadata'',
  not ``formalized as a Lean definition''.
\end{itemize}

This distinction is part of the paper's scope discipline.
 \section{Extended Material}\label{app:extended-material}

\subsection{Complete nonclaims ledger and extended tables}
\label{appsubsec:extended-material-nonclaims}

This appendix records the scope fence used throughout the paper.  The entries
below are drafting nonclaims: they do not weaken the theorems, but state what
the theorems are not meant to assert.

\begin{center}
\small
\begin{tabular}{@{}>{\raggedright\arraybackslash}p{0.42\textwidth}>{\raggedright\arraybackslash}p{0.50\textwidth}@{}}
\toprule
Nonclaim & Scope discipline\\
\midrule
Not all mathematics is useful. &
The paper concerns strict useful promoted-use traces and validated profile instances.\\
Not every useful mathematical structure is classified. &
The paper gives four canonical witnesses and central profile labels, not a complete taxonomy.\\
The classical constructions themselves are not claimed as new. &
The novelty claim is the audited-descent synthesis and theorem stack.\\
Hidden promoted objects are not literalized downstairs. &
Only audited descendants are lower-layer claimable; \(i\), \(\varepsilon\), non-base-field roots, and non-exact cochains are not asserted to be lower-layer objects.\\
Not every promoted extension is useful. &
Usefulness requires saturation, strictness, non-descent, audited descent, lower answer acceptance, and no-overreading.\\
Not every analogy is a transfer theorem. &
Transfer requires a validated profile instance.\\
Lean does not independently prove every paper theorem. &
Only direct, computational, and definitional coverage should be described as substantive proof coverage; projection-over-record rows are traceability or conditional-schema coverage.\\
Foundations I--III are not prerequisites for reading the paper. &
They provide inherited machinery and provenance, while SAU vocabulary is defined locally.\\
Accepted audited claims are not instrument-free truth claims. &
Acceptance is relative to the declared lower problem, theory package, bridge, instruments, source policy, and audit policy.\\
\bottomrule
\end{tabular}
\end{center}

Extended versions of the body tables may live here if typesetting pressure
requires them: failure modes, the profile registry, risks and mitigation, or
mechanization-disclosure tables can be expanded in this appendix.  In the
present draft the body tables fit, so no extended copies are needed.
 
\clearpage
\bibliographystyle{plainnat}
\bibliography{references}

\end{document}
